% ===========================================================================
%  KTH -- Team Reference Document  (NCPC / NWERC)
%  GENERATED FILE: edit src/*.tex, then run tools/build.sh
%  Paste notebook.tex into Overleaf and hit Recompile -- plain pdfLaTeX.
%
%  Layout matches KACTL: A4 LANDSCAPE, 3 columns. Code is 8pt, one point
%  ABOVE KACTL's (kactl uses \documentclass[9pt,a4paper]{extreport} + geometry
%   a4paper,landscape,hmargin={0.5cm,0.5cm},vmargin={1.0cm,0.4cm}
%   + multicols*{3} + \footnotesize listings; size9.clo makes that 7pt.)
%
%  NWERC rules encoded here: <= 25 pages, single-sided, A4, university name
%  in the UPPER LEFT, page number in the UPPER RIGHT.
% ===========================================================================
\documentclass[9pt,a4paper]{extarticle}

\newcommand{\uni}{KTH Royal Institute of Technology}
\newcommand{\unishort}{KTH}
\newcommand{\team}{$O(N^{\text{5h debugging}})$}
\newcommand{\members}{Ruben \quad Ruiming \quad Martin}

% Code size. 8pt is a point above KACTL's 7pt, which is printed for NWERC every year.
% DO NOT go below it -- raise both numbers if the print looks tight.
\newcommand{\codesize}{\fontsize{8}{8.8}\selectfont}

% includehead makes geometry reserve the header INSIDE the top margin. Without it
% LaTeX puts the header box top at 1in+\topmargin, which lands ABOVE the paper edge
% here and slices the header in half. (KACTL papers over the same problem by setting
% \headheight to 25pt; reserving the space properly is the honest fix.)
\usepackage[a4paper,landscape,hmargin={0.5cm,0.5cm},vmargin={0.4cm,0.4cm},
            includehead,headsep=4pt]{geometry}
\usepackage{multicol}
\usepackage{lmodern}
\usepackage[T1]{fontenc}
\usepackage{listings}
\usepackage{xcolor}
\usepackage{titlesec}
\usepackage{fancyhdr}
\usepackage{amsmath,amssymb}
\usepackage{tocloft}
\usepackage{graphicx}
\usepackage[expansion=false]{microtype}

\definecolor{secbg}{HTML}{1F3B57}
\definecolor{hdbg}{HTML}{E8ECF1}
\definecolor{kw}{HTML}{16507B}
\definecolor{cmt}{HTML}{6B7C6B}
\definecolor{str}{HTML}{8B3A2E}
\definecolor{rule}{HTML}{9AA7B4}

\setlength{\parindent}{0pt}
\setlength{\parskip}{0pt}
\setlength{\columnsep}{4mm}
\setlength{\columnseprule}{0.2pt}

\newcommand{\secbar}[1]{%
  \colorbox{secbg}{\parbox{\dimexpr\columnwidth-2\fboxsep}{%
    \strut\color{white}\bfseries #1}}}
\titleformat{\section}{\normalfont}{}{0pt}{\secbar}
\titlespacing{\section}{0pt}{1.1ex}{0.6ex}

\pagestyle{fancy}
\fancyhf{}
\renewcommand{\headrulewidth}{0.2pt}
\lhead{\footnotesize\uni}                 % NWERC: university, upper left
\chead{\footnotesize\team}
\rhead{\footnotesize\thepage}             % NWERC: page number, upper right
% headsep is set by geometry above; do not override it here or the header
% drifts back off the top of the page.

\lstset{
  language=C++,
  basicstyle=\ttfamily\codesize,
  keywordstyle=\color{kw}\bfseries,
  commentstyle=\color{cmt},
  stringstyle=\color{str},
  showstringspaces=false,
  keepspaces=true,
  columns=fullflexible,
  breaklines=false,
  tabsize=4,
  aboveskip=0.8ex,
  belowskip=0.3ex,
  xleftmargin=0pt,
  morekeywords={ll,ld,vi,pii,matrix,rep,down,all,i128,ull}
}

% Compact index: every \section and every \entry gets a table-of-contents line.
% secnumdepth 0 = unnumbered sections. The section bars show no number either,
% and with numbering on, \numberline would overprint the title in the index.
\setcounter{secnumdepth}{0}
\setcounter{tocdepth}{2}
\setlength{\cftbeforesecskip}{0.35ex}
\setlength{\cftbeforesubsecskip}{0pt}
\setlength{\cftsecindent}{0pt}
\setlength{\cftsubsecindent}{0.9em}
\setlength{\cftsecnumwidth}{0em}
\setlength{\cftsubsecnumwidth}{0em}
\renewcommand{\cftsecfont}{\bfseries\footnotesize}
\renewcommand{\cftsecpagefont}{\bfseries\footnotesize}
\renewcommand{\cftsecleader}{\normalfont\hfill}
\renewcommand{\cftsubsecfont}{\footnotesize}
\renewcommand{\cftsubsecpagefont}{\footnotesize}
\renewcommand{\cftsubsecleader}{\normalfont\hfill}
\renewcommand{\cftdot}{}

% #1 name, #2 complexity / note, #3 description
\newcommand{\entry}[3]{%
  \addcontentsline{toc}{subsection}{#1}%
  \vspace{0.8ex}%
  \noindent\colorbox{hdbg}{\parbox{\dimexpr\columnwidth-2\fboxsep}{%
    \strut\textbf{#1}\hfill{\footnotesize\textsf{#2}}}}\par
  \vspace{0.4ex}%
  {\footnotesize #3\par}%
}
% same, but no index line -- for a second listing belonging to the entry above
\newcommand{\entryn}[3]{%
  \vspace{0.8ex}%
  \noindent\colorbox{hdbg}{\parbox{\dimexpr\columnwidth-2\fboxsep}{%
    \strut\textbf{#1}\hfill{\footnotesize\textsf{#2}}}}\par
  \vspace{0.4ex}%
  {\footnotesize #3\par}%
}

\newcommand{\code}[1]{{\ttfamily\small #1}}
% compact two-column description table, sized for a 93mm column
\newenvironment{ctab}[2]{%
  \begingroup\setlength{\tabcolsep}{3pt}%
  \begin{tabular}{@{}p{#1\columnwidth}p{#2\columnwidth}@{}}}%
  {\end{tabular}\endgroup}

\begin{document}

% ------------------------------- front page --------------------------------
% The cover counts toward the 25 pages, and NWERC wants the university upper
% left and the page number upper right on the pages -- so the cover keeps the
% running header and is page 1, like every other page.
% Laid out like KACTL's cover: logo, university, team name, members, contest.
% kth.pdf must sit next to notebook.tex (upload it to Overleaf too).
\begin{center}
\vspace*{1.2cm}
\includegraphics[width=3.2cm]{kth.pdf}\\[0.7cm]
{\fontsize{24}{28}\selectfont \uni}\\[1.3cm]
{\fontsize{46}{50}\selectfont\bfseries \team}\\[1.2cm]
{\huge \members}\\[1.4cm]
\textcolor{rule}{\rule{0.55\textwidth}{0.8pt}}\\[0.9cm]
{\LARGE Team Reference Document \quad\textbullet\quad NCPC / NWERC}
\vfill
\end{center}
\newpage

% ------------------------------ index + body -------------------------------
{\large\bfseries Index}\hfill{\footnotesize\members}
\vspace{0.5ex}
\textcolor{rule}{\hrule height 0.6pt}
\vspace{1.0ex}

\begin{multicols*}{3}
\renewcommand{\contentsname}{\vspace{-3.5ex}}
\tableofcontents
% tocloft forces \thispagestyle{plain} on the table-of-contents page, which
% wipes the header off the index page. Undo it right here -- verified by
% bisecting the preamble against a minimal file.
\thispagestyle{fancy}
\vspace{1.5ex}

\section{Contest setup}

\entry{Header}{paste first}{%
Every entry in this document assumes this block: Ruben's and Ruiming's templates
merged. Ruben's code is \textbf{0-indexed} with \textbf{half-open} ranges
\code{[l, r)}; Ruiming's is \textbf{1-indexed} with global arrays of size
\code{MAXN} --- each entry says which.\\[0.3ex]
\textbf{Update \code{MAXN}} to fit the problem. Ruiming's per-problem globals
(\code{ll t,n,a[MAXN],ans;}) are left out so they don't clash with the entries;
add them back as you need. Ruiming's \code{struct matrix} clashes with the
\code{matrix} typedef --- delete the typedef line when you use it.}

%HEADER
\begin{lstlisting}
//UPDATERA ARRAY STORLEKEN!!!!!!!!!!!!!!!!!!!!!!!!!!!!!!!!!!!!
#include <bits/stdc++.h>
#include <cassert>   // bits/stdc++.h no longer pulls this in
using namespace std;
typedef long long ll;
typedef long double ld;
typedef unsigned long long ull;
typedef __int128 i128;
typedef pair<ll, ll> pii;
typedef vector<ll> vi;
typedef vector<vi> matrix;
#define rep(i, a, b) for(ll i = a; i < b; i++)
#define down(i, b, a) for(ll i = b - 1; i >= a; i--)
#define all(v) (v).begin(), (v).end()
ld eps = 1e-9;
const ll INF=1e18;
const ll MAXN=206;//UPDATERA ARRAY STORLEKEN!!!!!!!!!!!!!!!!!
int main(){
    ios_base::sync_with_stdio(false);
    cin.tie(NULL);
}
\end{lstlisting}

\entry{Compile and stress test}{}{%
Drop the sanitizers for the real submission --- they cost $\sim$3x.
\code{-Wconversion} additionally catches silent
\code{ll}$\rightarrow$\code{int} narrowing.\\[0.3ex]
Stress test: \code{gen.py} prints a small random case, \code{mine} is the fast
solution, \code{ans} the slow one you trust.}

%NOCOMPILE
\begin{lstlisting}
g++ -std=gnu++17 -O2 -Wall -Wextra -Wshadow \
    -fsanitize=address,undefined -g a.cpp
for i in $(seq 1000); do
  python3 gen.py $i > t.txt
  ./mine < t.txt > o1.txt; ./ans < t.txt > o2.txt
  cmp -s o1.txt o2.txt || { echo "WA $i"; cat t.txt; break; }
done
\end{lstlisting}

\entry{int128 I/O}{}{%
\code{cin} / \code{cout} cannot handle \code{\_\_int128}. Range is
$\pm1.7\cdot10^{38}$, so it is the fix for \code{a*b \% m} when $m$ is up to
$9\cdot10^{18}$: \code{(i128)a * b \% m}.}

\begin{lstlisting}
i128 read128() {
    i128 x = 0; int s = 1; char c = getchar();
    while (c == ' ' || c == '\n') c = getchar();
    if (c == '-') { s = -1; c = getchar(); }
    for (; c >= '0' && c <= '9'; c = getchar())
        x = x * 10 + (c - '0');
    return s * x;
}
void print128(i128 x) {
    if (x < 0) { putchar('-'); x = -x; }
    if (x > 9) print128(x / 10);
    putchar('0' + (int)(x % 10));
}
\end{lstlisting}

\section{Range structures}

All ranges are half-open: \code{[l, r)}, \code{r} not included.

\entry{Fenwick (BIT)}{add, sum, kth $O(\log n)$ \quad mem $n$}{%
Needs an operation with an inverse (sum, xor); min/max/gcd need a segment tree.
$\sim$3x faster than \code{Seg}, $\sim$15x faster than \code{LazySeg}.\\[0.3ex]
Range add + point get: \code{add(l, x)}, \code{add(r, -x)}; then \code{sum(i+1)}
\emph{is} \code{a[i]}. Range add + range sum: two BITs, or \code{LazySeg}.
\code{kth} needs every \code{a[i] >= 0} and \code{k >= 1}. Xor instead of sum:
every \code{+=} becomes \code{\textasciicircum=}, and drop the \code{-} in
\code{sum(l, r)}.}

\begin{lstlisting}
struct BIT {
    int n; vector<ll> t;
    BIT(int n = 0) : n(n), t(n + 1, 0) {}
    BIT(const vector<ll>& a) : BIT((int)a.size()) {
        for (int i = 1; i <= n; i++) {
            t[i] += a[i-1];
            int j = i + (i & -i);
            if (j <= n) t[j] += t[i];
        }
    }
    void add(int i, ll x) {                     // a[i] += x
        for (++i; i <= n; i += i & -i) t[i] += x;
    }
    ll sum(int r) {                             // a[0..r)
        ll s = 0;
        for (; r > 0; r -= r & -r) s += t[r];
        return s;
    }
    ll sum(int l, int r) {                      // a[l..r)
        return l < r ? sum(r) - sum(l) : 0;
    }
    int kth(ll k) {     // min i with sum(i+1) >= k, else n
        int p = 0;
        for (int w = 1 << __lg(max(n, 1)); w; w >>= 1)
            if (p + w <= n && t[p + w] < k) p += w, k -= t[p];
        return p;
    }
};
\end{lstlisting}

\entry{Segment tree}{build $O(n)$, set / query $O(\log n)$ \quad mem $2n$}{%
Point assign + range query.\\[0.3ex]
\begin{ctab}{0.15}{0.80}
\code{T} & what one element is, and what a query gives back.\\
\code{op(a,b)} & answers for two neighbouring pieces, \code{a} the left one;
return the answer for both together. Associative, need not commute.\\
\code{ID} & a \code{T} that \code{op} ignores: \code{op(x,ID) = op(ID,x) = x}.
An empty query returns it.\\
\end{ctab}\\[0.3ex]
{\ttfamily\footnotesize
\begin{ctab}{0.16}{0.79}
query= & op \hfill ID\\
sum & \code{a + b} \hfill \code{0}\\
min & \code{min(a,b)} \hfill \code{LLONG\_MAX}\\
max & \code{max(a,b)} \hfill \code{LLONG\_MIN}\\
gcd & \code{\_\_gcd(a,b)} \hfill \code{0}\\
xor & \code{a \textasciicircum{} b} \hfill \code{0}\\
and & \code{a \& b} \hfill \code{\textasciitilde 0LL}\\
or & \code{a | b} \hfill \code{0}\\
\end{ctab}}\\[0.3ex]
Count of min: \code{T = pair<ll,int>}, keep the smaller \code{.first} and sum
\code{.second} on ties, \code{ID = \{LLONG\_MAX, 0\}}. Struct \code{T}: write
\code{inline static const T ID = T\{...\};}\\[0.3ex]
Read-modify-write is \code{s.set(i, s[i] + v)}. Range updates need
\code{LazySeg}. Two different trees in one problem: copy, rename, edit.}

\begin{lstlisting}
struct Seg {
    using T = ll;
    static T op(T a, T b) { return a + b; }
    static constexpr T ID = 0;
    int n; vector<T> t;
    Seg(int n = 0) : n(n), t(2 * n, ID) {}
    Seg(const vector<T>& a)
            : n((int)a.size()), t(2*a.size()) {
        copy(a.begin(), a.end(), t.begin() + n);
        for (int i = n - 1; i > 0; i--)
            t[i] = op(t[2*i], t[2*i+1]);
    }
    T operator[](int i) const { return t[i + n]; }
    void set(int i, T v) {
        for (t[i += n] = v; i /= 2;)
            t[i] = op(t[2*i], t[2*i+1]);
    }
    T query(int l, int r) {               // combine a[l..r)
        T x = ID, y = ID;              // x = left, y = right
        for (l += n, r += n; l < r; l /= 2, r /= 2) {
            if (l & 1) x = op(x, t[l++]);
            if (r & 1) y = op(t[--r], y);
        }
        return op(x, y);
    }
};
\end{lstlisting}

\entry{Sparse table}{build $O(n \log n)$, query $O(1)$ \quad mem $n \log n$}{%
\textbf{Idempotent} \code{op} only --- min, max, gcd, and, or. \code{op(x,x)} must
be \code{x}, because the two halves of a query overlap. No \code{ID} needed, and
no updates ever. For sum use \code{BIT}, for updates use \code{Seg}.}

\begin{lstlisting}
struct SpTab {
    using T = ll;
    static T op(T a, T b) { return min(a, b); }
    int n; vector<vector<T>> t;
    SpTab(const vector<T>& a) : n((int)a.size()), t(1, a) {
        for (int k = 1; (1 << k) <= n; k++) {
            t.push_back(vector<T>(n - (1 << k) + 1));
            for (int i = 0; i < (int)t[k].size(); i++)
                t[k][i] = op(t[k-1][i],
                             t[k-1][i + (1 << (k-1))]);
        }
    }
    T query(int l, int r) {          // a[l..r), needs l < r
        int k = __lg(r - l);
        return op(t[k][l], t[k][r - (1 << k)]);
    }
};
\end{lstlisting}

\entry{Sparse segment tree}{set / add / query $O(\log C)$ \quad mem $V \log_2 (C/V)$ nodes}{%
Same \code{T / op / ID} as \code{Seg}, but \code{[lo, hi)} is given at
construction, may be $10^{18}$ wide or start negative, and nodes appear only where
you write. \code{ID} is \emph{also} what every untouched position reads as, so pick
the neutral element you want spread over the whole range.\\[0.3ex]
\begin{ctab}{0.19}{0.76}
\code{add(i,v)} & \code{a[i] = op(a[i], v)} in one descent; prefer it over
\code{set(i, s[i]+v)}, which pays two.\\
\code{hint} & 3rd ctor arg, how many \code{set}/\code{add} you expect.
\code{reserve} only; 0 or wrong is safe.\\
\code{kth(k)} & sum config, all \code{a[i] >= 0}, \code{k >= 1}: smallest
\code{i} with \code{query(lo, i+1) >= k}, else \code{hi}.\\
\end{ctab}\\[0.3ex]
Multiset on a value axis: count config, \code{add(v, 1)} / \code{add(v, -1)},
\code{query(a, b+1)} counts $[a,b]$, \code{kth} gives order statistics.\\[0.3ex]
Use it \emph{only} when the positions are not all known up front (interactive, or an
update that depends on an earlier answer). Otherwise coordinate-compress and use
\code{BIT}: 3.5x faster, 4.5x less memory. Nodes $\approx 2V + V\log_2(C/V)$, 16 B
each ($V$ = positions touched), so $V = 3\cdot10^5$ over $C = 10^9$ is $\sim$62 MB.}

\begin{lstlisting}
struct SparseSeg {
    using T = ll;
    static T op(T a, T b) { return a + b; }
    static constexpr T ID = 0;
    struct Nd { int l, r; T v; };
    ll lo, hi;
    vector<Nd> nd;               // node 0 = null, 1 = root
    SparseSeg(ll lo = 0, ll hi = 1, int hint = 0)
            : lo(lo), hi(hi) {
        ll C = max(hi - lo, 1LL), h = max(hint, 1);
        size_t cap = 2 + h * (3 + (C > h ? __lg(C / h) : 0));
        nd.reserve(cap);
        node(); node();
    }
    int node() {
        nd.push_back({0, 0, ID});
        return (int)nd.size() - 1;
    }
    void set(ll i, T v) { pt(i, v, true);  }
    void add(ll i, T v) { pt(i, v, false); }
    T operator[](ll i) { return query(i, i + 1); }
    T query(ll l, ll r) {                      // fold a[l..r)
        return l < r ? qry(1, lo, hi, l, r) : ID;
    }
    ll kth(T k) {
        if (nd[1].v < k) return hi;
        int x = 1; ll l = lo, r = hi;
        while (r - l > 1) {
            ll m = l + (r - l) / 2;
            if (nd[nd[x].l].v >= k) { x = nd[x].l; r = m; }
            else { k -= nd[nd[x].l].v; x = nd[x].r; l = m; }
        }
        return l;
    }
    void pt(ll i, T v, bool assign) {
        int st[64], d = 0, x = 1; ll l = lo, r = hi;
        while (r - l > 1) {
            st[d++] = x;
            ll m = l + (r - l) / 2;
            if (i < m) {
                if (!nd[x].l) { int c = node(); nd[x].l = c; }
                x = nd[x].l; r = m;
            } else {
                if (!nd[x].r) { int c = node(); nd[x].r = c; }
                x = nd[x].r; l = m;
            }
        }
        nd[x].v = assign ? v : op(nd[x].v, v);
        while (d--) {
            int p = st[d];
            nd[p].v = op(nd[nd[p].l].v, nd[nd[p].r].v);
        }
    }
    T qry(int x, ll l, ll r, ll ql, ll qr) {
        if (!x || qr <= l || r <= ql) return ID;
        if (ql <= l && r <= qr) return nd[x].v;
        ll m = l + (r - l) / 2;
        return op(qry(nd[x].l, l, m, ql, qr),
                  qry(nd[x].r, m, r, ql, qr));
    }
};
\end{lstlisting}

\entry{Lazy segment tree}{build $O(n)$, update / query $O(\log n)$ \quad mem $4n$}{%
A node holds \code{t}, the answer for its own piece, and \code{d}, an update already
counted into \code{t} but not yet given to its children.\\[0.3ex]
\begin{ctab}{0.15}{0.80}
\code{T} & what one element is, and what a query gives back.\\
\code{F} & one update: exactly what you hand to \code{update(l, r, f)}.\\
\code{op(a,b)} & merge two neighbouring answers, \code{a} the left one.
Associative, need not commute.\\
\code{ID} & neutral for \code{op}; an empty query returns it.\\
\code{ap(x,f,k)} & a node of \code{k} elements has answer \code{x} and \code{f} now
hits all \code{k}: return the node's new answer. You never see the elements, only
\code{x} and \code{k} --- if the new answer is not decided by those two, lazy
cannot do this update at all.\\
\code{cm(a,f)} & \code{a} was already waiting and \code{f} lands on top: return the
single update meaning ``\code{a} first, then \code{f}''.\\
\code{sh(f,k)} & the same update shifted \code{k} places right. An \code{f} starts at
the first element it covers, so the left child gets \code{f} and the right child
\code{sh(f, size of left child)}. Return \code{f} unchanged unless the update varies
along the range: ``add $m, m{+}1, \dots$'' becomes ``add $m{+}k, m{+}k{+}1, \dots$''.\\
\code{NOP} & an \code{F} meaning nothing waiting. For assigning updates use a value
you will never assign (\code{LLONG\_MIN}).\\
\end{ctab}\\[0.3ex]
{\ttfamily\footnotesize
\begin{ctab}{0.30}{0.65}
f does \hfill query= & ap(x,f,k) \hfill cm(a,f) \hfill ID \hfill NOP\\
x += f \hfill sum & x + f*k \hfill a+f \hfill 0 \hfill 0\\
x += f \hfill min & x + f \hfill a+f \hfill INF \hfill 0\\
x += f \hfill max & x + f \hfill a+f \hfill -INF \hfill 0\\
x = f \hfill sum & f * k \hfill f \hfill 0 \hfill NONE\\
x = f \hfill min & f \hfill f \hfill INF \hfill NONE\\
x = f \hfill max & f \hfill f \hfill -INF \hfill NONE\\
x \textasciicircum= f \hfill xor & k\&1 ? x\textasciicircum f : x \hfill a\textasciicircum f \hfill 0 \hfill 0\\
\end{ctab}}\\[0.3ex]
\code{NONE = LLONG\_MIN}. Affine (one problem mixing ``multiply a range'' and ``add a
range''): \code{F = pair<ll,ll>}, \code{f = \{p,q\}} means $x \mapsto px+q$,
\code{ap = f.first*x + f.second*k}, \code{cm(a,f) = \{a.first*f.first,
f.first*a.second + f.second\}}, \code{NOP = \{1,0\}}.\\[0.3ex]
Check a new row on paper with
$\mathrm{op}(\mathrm{ap}(a,f,k_a), \mathrm{ap}(b,\mathrm{sh}(f,k_a),k_b)) = \mathrm{ap}(\mathrm{op}(a,b), f, k_a{+}k_b)$
and $\mathrm{ap}(\mathrm{ap}(x,a,k),f,k) = \mathrm{ap}(x,\mathrm{cm}(a,f),k)$.
$\sim$4x slower than \code{Seg} --- use \code{Seg} for point updates.}

%LABEL LazySeg
\begin{lstlisting}
struct LazySeg {
    using T = ll;                            // node value
    using F = ll;                            // pending update
    static T op(T a, T b)        { return a + b; }
    static T ap(T x, F f, int k) { return x + f * k; }
    static F cm(F a, F f)        { return a + f; }
    static F sh(F f, int k)      { return f; }
    static constexpr T ID  = 0;
    static constexpr F NOP = 0;
    int n; vector<T> t; vector<F> d;
    LazySeg(int n = 0) : n(n), t(4*n, ID), d(4*n, NOP) {}
    LazySeg(const vector<T>& a) : LazySeg((int)a.size()) {
        if (n) build(1, 0, n, a);
    }
    void update(int l, int r, F f) {
        if (l < r) upd(1, 0, n, l, r, f);
    }
    T query(int l, int r) {
        return l < r ? qry(1, 0, n, l, r) : ID;
    }
    T operator[](int i) { return query(i, i + 1); }
    // node v covers [l, r)
    void build(int v, int l, int r, const vector<T>& a) {
        if (r - l == 1) { t[v] = a[l]; return; }
        int m = (l + r) / 2;
        build(2*v, l, m, a); build(2*v+1, m, r, a);
        t[v] = op(t[2*v], t[2*v+1]);
    }
    void put(int v, int k, F f) {
        t[v] = ap(t[v], f, k); d[v] = cm(d[v], f);
    }
    void push(int v, int l, int m, int r) {
        if (d[v] == NOP) return;
        put(2*v, m - l, d[v]);
        put(2*v+1, r - m, sh(d[v], m - l));
        d[v] = NOP;
    }
    void upd(int v, int l, int r, int ql, int qr, F f) {
        if (ql <= l && r <= qr) { put(v, r - l, f); return; }
        int m = (l + r) / 2;
        push(v, l, m, r);
        if (ql < m) upd(2*v, l, m, ql, qr, f);
        int c = max(0, m - max(ql, l));   // update left of m
        if (m < qr) upd(2*v+1, m, r, ql, qr, sh(f, c));
        t[v] = op(t[2*v], t[2*v+1]);
    }
    T qry(int v, int l, int r, int ql, int qr) {
        if (ql <= l && r <= qr) return t[v];
        int m = (l + r) / 2;
        push(v, l, m, r);
        if (qr <= m) return qry(2*v, l, m, ql, qr);
        if (m <= ql) return qry(2*v+1, m, r, ql, qr);
        return op(qry(2*v, l, m, ql, qr),
                  qry(2*v+1, m, r, ql, qr));
    }
};
\end{lstlisting}

\entry{Implicit treap}{insert / erase / at / query $O(\log n)$ \quad mem $n$}{%
A \emph{sequence} you can insert into and erase from the middle of, by index --- what
\code{vector} cannot do and \code{Seg} cannot do. Indices are positions in the
current sequence, so every insert renumbers everything after it.\\[0.3ex]
\code{agg} is the subtree fold; change the \code{+} in \code{pull} and the \code{T\{\}}
in \code{ag} together to get min/max instead of sum. \code{query(l, r)} folds
\code{[l, r)}.\\[0.3ex]
\code{Treap<ll> t; t.insert(0, 5); t.insert(1, 7); t.at(0); t.erase(0);}}

\begin{lstlisting}
mt19937 rng(chrono::steady_clock::now()
                .time_since_epoch().count());
template<typename T>
struct Treap {
    struct Node {
        T val, agg; ll prio, sz; Node *l, *r;
        Node(T v) : val(v), agg(v), prio(rng()), sz(1),
                    l(nullptr), r(nullptr) {}
    };
    Node* root = nullptr;
    static ll sz(Node* t) { return t ? t->sz : 0; }
    static T  ag(Node* t) { return t ? t->agg : T{}; }
    static void pull(Node* t) {
        if (!t) return;
        t->sz  = 1 + sz(t->l) + sz(t->r);
        t->agg = ag(t->l) + t->val + ag(t->r);   // op
    }
    // first k elements -> a, rest -> b
    static void split(Node* t, ll k, Node*& a, Node*& b) {
        if (!t) { a = b = nullptr; return; }
        if (sz(t->l) >= k) { split(t->l, k, a, t->l); b = t; }
        else { split(t->r, k-sz(t->l)-1, t->r, b); a = t; }
        pull(t);
    }
    static Node* merge(Node* a, Node* b) {
        if (!a || !b) return a ? a : b;
        if (a->prio > b->prio) {
            a->r = merge(a->r, b); pull(a); return a;
        }
        b->l = merge(a, b->l); pull(b); return b;
    }
    void insert(ll k, T val) {       // ends up at index k
        Node *a, *b;
        split(root, k, a, b);
        root = merge(merge(a, new Node(val)), b);
    }
    void erase(ll k) {
        Node *a, *b, *c;
        split(root, k, a, b); split(b, 1, b, c);
        delete b; root = merge(a, c);
    }
    T at(ll k) {
        Node* t = root;
        while (t) {
            ll ls = sz(t->l);
            if (k == ls) return t->val;
            if (k <  ls) t = t->l;
            else { k -= ls + 1; t = t->r; }
        }
        assert(false);
    }
    T query(ll l, ll r) {                  // fold [l, r)
        Node *a, *b, *c;
        split(root, l, a, b); split(b, r - l, b, c);
        T res = ag(b);
        root = merge(a, merge(b, c));
        return res;
    }
    ll size() const { return sz(root); }
};
\end{lstlisting}

\entry{Order-statistic tree}{insert / erase / rank / kth $O(\log n)$}{%
A \code{set} that also answers ``how many are below $x$'' and ``what is the $k$-th''.
\code{find\_by\_order(k)} is an iterator to the \code{k}-th smallest (0-indexed);
\code{order\_of\_key(x)} counts elements \emph{strictly} below \code{x}.\\[0.3ex]
It is a \emph{set}, so duplicates collapse. For a multiset store
\code{pair<val, unique\_id>} and query with \code{\{x, 0\}} / \code{\{x, INF\}}.}

\begin{lstlisting}
#include <ext/pb_ds/assoc_container.hpp>
using namespace __gnu_pbds;
template<class T> using Tree = tree<T, null_type,
    less<T>, rb_tree_tag,
    tree_order_statistics_node_update>;
// Tree<ll> s; s.insert(4);
// *s.find_by_order(0);   s.order_of_key(4);
\end{lstlisting}

\entry{Line container (CHT)}{add / query $O(\log n)$ \quad mem $n$}{%
Upper hull of lines: \code{add(k, m)} inserts $y = kx+m$, \code{query(x)} returns the
\emph{maximum} over all inserted lines at $x$. Lines may arrive in any order.
For a minimum, insert \code{(-k, -m)} and negate the answer.\\[0.3ex]
This is the tool for DP transitions of the form
$dp_i = \max_j (dp_j + a_j \cdot b_i)$ --- insert line
\code{(a[j], dp[j])} once $dp_j$ is known, then \code{query(b[i])}.
Integer arithmetic only; \code{div} floors towards $-\infty$.}

\begin{lstlisting}
struct Line {
    mutable ll k, m, p;
    bool operator<(const Line& o) const { return k < o.k; }
    bool operator<(ll x) const { return p < x; }
};
struct LineContainer : multiset<Line, less<>> {
    static const ll LINF = LLONG_MAX;
    ll fdiv(ll a, ll b) {      // floored division
        return a / b - ((a ^ b) < 0 && a % b);
    }
    bool isect(iterator x, iterator y) {
        if (y == end()) return x->p = LINF, 0;
        if (x->k == y->k) x->p = x->m > y->m ? LINF : -LINF;
        else x->p = fdiv(y->m - x->m, x->k - y->k);
        return x->p >= y->p;
    }
    void add(ll k, ll m) {
        auto z = insert({k, m, 0}), y = z++, x = y;
        while (isect(y, z)) z = erase(z);
        if (x != begin() && isect(--x, y))
            isect(x, y = erase(y));
        while ((y = x) != begin() && (--x)->p >= y->p)
            isect(x, erase(y));
    }
    ll query(ll x) {                   // max over lines
        assert(!empty());
        auto l = *lower_bound(x);
        return l.k * x + l.m;
    }
};
\end{lstlisting}

\entry{Mo's algorithm}{KACTL \quad $O((N + Q)\sqrt N)$ add/del calls}{%
Offline range queries where you can cheaply \textbf{add or remove one element} of the
current window but cannot merge two halves: number of distinct values, most frequent
value, number of equal pairs, $\sum \mathrm{cnt}^2$. If a segment tree can do it, use
the segment tree.\\[0.3ex]
Queries are \textbf{half-open} \code{\{l, r\}} $= [l, r)$. Read all of them first;
\code{mo} returns the answers in input order. It sorts queries by block of \code{l}
(alternating the \code{r} direction per block), so the two pointers move
$O(N\sqrt Q)$ in total. Set \code{blk} $\approx N/\sqrt Q$.\\[0.3ex]
You write \code{add}, \code{del}, \code{calc}; the example below counts distinct values
(compress \code{a} first, size \code{cnt} to the number of values). \code{end} says
which side moved (0 left, 1 right) --- only needed when order matters. KACTL also has a
tree-path version (\code{moTree}).}

%LABEL mo
\begin{lstlisting}
vi a, cnt; ll nd = 0;   // example: count distinct values
void add(int i, int end) { if (cnt[a[i]]++ == 0) nd++; }
void del(int i, int end) { if (--cnt[a[i]] == 0) nd--; }
ll calc() { return nd; }
vi mo(vector<pii> Q) {
    int L = 0, R = 0, blk = 350; // ~N/sqrt(Q)
    vi s(Q.size()), res = s;
#define K(x) pii(x.first/blk, x.second ^ -(x.first/blk & 1))
    iota(all(s), 0);
    sort(all(s), [&](int s, int t){
        return K(Q[s]) < K(Q[t]); });
    for (int qi : s) {
        pii q = Q[qi];
        while (L > q.first) add(--L, 0);
        while (R < q.second) add(R++, 1);
        while (L < q.first) del(L++, 0);
        while (R > q.second) del(--R, 1);
        res[qi] = calc();
    }
    return res;
}
\end{lstlisting}

\section{Union-find and MST}

\entry{DSU}{$O(\alpha(n))$ amortised \quad mem $n$}{%
\code{p[a] < 0} marks a root and then \code{-p[a]} is its component size, otherwise
\code{p[a]} is the parent. Union by size plus path compression.}

%LABEL DSU
\begin{lstlisting}
struct DSU {
    vector<int> p;
    DSU(int n = 0) : p(n, -1) {}
    int find(int a) { return p[a]<0 ? a : p[a]=find(p[a]); }
    int size(int a) { return -p[find(a)]; }
    bool same(int a, int b) { return find(a) == find(b); }
    bool join(int a, int b) {   // false if already joined
        a = find(a); b = find(b);
        if (a == b) return false;
        if (p[a] > p[b]) swap(a, b);         // a = bigger
        p[a] += p[b]; p[b] = a;
        return true;
    }
};
\end{lstlisting}

\entry{DSU with rollback}{join / find $O(\log n)$, undo $O(1)$ \quad mem $n+q$}{%
Union-find where every join can be taken back. Path compression is dropped because it
cannot be undone, so \code{find} costs a log. \code{snap()} records the current state,
\code{undo(t)} restores it exactly.\\[0.3ex]
Use it when merges happen on the way \emph{down} a recursion and must be un-merged on
the way back up: offline dynamic connectivity (segment tree over time, see
\code{Trees}), or a batch of MST/connectivity queries that each add a different set of
edges on top of a common prefix.}

\begin{lstlisting}
struct RDSU {
    vector<int> p;
    vector<pair<int,int>> hist;  // (joined root, its old p)
    RDSU(int n = 0) : p(n, -1) {}
    int find(int a) { while (p[a] >= 0) a = p[a]; return a; }
    int size(int a) { return -p[find(a)]; }
    bool same(int a, int b) { return find(a) == find(b); }
    bool join(int a, int b) {
        a = find(a); b = find(b);
        if (a == b) return false;
        if (p[a] > p[b]) swap(a, b);
        hist.push_back({b, p[b]});
        p[a] += p[b]; p[b] = a;
        return true;
    }
    int snap() { return (int)hist.size(); }
    void undo(int t) {
        while ((int)hist.size() > t) {
            auto [b, old] = hist.back(); hist.pop_back();
            p[p[b]] -= old; p[b] = old;
        }
    }
};
\end{lstlisting}

\entry{Weighted DSU (potentials)}{$O(\log n)$ \quad mem $n$}{%
Union-find that also remembers \emph{differences}. \code{join(a, b, w)} asserts
$w(a) - w(b) = w$ and returns \code{false} if that contradicts what is already known
--- so a run of \code{join} calls decides whether a system of
``$x_i - x_j = c$'' equations is consistent. \code{diff(a, b)} gives
$w(a) - w(b)$, meaningful only when \code{same(a, b)}.\\[0.3ex]
For $\bmod\ k$ versions (``these two are on opposite sides'') use $k = 2$ and keep
every weight in $\{0,1\}$: that is bipartiteness under edge insertions.}

%LABEL WDSU
\begin{lstlisting}
struct WDSU {
    vector<int> p, sz; vector<ll> d;  // d[x]=w(x)-w(root)
    WDSU(int n = 0) : p(n), sz(n, 1), d(n, 0) {
        iota(all(p), 0);
    }
    int find(int a) {
        if (p[a] == a) return a;
        int r = find(p[a]);
        d[a] += d[p[a]]; p[a] = r;     // parent first!
        return r;
    }
    ll pot(int a) { find(a); return d[a]; }
    ll diff(int a, int b) { return pot(a) - pot(b); }
    bool same(int a, int b) { return find(a) == find(b); }
    bool join(int a, int b, ll w) {    // w(a) - w(b) == w
        ll da = pot(a), db = pot(b);
        int ra = find(a), rb = find(b);
        if (ra == rb) return da - db == w;
        if (sz[ra] < sz[rb]) {
            swap(ra, rb); swap(da, db); w = -w;
        }
        p[rb] = ra; sz[ra] += sz[rb];
        d[rb] = da - db - w;
        return true;
    }
};
\end{lstlisting}

\entry{Kruskal}{$O(m \log m)$ \quad mem $n + m$}{%
Needs \code{DSU}. Sort by weight, keep an edge iff it joins two components.
\code{mst.size() < n-1} means the graph was disconnected.\\[0.3ex]
What you usually want is not the tree but the \emph{join order}: after processing all
edges of weight $\le w$ the components are exactly the connected components at that
threshold.\\[0.3ex]
The max edge on the MST path between $a$ and $b$ equals the minimum, over all
$a$--$b$ paths, of the max edge on it (minimax path) --- answer those with \code{LCA}
carrying a max table, not with a new search. For ``is this edge in some MST'' and for
the second-best MST, process all edges of equal weight as one block.}

%DEP DSU
\begin{lstlisting}
// e = {w, a, b}. returns {total cost, chosen edges}
pair<ll, vector<array<ll,3>>>
kruskal(int n, vector<array<ll,3>> e) {
    sort(all(e));
    DSU d(n); ll cost = 0; vector<array<ll,3>> mst;
    for (auto [w, a, b] : e)
        if (d.join(a, b)) {
            cost += w; mst.push_back({w, a, b});
        }
    return {cost, mst};
}
\end{lstlisting}

\section{Trees}

\entry{Euler tour}{$O(n)$ \quad mem $n$}{%
The subtree of \code{u} is the range \code{[tin[u], tout[u])} in any array indexed by
\code{tin}, and \code{u} is an ancestor of \code{v} iff
\code{tin[u] <= tin[v] < tout[u]}. Turns ``update / query a whole subtree'' into a
range operation, so pair it with \code{BIT} / \code{Seg} / \code{LazySeg}. Edge values
go on the \emph{lower} endpoint of the edge.\\[0.3ex]
Iterative, so depth never matters.}

\begin{lstlisting}
void euler(const vector<vector<int>>& g, int root,
           vector<int>& tin, vector<int>& tout) {
    int n = (int)g.size(), t = 0;
    tin.assign(n, 0); tout.assign(n, 0);
    vector<int> par(n, -1), it(n, 0), st{root};
    tin[root] = t++;
    while (!st.empty()) {
        int u = st.back();
        if (it[u] < (int)g[u].size()) {
            int v = g[u][it[u]++];
            if (v == par[u]) continue;
            par[v] = u; tin[v] = t++; st.push_back(v);
        } else { tout[u] = t; st.pop_back(); }
    }
}
\end{lstlisting}

\entry{LCA (binary lifting)}{build $O(n \log n)$, query $O(\log n)$ \quad mem $n \log n$}{%
\code{up[k][v]} is the $2^k$-th ancestor of \code{v}; the root is its own ancestor so
jumps past it clamp there. Lift the deeper vertex to the other's depth, then jump both
up together while the ancestors differ. Built by BFS, so no recursion depth
trouble.\\[0.3ex]
Max / min / sum edge on a path: add a second table \code{mx} with \code{mx[0][v]} =
weight of the edge to the parent, \code{mx[0][root] = ID},
\code{mx[k][v] = op(mx[k-1][v], mx[k-1][up[k-1][v]])}, and merge \code{mx[i][v]} into
the answer at every jump taken in \code{kth} and \code{lca}.}

\begin{lstlisting}
struct LCA {
    int LG; vector<vector<int>> up; vector<int> dep;
    LCA(const vector<vector<int>>& g, int root = 0) {
        int n = g.size(); LG = 1 + __lg(max(n, 2));
        up.assign(LG, vector<int>(n, root)); dep.assign(n, 0);
        vector<int> order{root};
        vector<bool> seen(n); seen[root] = true;
        for (int i = 0; i < (int)order.size(); i++) {
            int u = order[i];
            for (int v : g[u]) if (!seen[v]) {
                seen[v] = true; up[0][v] = u;
                dep[v] = dep[u] + 1; order.push_back(v);
            }
        }
        for (int k = 1; k < LG; k++)
            for (int v = 0; v < n; v++)
                up[k][v] = up[k-1][up[k-1][v]];
    }
    int kth(int v, int k) {        // root if k is too big
        k = min(k, dep[v]);
        for (int i = 0; k; k >>= 1, i++)
            if (k & 1) v = up[i][v];
        return v;
    }
    int lca(int a, int b) {
        if (dep[a] < dep[b]) swap(a, b);
        a = kth(a, dep[a] - dep[b]);
        if (a == b) return a;
        down(k, LG, 0) if (up[k][a] != up[k][b])
            a = up[k][a], b = up[k][b];
        return up[0][a];
    }
    int dist(int a, int b) {
        return dep[a] + dep[b] - 2 * dep[lca(a, b)];
    }
};
\end{lstlisting}

\entry{Heavy-light decomposition}{build $O(n)$, path / subtree $O(\log^2 n)$ \quad mem $n$}{%
Every vertex points at its heaviest child, so a root-to-leaf walk changes chain at most
$\log n$ times. Positions are handed out so that each chain and each subtree is a
contiguous range; feed those to \code{Seg} / \code{LazySeg} / \code{BIT}, which is the
second log. \code{lca} comes for free, no extra table.\\[0.3ex]
\code{path} calls \code{f} on $O(\log n)$ ranges in \emph{arbitrary} order, so
\code{op} must commute. For a non-commutative op keep two trees, one on \code{pos} and
one on reversed \code{pos}, and collect the two sides of the LCA separately.
\code{edges = true} stores each edge at its lower endpoint and skips the LCA itself;
then \code{sub} should start at \code{pos[u] + 1} too. Recursion depth is $O(n)$ ---
on a path of $2\cdot10^5$ vertices raise the stack or rewrite the two \code{dfs}
calls iteratively.\\[0.3ex]
\code{HLD h(g); Seg s(n); s.set(h.pos[u], val[u]);}\\
\code{h.path(a, b, [\&](int l, int r)\{ ... s.query(l, r) ... \});}}

\begin{lstlisting}
struct HLD {
    int n, t = 0;
    vector<vector<int>> g;      // heavy child kept first
    vector<int> par, dep, sz, head, pos;
    HLD(vector<vector<int>> adj, int root = 0)
        : n(adj.size()), g(move(adj)), par(n, -1), dep(n),
          sz(n), head(n, root), pos(n) {
        dfs(root); build(root, root);
    }
    void dfs(int u) {
        if (par[u] != -1) g[u].erase(find(all(g[u]), par[u]));
        sz[u] = 1;
        for (int& v : g[u]) {
            par[v] = u; dep[v] = dep[u] + 1;
            dfs(v); sz[u] += sz[v];
            if (sz[v] > sz[g[u][0]]) swap(v, g[u][0]);
        }
    }
    void build(int u, int h) {      // heavy child keeps head
        head[u] = h; pos[u] = t++;
        for (int i = 0; i < (int)g[u].size(); i++)
            build(g[u][i], i ? g[u][i] : h);
    }
    int lca(int a, int b) {
        for (; head[a] != head[b]; b = par[head[b]])
            if (dep[head[a]] > dep[head[b]]) swap(a, b);
        return dep[a] < dep[b] ? a : b;
    }
    pair<int,int> sub(int u) {              // subtree range
        return {pos[u], pos[u] + sz[u]};
    }
    template<class F>
    void path(int a, int b, F f, bool edges = false) {
        for (; head[a] != head[b]; b = par[head[b]]) {
            if (dep[head[a]] > dep[head[b]]) swap(a, b);
            f(pos[head[b]], pos[b] + 1);
        }
        if (dep[a] > dep[b]) swap(a, b);
        f(pos[a] + edges, pos[b] + 1);
    }
};
\end{lstlisting}

\entry{Tree diameter}{$O(n)$ \quad mem $n$}{%
Two BFS runs: the farthest vertex from anywhere is an endpoint of some diameter, and
the farthest vertex from \emph{that} is the other endpoint. Unweighted only --- with
weights use the same shape with Dijkstra, or a single DFS keeping the two deepest
children per vertex. Iterative.}

\begin{lstlisting}
// returns {length in edges, endpoint a, endpoint b}
array<int,3> diameter(const vector<vector<int>>& g) {
    int n = (int)g.size();
    auto bfs = [&](int s) {
        vector<int> d(n, -1); d[s] = 0;
        vector<int> q{s};
        for (int i = 0; i < (int)q.size(); i++)
            for (int v : g[q[i]]) if (d[v] < 0) {
                d[v] = d[q[i]] + 1; q.push_back(v);
            }
        int b = s;
        for (int i = 0; i < n; i++)
            if (d[i] > d[b]) b = i;
        return pair<int,int>{b, d[b]};
    };
    int a = bfs(0).first;
    auto [b, len] = bfs(a);
    return {len, a, b};
}
\end{lstlisting}

\entry{Centroid}{$O(n)$ \quad mem $n$}{%
The centroid minimises the largest component left after deleting it; that component is
always $\le n/2$. Every tree has one or two, and two only when they are
adjacent.\\[0.3ex]
Recurse on the pieces (removing the centroid each time) and you get centroid
decomposition: $O(\log n)$ levels, every path in the tree passes through the centroid
of exactly one level, which is how ``count paths with property X'' problems are
solved in $O(n \log n)$.}

\begin{lstlisting}
vector<int> centroids(const vector<vector<int>>& g) {
    int n = (int)g.size();
    vector<int> sub(n, 1), par(n, -1), it(n, 0), st{0}, ord;
    while (!st.empty()) {              // iterative dfs order
        int u = st.back();
        if (it[u] < (int)g[u].size()) {
            int v = g[u][it[u]++];
            if (v == par[u]) continue;
            par[v] = u; st.push_back(v);
        } else { ord.push_back(u); st.pop_back(); }
    }
    for (int u : ord)                        // children first
        if (par[u] >= 0) sub[par[u]] += sub[u];
    vector<int> res;
    for (int u = 0; u < n; u++) {
        int mx = n - sub[u];
        for (int v : g[u]) if (v != par[u])
            mx = max(mx, sub[v]);
        if (mx <= n / 2) res.push_back(u);
    }
    return res;
}
\end{lstlisting}

\entry{Rerooting DP}{$O(n)$ \quad mem $n$}{%
Answers a subtree DP for \emph{every} root at once, in two passes: up to collect each
subtree, then down to hand each child what the rest of the tree contributes. Shown for
$\mathrm{ans}[u] = \sum_v \mathrm{dist}(u,v)$.\\[0.3ex]
The pattern generalises: whatever \code{dp} accumulates from children, in the second
pass give child \code{v} the value ``everything except \code{v}'s subtree'', which for
a sum is \code{dp[u] - contribution(v)} and for a max needs the top two children.}

\begin{lstlisting}
// ans[u] = sum over v of dist(u, v)
vi sumOfDists(const vector<vector<int>>& g) {
    int n = (int)g.size();
    vi sub(n, 1), dp(n, 0), ans(n, 0);
    vector<int> par(n, -1), it(n, 0), st{0}, ord;
    while (!st.empty()) {
        int u = st.back();
        if (it[u] < (int)g[u].size()) {
            int v = g[u][it[u]++];
            if (v == par[u]) continue;
            par[v] = u; st.push_back(v);
        } else { ord.push_back(u); st.pop_back(); }
    }
    for (int u : ord) if (par[u] >= 0) {     // up: children
        sub[par[u]] += sub[u];
        dp[par[u]]  += dp[u] + sub[u];
    }
    ans[ord.back()] = dp[ord.back()];        // root
    down(i, (ll)ord.size(), 0) {             // down: parents
        int u = ord[i];
        if (par[u] >= 0)
            ans[u] = ans[par[u]] + n - 2 * sub[u];
    }
    return ans;
}
\end{lstlisting}

\entry{Offline DFS + rollback}{$O((n + q) \cdot \mathrm{apply})$ \quad mem $n + q$}{%
Not a structure, a shape. Bucket every query at the vertex it belongs to, then run
\emph{one} DFS from the root holding a global structure that always describes the
current root-to-\code{u} path. Apply \code{u}'s updates on entry, answer \code{u}'s
queries, recurse, and on exit undo exactly what was applied --- so the structure is
correct at every vertex without ever being rebuilt. Needs invertible updates and
offline queries.\\[0.3ex]
\begin{ctab}{0.33}{0.62}
array / counter & subtract what you added\\
diff array by depth & \code{d[dep] += x}, \code{d[dep+len+1] -= x}, prefix-summed
along the path\\
\code{BIT} & \code{add(i, x)} then \code{add(i, -x)}\\
\code{RDSU} & \code{snap()} then \code{undo(t)}\\
\end{ctab}\\[0.3ex]
Replace the recursion over the tree with a recursion over a segment tree of
\emph{time} and the same shape becomes offline dynamic connectivity: put each edge on
the $O(\log q)$ nodes covering the interval it is alive, DFS that tree, \code{join} on
entry, \code{undo} on exit. That is the only place \code{RDSU} is really needed ---
with an array or a BIT the undo is just the negated update.}

%NOCOMPILE
\begin{lstlisting}
void dfs(int u, int par) {
    int t = snapshot();                // e.g. rdsu.snap()
    for (auto& upd : at[u]) apply(upd);          // enter
    for (auto& q : qs[u]) ans[q.id] = get(q);
    for (int v : g[u]) if (v != par) dfs(v, u);
    undo(t);                     // exit: exact inverse
}
\end{lstlisting}

\section{Connectivity}

\entry{SCC (Tarjan, iterative)}{$O(n + m)$ \quad mem $n + m$}{%
Directed. Takes any adjacency of an integer type. No recursion, so depth never matters;
self-loops and multi-edges are fine.\\[0.3ex]
\code{comp} is in \textbf{topological} order --- every edge \code{u -> v} has
\code{comp[u] <= comp[v]} --- so a DP over the condensation is just a loop over
\code{0 .. nc-1}, no toposort. \code{members()} gives the vertices of each comp,
\code{dag()} the condensation with duplicates removed.\\[0.3ex]
A comp holds a cycle iff it has $\ge 2$ vertices, or 1 with a self-loop; the graph is a
DAG iff \code{nc == n} and no self-loop exists. Sources of the condensation have
in-degree 0 in \code{dag()}, sinks out-degree 0 --- count degrees, \code{comp[v] == 0}
is not enough.\\[0.3ex]
\textbf{2-SAT.} Node \code{2i} is $x_i$ false, \code{2i+1} is true, so \code{n = 2*vars}.
A clause $(a \lor b)$ adds \emph{both} \code{!a -> b} and \code{!b -> a}; forcing $a$
true is the single edge \code{!a -> a}. Satisfiable iff
\code{comp[2i] != comp[2i+1]} for every $i$, and then
$x_i$ = \code{comp[2i+1] > comp[2i]} (needs \code{comp} topological, as here).}

%LABEL SCC
\begin{lstlisting}
struct SCC {
    int n, nc = 0;              // nc = number of comps
    vector<vector<int>> g;
    vector<int> comp;           // topological order
    template<class G>          // any vector<vector<intlike>>
    SCC(const G& adj) : n(adj.size()), g(n), comp(n, -1) {
        for (int u = 0; u < n; u++)
            for (auto v : adj[u]) g[u].push_back((int)v);
        vector<int> tin(n, -1), low(n), ptr(n, 0), stk, cs;
        vector<char> on(n, 0);
        int timer = 0;
        for (int r = 0; r < n; r++) {
            if (tin[r] != -1) continue;
            tin[r] = low[r] = timer++;
            stk.push_back(r); on[r] = 1; cs.push_back(r);
            while (!cs.empty()) {
                int u = cs.back();
                if (ptr[u] < (int)g[u].size()) {
                    int v = g[u][ptr[u]++];
                    if (tin[v] == -1) {
                        tin[v] = low[v] = timer++;
                        stk.push_back(v); on[v] = 1;
                        cs.push_back(v);
                    } else if (on[v])
                        low[u] = min(low[u], tin[v]);
                } else {
                    cs.pop_back();
                    if (!cs.empty()) {
                        int p = cs.back();
                        low[p] = min(low[p], low[u]);
                    }
                    if (low[u] == tin[u]) {  // u roots comp
                        int w;
                        do {
                            w = stk.back(); stk.pop_back();
                            on[w] = 0; comp[w] = nc;
                        } while (w != u);
                        nc++;
                    }
                }
            }
        }
        for (int v = 0; v < n; v++)         // -> topological
            comp[v] = nc - 1 - comp[v];
    }
    vector<vector<int>> members() {
        vector<vector<int>> c(nc);
        for (int v = 0; v < n; v++) c[comp[v]].push_back(v);
        return c;
    }
    vector<vector<int>> dag() {      // condensation, dedup
        vector<vector<int>> d(nc);
        for (int u = 0; u < n; u++)
            for (int v : g[u])
                if (comp[u] != comp[v])
                    d[comp[u]].push_back(comp[v]);
        for (auto& a : d) {
            sort(all(a));
            a.erase(unique(all(a)), a.end());
        }
        return d;
    }
};
\end{lstlisting}

\entry{Bridges, cut vertices, 2ECC}{$O(n + m)$ \quad mem $n + m$}{%
Undirected, given as an \textbf{edge list} (pairs, 2-element arrays or tuples of any
integer type), so \code{bri[i]} indexes your own input. Multi-edges and self-loops are
safe, disconnected input is fine, no recursion.\\[0.3ex]
\begin{ctab}{0.15}{0.80}
\code{bri[i]} & input edge \code{i} is a bridge.\\
\code{art[v]} & \code{v} is a cut vertex.\\
\code{comp[v]} & 2-edge-connected comp: all you reach without crossing a bridge.\\
\code{tree()} & the bridge forest on those comps.\\
\end{ctab}\\[0.3ex]
Bridges are exactly the edges between comps, so \code{tree()} is a forest, one tree per
connected component --- feed it to \code{LCA} / \code{HLD} for ``how many bridges on
the \code{u}--\code{v} path''. Adding one edge kills every bridge on the path joining
its two comps.\\[0.3ex]
The parent edge is skipped by \textbf{edge id}, never by vertex --- by vertex reports
parallel edges as bridges, which is the usual wrong answer. A bridge endpoint need not
be a cut vertex (degree 1) and a cut vertex need not carry one (two cycles sharing it);
the DFS root is one iff it has $\ge 2$ children.}

\begin{lstlisting}
struct Bridges {
    int n, m, nc = 0;                  // nc = number of 2ecc
    vector<vector<pair<int,int>>> g;   // (to, edge id)
    vector<int> comp;                  // 2ecc id per vertex
    vector<char> bri, art;             // per edge / vertex
    vector<pair<int,int>> es;
    template<class E>      // list of pair / array / tuple
    Bridges(int n, const E& edges)
            : n(n), m(edges.size()), g(n), comp(n, -1),
              bri(m, 0), art(n, 0) {
        es.reserve(m);
        for (auto& e : edges)
            es.push_back({(int)get<0>(e), (int)get<1>(e)});
        vector<int> deg(n, 0);
        for (int i = 0; i < m; i++) {
            deg[es[i].first]++; deg[es[i].second]++;
        }
        for (int v = 0; v < n; v++) g[v].reserve(deg[v]);
        for (int i = 0; i < m; i++) {
            g[es[i].first].push_back({es[i].second, i});
            g[es[i].second].push_back({es[i].first, i});
        }
        vector<int> tin(n, -1), low(n), ptr(n, 0), cs;
        vector<int> pe(n, -1);
        int timer = 0;
        for (int r = 0; r < n; r++) {
            if (tin[r] != -1) continue;
            int kids = 0;
            tin[r] = low[r] = timer++; cs.push_back(r);
            while (!cs.empty()) {
                int u = cs.back();
                if (ptr[u] < (int)g[u].size()) {
                    auto [v, id] = g[u][ptr[u]++];
                    if (id == pe[u]) continue;   // skip by ID
                    if (tin[v] == -1) {
                        pe[v] = id; tin[v] = low[v] = timer++;
                        cs.push_back(v); kids += (u == r);
                    } else low[u] = min(low[u], tin[v]);
                } else {
                    cs.pop_back();
                    if (cs.empty()) break;
                    int p = cs.back();
                    low[p] = min(low[p], low[u]);
                    if (low[u] > tin[p]) bri[pe[u]] = 1;
                    if (low[u] >= tin[p]) art[p] = 1;
                }
            }
            art[r] = kids >= 2;                  // root rule
        }
        vector<int> q;                // 2ecc: drop bridges
        for (int r = 0; r < n; r++) {
            if (comp[r] != -1) continue;
            comp[r] = nc; q.clear(); q.push_back(r);
            while (!q.empty()) {
                int u = q.back(); q.pop_back();
                for (auto [v, id] : g[u])
                    if (!bri[id] && comp[v] == -1) {
                        comp[v] = nc; q.push_back(v);
                    }
            }
            nc++;
        }
    }
    Bridges(int n, initializer_list<pair<int,int>> e)
            : Bridges(n, vector<pair<int,int>>(e)) {}
    vector<vector<int>> tree() {     // bridge forest
        vector<vector<int>> t(nc);
        for (int i = 0; i < m; i++) if (bri[i]) {
            int a = comp[es[i].first], b = comp[es[i].second];
            t[a].push_back(b); t[b].push_back(a);
        }
        return t;
    }
};
\end{lstlisting}

\entry{Topological sort (Kahn)}{$O(n + m)$ \quad mem $n$}{%
Directed. Returns an empty vector iff the graph has a cycle, so it doubles as a cycle
test. If you already built \code{SCC} you have a topological order for free in
\code{comp} --- use this only when you do not need the components.\\[0.3ex]
Pop from the back instead of the front and you get a DFS-flavoured order; use a
\code{priority\_queue} for the lexicographically smallest one.}

\begin{lstlisting}
vector<int> topo(const vector<vector<int>>& g) {
    int n = (int)g.size();
    vector<int> deg(n, 0), q;
    for (int u = 0; u < n; u++)
        for (int v : g[u]) deg[v]++;
    for (int u = 0; u < n; u++) if (!deg[u]) q.push_back(u);
    for (int i = 0; i < (int)q.size(); i++)
        for (int v : g[q[i]])
            if (--deg[v] == 0) q.push_back(v);
    return (int)q.size() == n ? q : vector<int>{};
}
\end{lstlisting}

\entry{2-SAT}{KACTL \quad $O(N + E)$}{%
The ready-made version of the 2-SAT note under \emph{SCC}. Recognise it: every item
has \textbf{two choices} (on/off, left/right, colour A/B) and every constraint mentions
\textbf{at most two} items (``not both'', ``at least one of'', ``same'', ``if this then
that'').\\[0.3ex]
Variables are \code{0..N-1}; \code{\textasciitilde x} is ``x is false''.
\code{either(a, b)} is $a \lor b$, so: implication $a \Rightarrow b$ is
\code{either(\textasciitilde a, b)}; not both is
\code{either(\textasciitilde a, \textasciitilde b)}; $a = b$ is
\code{either(\textasciitilde a, b)} $+$ \code{either(a, \textasciitilde b)}.
\code{setValue(x)} forces x true. \code{atMostOne(\{..\})} adds helper variables so it
stays linear --- do not write the $O(k^2)$ pairs.\\[0.3ex]
\code{solve()} returns satisfiable; \code{values[i]} is then an assignment. Recursive
DFS: depth up to $2N$, fine to $\sim10^5$ variables.}

%LABEL twosat
\begin{lstlisting}
struct TwoSat {
    int N;
    vector<vi> gr;
    vi values; // 0 = false, 1 = true
    TwoSat(int n = 0) : N(n), gr(2*n) {}
    int addVar() { // (optional)
        gr.emplace_back();
        gr.emplace_back();
        return N++;
    }
    void either(int f, int j) {
        f = max(2*f, -1-2*f);
        j = max(2*j, -1-2*j);
        gr[f].push_back(j^1);
        gr[j].push_back(f^1);
    }
    void setValue(int x) { either(x, x); }
    void atMostOne(const vi& li) { // (optional)
        if (li.size() <= 1) return;
        int cur = ~li[0];
        rep(i,2,(ll)li.size()) {
            int next = addVar();
            either(cur, ~li[i]);
            either(cur, next);
            either(~li[i], next);
            cur = ~next;
        }
        either(cur, ~li[1]);
    }
    vector<int> val, comp, z; int time = 0;
    int dfs(int i) {
        int low = val[i] = ++time, x; z.push_back(i);
        for(int e : gr[i]) if (!comp[e])
            low = min(low, val[e] ?: dfs(e));
        if (low == val[i]) do {
            x = z.back(); z.pop_back();
            comp[x] = low;
            if (values[x>>1] == -1)
                values[x>>1] = x&1;
        } while (x != i);
        return val[i] = low;
    }
    bool solve() {
        values.assign(N, -1);
        val.assign(2*N, 0); comp = val;
        rep(i,0,2*N) if (!comp[i]) dfs(i);
        rep(i,0,N) if (comp[2*i] == comp[2*i+1]) return 0;
        return 1;
    }
};
\end{lstlisting}

\entry{Maximum clique}{KACTL \quad $\sim$1 s for $n = 155$, dense}{%
Largest set of pairwise adjacent vertices, exponential but heavily pruned (greedy
colouring bounds). Input: symmetric \code{bitset} adjacency, \textbf{no self-loops};
raise \code{200} to your $n$. Returns the vertices.\\[0.3ex]
\textbf{Maximum independent set} of a general graph = maximum clique of the
\emph{complement} (flip every off-diagonal bit). Bipartite graphs do not need this ---
see the vertex cover entry under matching. Typical bounds: $n \le 40$--$150$, ``pick as
many as possible, no two conflicting''.\\[0.3ex]
KACTL uses \code{sz} everywhere; it is defined and undefined around the struct so
\code{Treap} still compiles.}

%LABEL clique
\begin{lstlisting}
#define sz(x) (int)(x).size()
typedef vector<bitset<200>> vb;
struct Maxclique {
    double limit=0.025, pk=0;
    struct Vertex { int i, d=0; };
    typedef vector<Vertex> vv;
    vb e;
    vv V;
    vector<vi> C;
    vi qmax, q, S, old;
    void init(vv& r) {
        for (auto& v : r) v.d = 0;
        for (auto& v : r) for (auto j : r) v.d += e[v.i][j.i];
        sort(all(r), [](auto a, auto b) { return a.d>b.d; });
        int mxD = r[0].d;
        rep(i,0,sz(r)) r[i].d = min((int)i, mxD) + 1;
    }
    void expand(vv& R, int lev = 1) {
        S[lev] += S[lev - 1] - old[lev];
        old[lev] = S[lev - 1];
        while (sz(R)) {
            if (sz(q) + R.back().d <= sz(qmax)) return;
            q.push_back(R.back().i);
            vv T;
            for(auto v:R) if (e[R.back().i][v.i])
                T.push_back({v.i});
            if (sz(T)) {
                if (S[lev]++ / ++pk < limit) init(T);
                int j = 0, mxk = 1;
                int mnk = max(sz(qmax) - sz(q) + 1, 1);
                C[1].clear(), C[2].clear();
                for (auto v : T) {
                    int k = 1;
                    auto f = [&](int i) { return e[v.i][i]; };
                    while (any_of(all(C[k]), f)) k++;
                    if (k > mxk) mxk = k, C[mxk + 1].clear();
                    if (k < mnk) T[j++].i = v.i;
                    C[k].push_back(v.i);
                }
                if (j > 0) T[j - 1].d = 0;
                rep(k,mnk,mxk + 1) for (int i : C[k])
                    T[j].i = i, T[j++].d = k;
                expand(T, lev + 1);
            } else if (sz(q) > sz(qmax)) qmax = q;
            q.pop_back(), R.pop_back();
        }
    }
    vi maxClique() { init(V), expand(V); return qmax; }
    Maxclique(vb conn) : e(conn), C(sz(e)+1), S(sz(C)),
        old(S) {
        rep(i,0,sz(e)) V.push_back({(int)i});
    }
};
#undef sz
\end{lstlisting}

\section{Shortest paths}

\entry{Dijkstra}{$O((n + m) \log n)$ \quad mem $n + m$}{%
\textbf{Non-negative} weights only --- one negative edge and the answer is silently
wrong, use \code{bellman}. \code{d[v] == INF} means unreachable. Pass \code{\&par} to
recover a path: walk \code{par} back from the target and reverse.\\[0.3ex]
Variants that are the same loop with a different relax step: maximise the bottleneck
(\code{max(d[u], w)} with a max-heap), count shortest paths (a second array,
\code{cnt[v] += cnt[u]} on tie), fewest edges among shortest paths (compare
\code{pair<dist, edges>}). For 0/1 weights use a \code{deque} instead of the heap
(push\_front on 0) and drop the log.}

%LABEL dijkstra
\begin{lstlisting}
// g[u] = {(v, w)}
vi dijkstra(const vector<vector<pii>>& g, int s,
            vector<int>* par = nullptr) {
    int n = (int)g.size();
    vi d(n, INF);
    if (par) par->assign(n, -1);
    priority_queue<pii, vector<pii>, greater<pii>> q;
    d[s] = 0; q.push({0, s});
    while (!q.empty()) {
        auto [dv, u] = q.top(); q.pop();
        if (dv > d[u]) continue;            // stale entry
        for (auto [v, w] : g[u])
            if (dv + w < d[v]) {
                d[v] = dv + w;
                if (par) (*par)[v] = (int)u;
                q.push({d[v], v});
            }
    }
    return d;
}
\end{lstlisting}

\entry{Bellman-Ford}{$O(nm)$ \quad mem $n + m$}{%
Handles negative weights and reports a negative cycle reachable from \code{s}. Stops
early when a round changes nothing, so it is fast on easy inputs.\\[0.3ex]
\textbf{Difference constraints.} A system of $x_v - x_u \le w$ is exactly a shortest
path problem: add edge \code{u -> v} of weight \code{w}, add a super-source with a
0-edge to every vertex, run this from it. The resulting \code{d} is a feasible
assignment, and the system is infeasible iff there is a negative cycle. $x_v - x_u \ge
w$ becomes $x_u - x_v \le -w$; $x_v - x_u = w$ is both.\\[0.3ex]
To mark every vertex whose distance is truly $-\infty$, run $n$ more rounds and set
\code{d[v] = -INF} whenever it still relaxes.}

\begin{lstlisting}
// e = {u, v, w}; neg = a negative cycle is reachable from s
vi bellman(int n, const vector<array<ll,3>>& e, int s,
           bool& neg) {
    vi d(n, INF); d[s] = 0; neg = false;
    rep(i, 0, n) {
        bool ch = false;
        for (auto [u, v, w] : e)
            if (d[u] < INF && d[u] + w < d[v]) {
                d[v] = d[u] + w; ch = true;
                if (i == n - 1) neg = true;
            }
        if (!ch) break;
    }
    return d;
}
\end{lstlisting}

\entry{Floyd-Warshall}{$O(n^3)$ \quad mem $n^2$}{%
All pairs, negative weights allowed. Fill \code{d[i][i] = 0}, \code{d[i][j]} = the
\emph{cheapest} edge \code{i -> j} (careful with multi-edges: take a \code{min}), and
\code{INF} elsewhere. Rewrites \code{d} in place.\\[0.3ex]
A negative cycle exists iff \code{d[i][i] < 0} for some \code{i} afterwards; then
\code{d[a][b]} is $-\infty$ whenever \code{d[a][i] < INF} and \code{d[i][b] < INF} for
such an \code{i}. The \code{k} loop must be outermost. The \code{< INF} guards are what
stop \code{INF + INF} from wrapping.\\[0.3ex]
Same triple loop, other semirings: transitive closure (\code{|=} with \code{\&\&}),
maximum bottleneck (\code{max}/\code{min}), minimax path.}

\begin{lstlisting}
void floyd(matrix& d) {
    int n = (int)d.size();
    rep(k, 0, n) rep(i, 0, n) if (d[i][k] < INF)
        rep(j, 0, n) if (d[k][j] < INF)
            d[i][j] = min(d[i][j], d[i][k] + d[k][j]);
}
\end{lstlisting}

\entry{K shortest walks}{$O((m + k n)\log)$ \quad mem $n + m$}{%
The \code{k} smallest \code{s}$\to$\code{t} lengths in ascending order, \textbf{walks}
--- vertices may repeat, so on a graph with a cycle you get the same path plus loops.
That is what most ``k-th shortest route'' problems want. Non-negative weights.
Fewer than \code{k} results means fewer exist.\\[0.3ex]
For genuinely \emph{simple} k shortest paths you need Yen's algorithm --- much longer,
and rarely what is asked.}

\begin{lstlisting}
vi kShortest(const vector<vector<pii>>& g, int s, int t,
             int k) {
    int n = (int)g.size();
    vector<int> cnt(n, 0); vi res;
    priority_queue<pii, vector<pii>, greater<pii>> q;
    q.push({0, s});
    while (!q.empty() && (int)res.size() < k) {
        auto [d, u] = q.top(); q.pop();
        if (cnt[u] >= k) continue;
        cnt[u]++;
        if (u == t) res.push_back(d);
        for (auto [v, w] : g[u])
            if (cnt[v] < k) q.push({d + w, v});
    }
    return res;
}
\end{lstlisting}

\section{Flow and matching}

\entry{Dinic (max flow)}{$O(n^2 m)$, unit caps $O(m\sqrt m)$ \quad mem $n + m$}{%
\code{addEdge} returns the edge id, and after \code{maxflow} the flow actually pushed
along it is \code{flowOn(id)} (only meaningful when it was added with \code{rcap = 0}).
For an undirected edge pass \code{rcap = cap}.\\[0.3ex]
\begin{ctab}{0.22}{0.73}
\code{cut()} & after \code{maxflow}: \code{c[v]} true iff \code{v} is on the
\code{s} side of a minimum cut. The cut edges are those with \code{c[u] \&\& !c[v]}.\\
\end{ctab}\\[0.3ex]
\textbf{Recipes.} Bipartite matching: source$\to$left cap 1, left$\to$right cap 1,
right$\to$sink cap 1; with unit caps this runs in $O(m\sqrt n)$, i.e. the same bound as
Hopcroft-Karp, so there is no separate template here.
Min vertex cover in a bipartite graph = max matching (K\"onig), and max independent set
= $n$ minus that.
Vertex capacities: split \code{v} into \code{v\_in -> v\_out} with that capacity.
Minimum path cover of a DAG = $n$ minus max matching on the split graph.
Project selection / max closure: source$\to$profit nodes, cost nodes$\to$sink,
prerequisites as \code{INF} edges, answer = total profit $-$ max flow.\\[0.3ex]
\code{dfs} recurses to depth $O(n)$; fine at contest sizes, but it is the thing to
rewrite if $n$ is huge.}

%LABEL Dinic
\begin{lstlisting}
struct Dinic {
    struct E { int to; ll cap; };
    int n; vector<E> e; vector<vector<int>> g;
    vector<int> lvl, it;
    Dinic(int n = 0) : n(n), g(n) {}
    int addEdge(int u, int v, ll cap, ll rcap = 0) {
        g[u].push_back((int)e.size()); e.push_back({v, cap});
        g[v].push_back((int)e.size()); e.push_back({u, rcap});
        return (int)e.size() - 2;               // edge id
    }
    ll dfs(int u, int t, ll f) {
        if (u == t || !f) return f;
        for (int& i = it[u]; i < (int)g[u].size(); i++) {
            int id = g[u][i];
            if (lvl[e[id].to] != lvl[u] + 1 || !e[id].cap)
                continue;
            if (ll d = dfs(e[id].to, t, min(f, e[id].cap))) {
                e[id].cap -= d; e[id ^ 1].cap += d;
                return d;
            }
        }
        return 0;
    }
    ll maxflow(int s, int t) {
        ll fl = 0;
        while (true) {
            lvl.assign(n, -1); lvl[s] = 0;
            vector<int> q{s};
            for (int i = 0; i < (int)q.size(); i++)
                for (int id : g[q[i]])
                    if (e[id].cap && lvl[e[id].to] < 0) {
                        lvl[e[id].to] = lvl[q[i]] + 1;
                        q.push_back(e[id].to);
                    }
            if (lvl[t] < 0) return fl;
            it.assign(n, 0);
            while (ll d = dfs(s, t, INF)) fl += d;
        }
    }
    vector<char> cut() {          // true = s side of min cut
        vector<char> c(n);
        for (int i = 0; i < n; i++) c[i] = lvl[i] >= 0;
        return c;
    }
    ll flowOn(int id) { return e[id ^ 1].cap; }
};
\end{lstlisting}

\entry{Min cost max flow}{$O(\mathrm{flow} \cdot nm)$ \quad mem $n + m$}{%
SPFA-based, so \textbf{negative edge costs are fine} (negative \emph{cycles} are not).
Returns \code{\{flow, cost\}}. Pass \code{maxf} to cap the flow, which is how you ask
for ``cheapest way to ship exactly $k$ units''.\\[0.3ex]
For \emph{min cost flow of any size} rather than min cost \emph{max} flow, stop as soon
as \code{d[t] >= 0}: further augmentation only adds cost.\\[0.3ex]
Assignment problems fit here (left$\to$right with cost, unit caps) but
\code{hungarian} is $O(n^3)$ and much faster for a dense square matrix.}

%LABEL MCMF
\begin{lstlisting}
struct MCMF {
    struct E { int to; ll cap, cost; };
    int n; vector<E> e; vector<vector<int>> g;
    MCMF(int n = 0) : n(n), g(n) {}
    void addEdge(int u, int v, ll cap, ll cost) {
        g[u].push_back((int)e.size());
        e.push_back({v, cap, cost});
        g[v].push_back((int)e.size());
        e.push_back({u, 0, -cost});
    }
    pair<ll,ll> run(int s, int t, ll maxf = INF) {
        ll fl = 0, cost = 0;
        while (fl < maxf) {
            vi d(n, INF); vector<int> pe(n, -1);
            vector<char> inq(n, 0);
            deque<int> q{s}; d[s] = 0; inq[s] = 1;
            while (!q.empty()) {              // SPFA
                int u = q.front(); q.pop_front(); inq[u] = 0;
                for (int id : g[u]) {
                    if (!e[id].cap) continue;
                    int v = e[id].to;
                    if (d[u] + e[id].cost < d[v]) {
                        d[v] = d[u] + e[id].cost; pe[v] = id;
                        if (!inq[v]) {
                            inq[v] = 1; q.push_back(v);
                        }
                    }
                }
            }
            if (d[t] >= INF) break;
            ll push = maxf - fl;
            for (int v = t; v != s; v = e[pe[v] ^ 1].to)
                push = min(push, e[pe[v]].cap);
            for (int v = t; v != s; v = e[pe[v] ^ 1].to) {
                e[pe[v]].cap -= push;
                e[pe[v] ^ 1].cap += push;
            }
            fl += push; cost += push * d[t];
        }
        return {fl, cost};
    }
};
\end{lstlisting}

\entry{Bipartite matching (Kuhn)}{$O(nm)$ \quad mem $n + m$}{%
Left side \code{[0, n1)}, right side \code{[0, n2)}, \code{addEdge(u, v)} with
\code{u} left and \code{v} right. After \code{maxMatching()}, \code{matchL[u]} /
\code{matchR[v]} give the partner or \code{-1}.\\[0.3ex]
Short and fast enough for a few thousand vertices. Above that build a unit-capacity
\code{Dinic} instead --- same graph, $O(m\sqrt n)$.\\[0.3ex]
\textbf{K\"onig, minimum vertex cover.} Let $Z$ be everything reachable by alternating
paths from the \emph{unmatched} left vertices. Then a minimum vertex cover is
$(L \setminus Z) \cup (R \cap Z)$ and a maximum independent set is its complement; both
have the expected sizes.}

%LABEL Matching
\begin{lstlisting}
struct Matching {
    int n1, n2;
    vector<vector<int>> g;
    vector<int> matchL, matchR;
    vector<char> vis;
    Matching(int n1 = 0, int n2 = 0) : n1(n1), n2(n2), g(n1),
        matchL(n1, -1), matchR(n2, -1) {}
    void addEdge(int u, int v) { g[u].push_back(v); }
    bool tryKuhn(int u) {
        for (int v : g[u]) if (!vis[v]) {
            vis[v] = 1;
            if (matchR[v] == -1 || tryKuhn(matchR[v])) {
                matchL[u] = v; matchR[v] = u;
                return true;
            }
        }
        return false;
    }
    int maxMatching() {
        int r = 0;
        for (int u = 0; u < n1; u++) {
            vis.assign(n2, 0);
            if (tryKuhn(u)) r++;
        }
        return r;
    }
};
\end{lstlisting}

\entry{Minimum vertex cover (K\"onig)}{KACTL \quad $O(n + m)$ after matching}{%
The construction above as code, KACTL's \code{cover} on our \code{Matching}. Call
\code{maxMatching()} first. Returns the cover as \{left vertices, right vertices\};
its size equals the matching. Everything \emph{not} in it is a maximum independent set
--- ``choose the most items with no conflicting pair'' when conflicts only go between
two groups.}

%DEP Matching
%LABEL cover
\begin{lstlisting}
pair<vi, vi> vertexCover(Matching& M) {
    vector<char> lf(M.n1, 1), seen(M.n2, 0);
    for (int u : M.matchR) if (u != -1) lf[u] = 0;
    vi q, L, R;
    rep(i,0,M.n1) if (lf[i]) q.push_back(i);
    while (!q.empty()) {
        int i = q.back(); q.pop_back();
        lf[i] = 1;
        for (int e : M.g[i]) if (!seen[e] && M.matchR[e]+1) {
            seen[e] = 1;
            q.push_back(M.matchR[e]);
        }
    }
    rep(i,0,M.n1) if (!lf[i]) L.push_back(i);
    rep(i,0,M.n2) if (seen[i]) R.push_back(i);
    return {L, R};
}
\end{lstlisting}

\entry{Hungarian (assignment)}{$O(n^2 m)$ \quad mem $nm$}{%
Minimum-cost \textbf{perfect} matching on a cost matrix, rows $\le$ columns: every row
gets its own column, cheapest in total. Returns \code{\{cost, asg\}} with
\code{asg[row] = column}. Input is 0-indexed; the internals are 1-indexed, which is why
the matrix is read as \code{a[i0-1][j-1]}.\\[0.3ex]
\textbf{Maximise} instead: negate every cost, then negate the answer. Forbidden pairs:
a large finite cost, not \code{INF} (\code{INF} arithmetic overflows here) --- something
like $10^{12}$ with real costs below $10^6$.\\[0.3ex]
Not the same thing as \code{Matching}: that one maximises the \emph{number} of pairs,
this one fixes the number and minimises the \emph{cost}. Costs may be negative.}

%LABEL hungarian
\begin{lstlisting}
// a = cost matrix, a.size() <= a[0].size()
pair<ll, vector<int>> hungarian(const matrix& a) {
    int n = (int)a.size(), m = (int)a[0].size();
    // 1-indexed inside; column 0 is virtual, p[j]==0 => free
    vi u(n + 1, 0), v(m + 1, 0), minv(m + 1);
    vector<int> p(m + 1, 0), way(m + 1, 0);
    for (int i = 1; i <= n; i++) {
        p[0] = i; int j0 = 0;
        fill(all(minv), INF);
        vector<char> used(m + 1, 0);
        do {
            used[j0] = 1;
            int i0 = p[j0], j1 = 0; ll delta = INF;
            for (int j = 1; j <= m; j++) if (!used[j]) {
                ll cur = a[i0-1][j-1] - u[i0] - v[j];
                if (cur < minv[j]) {
                    minv[j] = cur; way[j] = j0;
                }
                if (minv[j] < delta) {
                    delta = minv[j]; j1 = j;
                }
            }
            for (int j = 0; j <= m; j++) {
                if (used[j]) {
                    u[p[j]] += delta; v[j] -= delta;
                } else minv[j] -= delta;
            }
            j0 = j1;
        } while (p[j0] != 0);
        do {                          // flip the path back
            int j1 = way[j0]; p[j0] = p[j1]; j0 = j1;
        } while (j0);
    }
    vector<int> asg(n);
    for (int j = 1; j <= m; j++) if (p[j]) asg[p[j]-1] = j-1;
    ll cost = 0;
    for (int i = 0; i < n; i++) cost += a[i][asg[i]];
    return {cost, asg};
}
\end{lstlisting}

\entry{Eulerian path}{$O(n + m)$ \quad mem $n + m$}{%
Hierholzer, iterative, handles both directions with the \code{dir} flag. Returns the
\textbf{vertex} sequence, length \code{m+1}, using every edge exactly once --- or an
empty vector when none exists (wrong degrees, or the edges are disconnected). A closed
tour comes back with \code{path.front() == path.back()}.\\[0.3ex]
Existence: \emph{directed} needs \code{out-in} equal to 0 everywhere except possibly one
$+1$ (the start) and one $-1$ (the end); \emph{undirected} needs zero or exactly two odd
vertices. Both also need all edges in one connected component --- which the final
\code{size() != m+1} test catches, so isolated vertices do not matter.\\[0.3ex]
For the edge sequence instead, record \code{id} alongside each pop.}

%LABEL eulerPath
\begin{lstlisting}
vector<int> eulerPath(int n, const vector<pii>& es,
                      bool dir) {
    int m = (int)es.size();
    if (!m) return {};
    vector<vector<pii>> g(n);          // (to, edge id)
    vi outd(n, 0), ind(n, 0);
    for (int i = 0; i < m; i++) {
        int a = (int)es[i].first, b = (int)es[i].second;
        g[a].push_back({b, i}); outd[a]++; ind[b]++;
        if (!dir) {
            g[b].push_back({a, i}); outd[b]++; ind[a]++;
        }
    }
    int s = -1;
    if (dir) {
        int plus = 0, minus = 0;
        for (int v = 0; v < n; v++) {
            ll df = outd[v] - ind[v];
            if (df == 1) { plus++; s = v; }
            else if (df == -1) minus++;
            else if (df) return {};
        }
        if (plus > 1 || minus > 1) return {};
    } else {
        int odd = 0;
        for (int v = 0; v < n; v++)
            if (outd[v] & 1) { odd++; s = v; }
        if (odd != 0 && odd != 2) return {};
    }
    if (s < 0)
        for (int v = 0; v < n; v++)
            if (outd[v]) { s = v; break; }
    vector<int> it(n, 0), st{s}, path;
    vector<char> used(m, 0);
    while (!st.empty()) {
        int u = st.back();
        while (it[u] < (int)g[u].size()
               && used[g[u][it[u]].second]) it[u]++;
        if (it[u] == (int)g[u].size()) {
            path.push_back(u); st.pop_back();
        } else {
            auto [v, id] = g[u][it[u]++];
            used[id] = 1; st.push_back((int)v);
        }
    }
    if ((int)path.size() != m + 1) return {};   // disjoint
    reverse(all(path));
    return path;
}
\end{lstlisting}

\entry{Global min cut (Stoer--Wagner)}{KACTL \quad $O(V^3)$}{%
The cheapest set of edges whose removal disconnects an \textbf{undirected} graph, with
\emph{no} \code{s} and \code{t} given. (If they are given, that is max flow: use
\code{Dinic}.) Running Dinic for every pair would be $V^2$ flows; this is one $O(V^3)$
pass, fine for $V \le 500$. With unit weights it is the edge connectivity.\\[0.3ex]
Input is an adjacency \textbf{matrix} of total weights --- add each edge to both
\code{mat[u][v]} and \code{mat[v][u]}, summing parallel edges. Returns
\code{\{weight, vertices on one side\}}. Each phase merges the last two vertices of a
maximum-adjacency order; the cut ``last vertex alone'' is a candidate.
KACTL's \code{int}s are \code{ll} here (\code{INT\_MIN} became \code{-INF}).}

%LABEL mincut
\begin{lstlisting}
pair<ll, vi> globalMinCut(vector<vi> mat) {
    pair<ll, vi> best = {LLONG_MAX, {}};
    int n = (int)mat.size();
    vector<vi> co(n);
    rep(i,0,n) co[i] = {i};
    rep(ph,1,n) {
        vi w = mat[0];
        size_t s = 0, t = 0;
        rep(it,0,n-ph) { // O(V^2); O(E log V) with a heap
            w[t] = -INF;
            s = t, t = max_element(all(w)) - w.begin();
            rep(i,0,n) w[i] += mat[t][i];
        }
        best = min(best, {w[t] - mat[t][t], co[t]});
        co[s].insert(co[s].end(), all(co[t]));
        rep(i,0,n) mat[s][i] += mat[t][i];
        rep(i,0,n) mat[i][s] = mat[s][i];
        mat[0][t] = -INF;
    }
    return best;
}
\end{lstlisting}

\section{Bitmask and DP}

\entry{PathDP (bitmask)}{$O(2^n n^2)$ \quad mem $2^n n$}{%
Folds every simple path covering \emph{exactly} a mask into one value, $n \le 20$.
Directed --- fill both directions for undirected. Pick a row of the table, edit the 5
config lines, done. The constructor does all the work; every call below is a
lookup.\\[0.3ex]
{\ttfamily\footnotesize
\begingroup\setlength{\tabcolsep}{2.5pt}
\begin{tabular}{@{}llllll@{}}
want & T & UNIT & NONE & ext(x,wt) & merge(a,b)\\
shortest path & ll & 0 & INF & x + wt & min(a, b)\\
longest path & ll & 0 & -INF & x + wt & max(a, b)\\
count paths & ll & 1 & 0 & x & a + b\\
count mod p & ll & 1 & 0 & x & (a + b) \% p\\
min bottleneck & ll & 0 & INF & max(x, wt) & min(a, b)\\
max probability & double & 1 & 0 & x * wt & max(a, b)\\
reachable lengths & bitset & 1<<0 & 0 & x << wt & a | b\\
\end{tabular}\endgroup}\\[0.3ex]
\begin{ctab}{0.25}{0.70}
\code{PathDP p(w,st,cn)} & \code{w[u][v]} = weight, \code{INF} = no edge.\\
\code{val(m,v)} & paths covering exactly \code{m}, ending at \code{v}.\\
\code{hamPath()} & over \code{v} of \code{val(full, v)}.\\
\code{hamPathTo(v)} & \code{val(full, v)}.\\
\code{hamCycle(s)} & closes \code{w[v][s]}; build with \code{st = 1LL<<s}.\\
\code{endingAt(v)} & over every mask, ending at \code{v}.\\
\code{bySize(k)} & over the masks of exactly \code{k} vertices.\\
\code{anyPath()} & over every mask and every \code{v}.\\
\code{cycles()} & over every simple cycle; build with \code{cn = true}.\\
\code{path(m,v)} & vertex order; \code{MIN} / \code{MAX} only.\\
\code{induced(g,vs)} & \code{w} for the vertices \code{vs}, weights all 1, local
\code{i} is \code{vs[i]}; translate back with \code{vs[x]}.\\
\end{ctab}\\[0.3ex]
\code{st} is a \textbf{mask}, not a vertex: \code{-1} = any start, \code{1LL<<s} pins
it. \code{cn} makes every path begin at the lowest vertex of its mask, for
\code{cycles()} only.\\[0.3ex]
\code{val(m,v)} is exactly \code{m}, not a subset, and \code{v} must be inside
\code{m}. Undirected counts come out $2\times$. \code{ext} must distribute over
\code{merge}, and \code{merge} must commute. Struct or bitset \code{T}:
\code{inline static const}, not \code{constexpr}. Unmodded counts overflow \code{ll} by
dense $n = 18$. At \code{-O2}, 30\% density: $n{=}16$ 0.01 s 13 MB, $n{=}18$ 0.05 s
44 MB, $n{=}20$ 0.22 s 184 MB.}

\begin{lstlisting}
struct PathDP {
    using T = ll;
    static T ext(T x, ll wt) { return x + wt; }
    static T merge(T a, T b) { return min(a, b); }
    static constexpr T UNIT = 0;
    static constexpr T NONE = 1e18;
    static constexpr ll INF = 1e18;      // missing edge in w
    int n; ll starts, full;
    matrix w; vector<T> dp; vector<char> par;
    PathDP(const matrix& wt, ll ss = -1, bool canon = false)
            : n(wt.size()), starts(ss), full((1LL << n) - 1),
              w(wt) {
        dp.assign((full + 1) * n, NONE);
        par.assign((full + 1) * n, -1);
        rep(v, 0, n) if (starts >> v & 1)
            dp[(1LL << v) * n + v] = UNIT;
        vector<ll> nb(n, 0);        // out-edges of u as bits
        rep(u, 0, n) rep(v, 0, n)
            if (w[u][v] < INF) nb[u] |= 1LL << v;
        rep(m, 1, full + 1) {
            ll base = m * n;
            for (ll b = m; b; b &= b - 1) {   // set bits of m
                int u = __builtin_ffsll(b) - 1;
                T cur = dp[base + u];
                if (cur == NONE) continue;
                for (ll c = nb[u] & ~m; c; c &= c - 1) {
                    int v = __builtin_ffsll(c) - 1;
                    if (canon && (1LL << v) < (m & -m))
                        continue;
                    ll t = (m | 1LL << v) * n + v;
                    T nv = merge(dp[t], ext(cur, w[u][v]));
                    if (nv != dp[t]) {
                        dp[t] = nv; par[t] = (char)u;
                    }
                }
            }
        }
    }
    T val(ll m, int v) { return dp[m * n + v]; }
    T hamPathTo(int v) { return val(full, v); }
    T hamPath() {
        T r = NONE;
        rep(v, 0, n) r = merge(r, val(full, v));
        return r;
    }
    T hamCycle(int s) {          // build with 1LL << s
        T r = NONE;
        rep(v, 0, n) {
            T d = val(full, v);
            if (d != NONE && w[v][s] < INF)
                r = merge(r, ext(d, w[v][s]));
        }
        return r;
    }
    T endingAt(int v) {
        T r = NONE;
        rep(m, 1, full + 1) r = merge(r, val(m, v));
        return r;
    }
    T bySize(int k) {            // exactly k vertices
        T r = NONE;
        rep(m, 1, full + 1) if (__builtin_popcountll(m) == k)
            rep(v, 0, n) r = merge(r, val(m, v));
        return r;
    }
    T anyPath() {
        T r = NONE;
        rep(m, 1, full + 1) rep(v, 0, n)
            r = merge(r, val(m, v));
        return r;
    }
    T cycles() {                 // build with canon
        T r = NONE;
        rep(m, 1, full + 1) {
            if (__builtin_popcountll(m) < 3) continue;
            int s = 0;
            while (!(m >> s & 1)) s++;   // lowest of m
            rep(v, 0, n) {
                T d = val(m, v);
                if (v != s && d != NONE && w[v][s] < INF)
                    r = merge(r, ext(d, w[v][s]));
            }
        }
        return r;
    }
    vi path(ll m, int v) {               // MIN / MAX only
        vi p;
        while (v >= 0) {
            p.push_back(v);
            int u = par[m * n + v]; m ^= 1LL << v; v = u;
        }
        reverse(p.begin(), p.end());
        return p;
    }
    template<class G, class V>    // g, vs: any integer type
    static matrix induced(const G& g, const V& vs) {
        static vi idx;        // reused, all -1 on entry
        if (idx.size() < g.size()) idx.assign(g.size(), -1);
        int k = (int)vs.size();
        rep(i, 0, k) idx[vs[i]] = i;
        matrix w(k, vi(k, INF));
        rep(i, 0, k) for (ll u : g[vs[i]])
            if (idx[u] >= 0) w[i][idx[u]] = 1;
        rep(i, 0, k) idx[vs[i]] = -1;   // wipe what we set
        return w;
    }
};
\end{lstlisting}

\entry{Subset sums (SOS / zeta)}{$O(2^n n)$ \quad mem $2^n$}{%
\code{sos(n, a)[m]} is the sum of \code{a[s]} over \emph{all} submasks \code{s} of
\code{m}. Replace \code{+=} with \code{|=} / \code{max} for ``or over submasks'' /
``best submask''. Supersets instead: flip the condition to
\code{if (!(m >> b \& 1)) f[m] += f[m | 1LL<<b]}.\\[0.3ex]
Inverse (M\"obius, recovers \code{a} from \code{f}): same loops with \code{-=}.
Subset convolution and ``count pairs with \code{a \& b == 0}'' both reduce to
this.\\[0.3ex]
To walk the submasks of one mask only: \code{for (ll s = m; s; s = (s-1) \& m)}, which
visits every nonempty submask once --- $O(3^n)$ over all masks.}

\begin{lstlisting}
// f has size 1<<n and holds a on entry
vi sos(int n, vi f) {
    rep(b, 0, n) rep(m, 0, 1LL << n)
        if (m >> b & 1) f[m] += f[m ^ (1LL << b)];
    return f;
}
\end{lstlisting}

\entry{LIS}{$O(n \log n)$ \quad mem $n$}{%
\code{c[k]} ends up as the smallest possible tail of an increasing subsequence of
length \code{k+1}, so \code{c.size()} is the answer. \code{c} itself is \emph{not} a
valid subsequence.\\[0.3ex]
\code{lower\_bound} gives \textbf{strictly} increasing; \code{upper\_bound} gives
non-decreasing. Longest decreasing: reverse the input, or negate it. Minimum number of
non-increasing chains covering the sequence = length of the LIS (Dilworth).\\[0.3ex]
To reconstruct, store for each element the index it was placed at plus the predecessor
\code{c}-index, then walk back from the last placement.}

\begin{lstlisting}
ll lis(const vi& a) {
    vi c;
    for (ll e : a) {
        auto it = lower_bound(all(c), e);
        if (it == c.end()) c.push_back(e);
        else *it = e;
    }
    return (ll)c.size();
}
\end{lstlisting}

\entry{Interval DP}{$O(n^3)$, optimised $O(n^2)$ \quad mem $n^2$}{%
Shortest interval first, so both halves are ready when you need them. Circular
versions: double the array to length $2n$ and take the best window of length
$n$.\\[0.3ex]
\textbf{Knuth} cuts it to $O(n^2)$ when \code{opt[i][j-1] <= opt[i][j] <= opt[i+1][j]}
--- restrict \code{k} to that window and keep an \code{opt} table. It holds when the
cost satisfies the quadrangle inequality, which is true for the classic ``merge
adjacent piles, cost = sum of the pile'' problem. \textbf{Divide and conquer
optimisation} is the same idea one layer at a time, for
\code{dp[i][j] = min over k of dp[i-1][k] + C(k,j)} with monotone \code{opt}:
recurse on \code{(jlo,jhi,klo,khi)}, evaluate the midpoint, split. Both are worth it
only when $n^3$ is too slow --- check the bound first.}

%NOCOMPILE
\begin{lstlisting}
rep(len, 2, n + 1)
  rep(i, 0, n - len + 1) {
    ll j = i + len - 1;                  // inclusive
    dp[i][j] = INF;
    rep(k, i, j)                         // split point
      dp[i][j] = min(dp[i][j],
                     dp[i][k] + dp[k+1][j] + cost(i, j));
  }
\end{lstlisting}

\entry{Divide and conquer DP}{KACTL \quad $O(n \log n)$ per layer}{%
The optimisation described above, as code. Shape: split \code{1..n} into \code{G}
consecutive groups, \code{cur[i] = min over k < i of prv[k] + cost(k, i)} where
\code{cost(k, i)} is the price of the group \code{(k, i]}, one layer per group. Naively
$O(Gn^2)$; this is $O(Gn\log n)$.\\[0.3ex]
Valid when the best \code{k} never moves left as \code{i} grows --- true when
\code{cost} satisfies the quadrangle inequality, e.g.\ (sum of the group)$^2$, or the
number of equal pairs inside it. Not sure? Compare with the naive DP on random small
inputs.\\[0.3ex]
Per layer: \code{prv} = previous layer (\code{prv[0] = 0}, the rest \code{INF} for
layer 1), \code{cur.assign(n+1, INF); DP().solve(1, n+1); prv = cur;}. \code{lo},
\code{hi} bound \code{k} to \code{[lo, hi)}. \code{cost} must be $O(1)$: prefix
sums.}

%LABEL dncdp
\begin{lstlisting}
vi prv, cur;          // previous layer, layer being built
ll cost(int k, int i);              // cost of group (k, i]
struct DP { // Modify at will:
    int lo(int ind) { return 0; }
    int hi(int ind) { return ind; }
    ll f(int ind, int k) { return prv[k] + cost(k, ind); }
    void store(int ind, int k, ll v) { cur[ind] = v; }
    void rec(int L, int R, int LO, int HI) {
        if (L >= R) return;
        int mid = (L + R) >> 1;
        pair<ll, ll> best(LLONG_MAX, LO);
        rep(k, max(LO,lo(mid)), min(HI,hi(mid)))
            best = min(best, make_pair(f(mid, k), k));
        store(mid, best.second, best.first);
        rec(L, mid, LO, best.second+1);
        rec(mid+1, R, best.second, HI);
    }
    void solve(int L, int R) { rec(L, R, INT_MIN, INT_MAX); }
};
\end{lstlisting}

\section{Number theory}

\entry{Modular basics}{$O(\log)$ each}{%
\code{power(a, b, m)} is $a^b \bmod m$; it needs $m < 2^{31}$ or the \code{a*a} inside
overflows --- for bigger moduli cast to \code{i128}.\\[0.3ex]
\code{ext\_gcd(a, b, x, y)} returns $g=\gcd$ and fills $ax+by=g$. \code{inv(a, m)}
gives $a^{-1} \bmod m$ for \emph{any} \code{m} with $\gcd(a,m)=1$, and \code{-1} when
none exists. When \code{m} is prime, \code{power(a, m-2, m)} is the same thing and
shorter.\\[0.3ex]
$a^b \bmod m$ with a \emph{gigantic} $b$: read $b$ as a string, reduce it mod
$\varphi(m)$ --- but only add $\varphi(m)$ back (i.e. use $b \bmod \varphi(m) +
\varphi(m)$) when the true $b \ge \varphi(m)$, otherwise you get the wrong answer for
small $b$. That correction is needed exactly when $\gcd(a,m) \ne 1$.}

%LABEL modbasics
\begin{lstlisting}
ll power(ll a, ll b, ll mod) {         // a^b % mod
    a %= mod; ll res = 1;
    for (; b > 0; b >>= 1, a = a * a % mod)
        if (b & 1) res = res * a % mod;
    return res;
}
ll ext_gcd(ll a, ll b, ll &x, ll &y) {     // a*x + b*y = g
    if (b == 0) { x = 1; y = 0; return a; }
    ll x1, y1;
    ll g = ext_gcd(b, a % b, x1, y1);
    x = y1;
    y = x1 - (a / b) * y1;
    return g;
}
ll inv(ll a, ll m) {                 // -1 if none exists
    ll x, y;
    if (ext_gcd(a, m, x, y) != 1) return -1;
    return ((x % m) + m) % m;
}
\end{lstlisting}

\entry{Fast power}{Ruiming \quad $O(\log p)$}{%
\code{fast(x, p, m)} is $x^p \bmod m$. Needs $m < 2^{31}$ or \code{x*x} overflows.
\emph{Euler's theorem}, \emph{Binomials (inverse)} below call it.}

%LABEL fastmod
%FROM Ruiming/math/Fastmod.cpp
\begin{lstlisting}
ll fast(ll x,ll p,ll m){
    x%=m;
    ll ans=1;
    while(p>0){
        if(p&1){
            ans*=x;
            ans%=m;
        }
        x*=x;
        x%=m;
        p>>=1;
    }
    return ans;
}
\end{lstlisting}

\entry{Fast gcd}{Ruiming}{%
Binary (Stein) gcd with \code{\_\_builtin\_ctz}. Both arguments must be
\textbf{non-zero}: \code{\_\_builtin\_ctz(0)} is undefined.}

%FROM Ruiming/math/Fast_gcd.cpp
\begin{lstlisting}
ll gcd(ll a,ll b)
{
    ll az=__builtin_ctz(a),bz=__builtin_ctz(b);
    ll z=min(az,bz);
    ll dif;
    b>>=bz;
    while(a)
    {
        a>>=az;
        dif=b-a;
        az=__builtin_ctz(dif);
        if(a<b) b=a;
        if(dif<0) a=-dif;
        else a=dif;
    }
    return b<<z;
}
\end{lstlisting}

\entry{CRT}{Ruiming \quad $O(\sum m_i)$}{%
A whole program, \textbf{1-indexed}: reads \code{n}, then \code{n} pairs
\code{m[i] a[i]} meaning $x \equiv a_i \pmod{m_i}$, prints the smallest such $x$.
Moduli must be \textbf{pairwise coprime}. Merges one congruence at a time by trying
\code{j = 0..m[i]-1}, so it is linear in the moduli --- use \emph{CRT (general)} when
they are large or not coprime.}

%FROM Ruiming/math/CRT.cpp
\begin{lstlisting}
ll n,lst[MAXN],a[MAXN],m[MAXN],ans[MAXN],k,temp;
int main(){
    ios_base::sync_with_stdio(false);
    cin.tie(NULL);
    cin>>n;
    for(int i=1;i<=n;i++){
        cin>>m[i]>>a[i];
    }
    ans[1]=a[1];
    for(int i=2;i<=n;i++){
        for(ll j=0;j<=m[i]-1;j++){
            if(((j*m[i-1]+ans[i-1])%m[i])==a[i]){
                ans[i]=j*m[i-1]+ans[i-1];
                break;
            }
        }
        m[i]*=m[i-1];
    }
    cout<<ans[n]<<'\n';
}
\end{lstlisting}

\entry{CRT (general)}{Martin \quad $O(\log)$}{%
Merges $x \equiv r_1 \pmod{m_1}$ and $x \equiv r_2 \pmod{m_2}$ into one congruence
modulo $\mathrm{lcm}(m_1,m_2)$. Moduli need \emph{not} be coprime; returns
\code{\{-1,-1\}} when the two are incompatible. Fold a list by merging one at a
time.\\[0.3ex]
Watch the size: the intermediate \code{m1 / g * m2} must fit, so with several moduli
near $10^9$ switch to \code{i128}.}

%LABEL crt
%DEP modbasics
\begin{lstlisting}
// x = r1 (mod m1), x = r2 (mod m2)  ->  {r, lcm} or {-1,-1}
pair<ll,ll> crt(ll r1, ll m1, ll r2, ll m2) {
    ll x, y, g = ext_gcd(m1, m2, x, y);
    if ((r2 - r1) % g) return {-1, -1};
    ll l = m1 / g * m2, t = m2 / g;
    ll k = ((r2 - r1) / g % t) * (x % t) % t;
    ll r = (r1 + (i128)m1 * k) % l;
    return {(r % l + l) % l, l};
}
\end{lstlisting}

\entry{Sieves}{$O(n \log\log n)$ \quad mem $n$}{%
\code{sieve(n)} lists the primes up to \code{n} inclusive. \code{spf(n)} gives the
smallest prime factor of every number, which factorises any \code{x <= n} in
$O(\log x)$ by dividing repeatedly --- that is the one to build when you must factorise
many numbers.}

%LABEL sieves
\begin{lstlisting}
vector<int> sieve(int n) {
    vector<char> comp(n + 1, 0); vector<int> ps;
    for (int i = 2; i <= n; i++) if (!comp[i]) {
        ps.push_back(i);
        for (ll j = (ll)i * i; j <= n; j += i) comp[j] = 1;
    }
    return ps;
}
vector<int> spf(int n) {         // smallest prime factor
    vector<int> s(n + 1);
    iota(all(s), 0);
    for (int i = 2; (ll)i * i <= n; i++) if (s[i] == i)
        for (int j = i * i; j <= n; j += i)
            if (s[j] == j) s[j] = i;
    return s;
}
\end{lstlisting}

\entry{Euler sieve}{Ruiming \quad $O(n)$}{%
A whole program: reads \code{n q}, sieves the primes up to \code{n} into
\code{primes} (linear sieve), then answers \code{q} queries ``the $a$-th prime''.
\code{MAXN} must exceed \code{n}.}

%FROM Ruiming/math/Euler_sieve.cpp
\begin{lstlisting}
vector<ll> primes;
bool p[MAXN];
ll n,q;
int main(){
    ios_base::sync_with_stdio(false);
    cin.tie(NULL);
    cin>>n>>q;
    for(ll i=2;i<=n;i++){
        if(!p[i]){
            primes.push_back(i);
        }
        for(ll j=0;j<primes.size();j++){
            if(primes[j]*i>n) break;
            p[primes[j]*i]=1;
            if(i%primes[j]==0) break;
        }
    }
    for(int i=1;i<=q;i++){
        ll a;
        cin>>a;
        cout<<primes[a-1]<<'\n';
    }
}
\end{lstlisting}

\entry{Linear sieve + Euler phi}{Ruiming \quad $O(n)$ \quad mem $n$}{%
A whole program: reads \code{n}, fills \code{phi[1..n]} and \code{primes} (the
Euler / linear sieve, each composite crossed out exactly once). Drop the
\code{phi} lines for a plain linear sieve. \code{MAXN} must exceed \code{n}.\\[0.3ex]
$\varphi(n) = n \prod_{p \mid n} (1 - 1/p)$; $\sum_{d \mid n} \varphi(d) = n$. Number of
divisors and their sum are multiplicative, so the same sieve shape computes them.\\[0.3ex]
$a^b \bmod m$ with a gigantic $b$: see \emph{Euler's theorem} below.}

%FROM Ruiming/math/Linearphi.cpp
\begin{lstlisting}
ll n,phi[MAXN],ans=1;
bool p[MAXN];
vector<ll> primes;
int main(){
    ios_base::sync_with_stdio(false);
    cin.tie(NULL);
    cin>>n;
    if(n==1) cout<<0<<'\n';
    else{
        phi[1]=1;
        for(int i=2;i<=n;i++){//modiying euler sieve
            if(!p[i]){
                primes.push_back(i);
                phi[i]=i-1;
            }
            for(int j=0;j<primes.size();j++){
                if(primes[j]*i>n) break;
                p[primes[j]*i]=1;
                if(i%primes[j]==0){
                    phi[i*primes[j]]=primes[j]*phi[i];
                    break;
                }else{
                    phi[i*primes[j]]=(primes[j]-1)*phi[i];
                }
            }
        }
    }
}
\end{lstlisting}

\entry{Euler phi of one number}{Ruiming \quad $O(\sqrt n)$}{%
\code{phi(n)} by trial division. For every value up to $n$ at once, use the linear
sieve above.}

%FROM Ruiming/math/Phi.cpp
\begin{lstlisting}
ll phi(ll n) {//Euler's phi function
    ll ans = 1;
    for(ll i = 2; i * i <= n; ++ i) {
        if(n % i == 0) {
            n /= i;
            ans *= i - 1;
            while(n % i == 0) {
                n /= i;
                ans *= i;
            }
        }
    }
    if(n > 1)//optimisation: n>1 means one prime is left
        ans *= n - 1;
    return ans;
}
\end{lstlisting}

\entry{Prime factorization}{Ruiming \quad $O(\sqrt n)$}{%
\code{factor(n)} returns the prime factors with multiplicity, in increasing order, by
trial division. Fine up to $\sim10^{12}$; beyond that use Pollard rho.}

%FROM Ruiming/math/Prime_factorization.cpp
\begin{lstlisting}
vector<ll> factor(ll n) {
    vector<ll> ret;
    for (ll i = 2; i * i <= n; i++) {
        while (n % i == 0) {
            ret.push_back(i);
            n /= i;
        }
    }
    if (n > 1) { ret.push_back(n); }
    return ret;
}
\end{lstlisting}

\entry{Factorization table}{Ruiming \quad $O(\texttt{MAXN}\sqrt{\texttt{MAXN}})$}{%
A whole program: \code{fac[i]} is the factorization of every \code{i < MAXN} as a
\code{map} prime $\rightarrow$ exponent, then reads \code{a[1..n]}. His file uses
\code{MAXN = 36}.}

%FROM Ruiming/math/Factorization.cpp
\begin{lstlisting}
ll n,a[MAXN],ans;
map<ll,ll> fac[MAXN];
int main(){
    ios_base::sync_with_stdio(false);
    cin.tie(NULL);
    cin>>n;
    for(int i=1;i<MAXN;i++){
        ll temp=i;
        for(int j=2;j<=sqrt(temp+0.5);j++){
            while(temp%j==0){
                fac[i][j]++;
                temp/=j;
            }
        }
        if(temp>1){
            fac[i][temp]++;
        }
    }
    for(int i=1;i<=n;i++){
        cin>>a[i];
    }
}
\end{lstlisting}

\entry{Euler's theorem (huge exponent)}{Ruiming \quad $O(m + |b|)$}{%
A whole program: reads \code{a m}, then the exponent \code{b} as a string, prints
$a^b \bmod m$. Computes $\varphi(m)$, reduces \code{b} mod $\varphi(m)$ while reading,
and adds $\varphi(m)$ back only if the true $b \ge \varphi(m)$ (\code{flag}) --- which
is what makes it correct when $\gcd(a,m) \ne 1$. Uses \code{fast} from
\emph{Fast power}.}

%DEP fastmod
%FROM Ruiming/math/Eulerthm.cpp
\begin{lstlisting}
bool flag;
ll gcd(ll x,ll y){
    while(x!=y){
        if(x>y) swap(x,y);
        if(y%x==0) return x;
        y-=x*(y/x);
    }
    return x;
}
ll a,m,phi=1;
string b;
ll read(){
    string x;
    cin>>x;
    ll ans=0;
    for(int i=0;i<x.size();i++){
        ans=(ans*10+(x[i]-'0'));
        if(ans>=phi) flag=true;
        ans%=phi;
    }
    return ans;
}
int main(){
    cin>>a>>m;
    //find the euler function to m
    ll d=gcd(a,m);
    ll temp=m;
    for(int i=2;i<=temp;i++){
        //here uses temp still makes sure i is a prime
        if(temp%i==0){
            ll times=0;
            while(temp%i==0){
                times++;
                temp/=i;//makes sure i is a prime
            }
            phi*=pow(i,times-1)*(i-1);
        }
    }
    //if(phi==1) phi=m-1;
    //b needs to be stored as a string and takes mod phi
    //then we can calculate it with fast mod
    ll b=read();
    ll ans;
    if(flag) ans=fast(a,b+phi,m);
    else ans=fast(a,b,m);
    //deal with d d|m, will lead up to a ,
    //should be some period
    cout<<ans<<'\n';
}
\end{lstlisting}

\entry{Miller-Rabin (standalone)}{Ruiming \quad $O(12 \log n)$}{%
A complete file on its own --- \textbf{do not paste it under the header}: its
\code{typedef \_\_int128 ll} replaces \code{ll}, so \code{n} can exceed
\code{long long}, with \code{read}/\code{write}/\code{print} fast I/O for
\code{\_\_int128}. Reads \code{t} numbers, prints \code{Yes}/\code{No}. \code{a * a}
in \code{qpow} overflows once \code{n} passes $\sim1.3\cdot10^{19}$.}

%STANDALONE
%FROM Ruiming/math/miller_rabin.cpp
\begin{lstlisting}
#include <bits/stdc++.h>
using namespace std;
typedef __int128 ll;
typedef pair<int, int> pii;
template<typename T> inline T read() {
    T x = 0, f = 1; char ch = 0;
    for(; !isdigit(ch); ch = getchar()) if(ch == '-') f = -1;
    for(; isdigit(ch); ch = getchar())
        x = (x << 3) + (x << 1) + (ch - '0');
    return x * f;
}
template<typename T> inline void write(T x) {
    if(x < 0) putchar('-'), x = -x;
    if(x > 9) write(x / 10);
    putchar(x % 10 + '0');
}
template<typename T> inline void print(T x, char ed = '\n') {
    write(x), putchar(ed);
}
ll t, n;
ll qpow(ll a, ll b, ll mod) {
    ll ret = 1;
    while(b) {
        if(b & 1) ret = (ret * a) % mod;
        a = (a * a) % mod;
        b >>= 1;
    }
    return ret % mod;
}
vector<ll> p = {2, 3, 5, 7, 11, 13, 17, 19, 23, 29, 31, 37};
bool miller_rabin(ll n) {
    if(n < 3 || n % 2 == 0) return n == 2;
    ll u = n - 1, t = 0;
    while(u % 2 == 0) u /= 2, ++ t;
    for(auto a : p) {
        if(n == a) return 1;
        if(n % a == 0) return 0;
        ll v = qpow(a, u, n);
        if(v == 1) continue;
        ll s = 1;
        for(; s <= t; ++ s) {
            if(v == n - 1) break;
            v = v * v % n;
        }
        if(s > t) return 0;
    }
    return 1;
}
int main() {
    t = read<ll>();
    while(t --) {
        n = read<ll>();
        if(miller_rabin(n)) puts("Yes");
        else puts("No");
    }
    return 0;
}
\end{lstlisting}

\entry{Miller-Rabin + Pollard rho}{Ruiming \quad $O(n^{1/4})$}{%
\code{miller\_rabin(n)} is \textbf{deterministic} for every \code{n} that fits in
\code{long long} (those 12 bases suffice below $3.3\cdot10^{24}$).\\[0.3ex]
\code{p.clear(); pollard\_rho(n);} leaves the prime factors of \code{n} in the global
\code{p}, with multiplicity, unsorted --- sort to group them. Multiplications go through
\code{i128}, so \code{n} up to $9\cdot10^{18}$ is safe.\\[0.3ex]
For $n \le 10^{12}$ or so, plain trial division up to $\sqrt n$ is shorter.}

%LABEL pollard
%FROM Ruiming/math/Pollard_Rho.cpp
\begin{lstlisting}
#define LL long long
#define i128 __int128
mt19937_64 rd(time(0));
i128 ksm(i128 x, i128 y, LL mod)
{
    i128 res = 1;
    x %= mod;
    while(y)
    {
        if(y & 1) res = res * x % mod;
        y >>= 1;
        x = x * x % mod;
    }
    return res;
}
bool miller_rabin(LL n)
{
    if(n <= 1 || n > 2 && ~n & 1) return false;
    LL m = n - 1, k = 0;
    while(~m & 1) m >>= 1, k++;
    vector<int> s = {2, 3, 5, 7, 11, 13, 17, 19, 23, 29,
                     31, 37};
    for(int a : s)
    {
        if(a >= n) break;
        i128 x = ksm(a, m, n), y;
        for(int i = 1; i <= k; i++, x = y)
        {
            y = x * x % n;
            if(y == 1 && x != 1 && x != n - 1) return false;
        }
        if(x != 1) return false;
    }
    return true;
}
vector<LL> p;
void pollard_rho(LL n)
{
    if(n <= 1) return;
    if(miller_rabin(n))
    {
        p.push_back(n);
        return;
    }
    LL d = rd() % (n - 1) + 1;
    auto f = [&](LL x)
    {
        return ((i128)x * x + d) % n;
    };
    LL x = rd() % (n - 1) + 1, y = f(x), t = 1;
    while(t == 1)
    {
        for(int i = 1; i <= 32; i++)
        {
            while(x == y)
                d = rd() % (n - 1) + 1,
                x = rd() % (n - 1) + 1, y = f(x);
            if((i128)t * abs(x - y) % n == 0) break;
            t = (i128)t * abs(x - y) % n;
            x = f(x), y = f(f(y));
        }
        t = __gcd(n, t);
    }
    pollard_rho(t), pollard_rho(n / t);
}
\end{lstlisting}

\entry{Binomials mod p}{precompute $O(n)$, query $O(1)$ \quad mem $n$}{%
Call \code{precompute()} once, then \code{C(n, k)}. Uses the header's \code{MAXN}, so
\textbf{raise \code{MAXN}} above the largest \code{n}. \code{MOD} must be
\textbf{prime} and larger than \code{MAXN}. Returns 0 outside the triangle, so no
guards at the call site.\\[0.3ex]
Non-prime modulus, or $n$ small: build Pascal's triangle instead
(\code{c[i][j] = c[i-1][j-1] + c[i-1][j]}), which needs no inverses. Huge $n$ with small
prime $p$: Lucas' theorem, multiply \code{C} of the base-$p$ digits.\\[0.3ex]
Stars and bars: ways to write $n$ as an ordered sum of $k$ non-negative parts is
\code{C(n+k-1, k-1)}. Catalan $C_n = $ \code{C(2n,n) * inv(n+1)}.}

%DEP modbasics
%LABEL binom
\begin{lstlisting}
const ll MOD = 1000000007;      // MAXN from the header
ll fact[MAXN], inv_fact[MAXN];
void precompute() {
    fact[0] = 1;
    rep(i, 1, MAXN) fact[i] = fact[i-1] * i % MOD;
    inv_fact[MAXN-1] = power(fact[MAXN-1], MOD - 2, MOD);
    down(i, MAXN - 1, 0)
        inv_fact[i] = inv_fact[i+1] * (i+1) % MOD;
}
ll C(ll n, ll k) {
    if (n < 0 || k < 0 || k > n || n >= MAXN) return 0;
    return fact[n] * inv_fact[k] % MOD
                   * inv_fact[n-k] % MOD;
}
\end{lstlisting}

\entry{Binomials (Pascal)}{Ruiming \quad $O(\texttt{MAXN} \cdot 106)$}{%
A whole program: Pascal's triangle \code{c[i][j]} mod \code{M} for \code{i < MAXN-1},
\code{j < 106}. Needs no inverses, so it works for any modulus.}

%FROM Ruiming/math/Combinumber.cpp
\begin{lstlisting}
const ll M=1000000007;
ll c[MAXN][106],t,n,m,k,lst[MAXN],ans;
int main(){
    ios_base::sync_with_stdio(false);
    cin.tie(NULL);
    c[0][0]=1;
    for(int i=1;i<MAXN-1;i++){
        c[i][0]=1;
        if(i<106) c[i][i]=1;
        for(int j=1;j<min(i,105);j++){
            c[i][j]=c[i-1][j]+c[i-1][j-1];
            c[i][j]%=M;
        }
    }
}
\end{lstlisting}

\entry{Binomials (inverse)}{Ruiming \quad $O(n \log \texttt{MOD})$}{%
A whole program: reads \code{n k}, fills \code{fac[i]} $= i!$ and \code{inv[i]}
$= (i!)^{-1}$ mod \code{MOD}, then $\binom{n}{i}$ is \code{comb} in the last loop. Uses
\code{fast} from \emph{Fast power}. His file uses \code{MAXN = 26}.}

%DEP fastmod
%FROM Ruiming/math/Combinumber_inv.cpp
\begin{lstlisting}
const ll MOD=1e9+7;
ll t,n,fac[MAXN],inv[MAXN],ans,k;
int main(){
    ios_base::sync_with_stdio(false);
    cin.tie(NULL);
    cin>>n>>k;
    ll temp=1;
    for(int i=1;i<=n;i++){
        temp*=i;
        temp%=MOD;
        inv[i]=fast(temp,MOD-2,MOD);
        fac[i]=temp;
    }
    fac[0]=1;
    inv[0]=1;
    for(int i=0;i<=n;i++){
        ll comb=((fac[n]*inv[i])%MOD*inv[n-i])%MOD;
    }
}
\end{lstlisting}

\entry{Linear inverses}{Ruiming \quad $O(n)$}{%
A whole program: reads \code{n p} ($p$ prime), prints $i^{-1} \bmod p$ for
$i = 1..n$, using one modular exponentiation and the identity
$i^{-1} = (i-1)! \cdot (i!)^{-1}$.}

%FROM Ruiming/math/Linearinverse.cpp
\begin{lstlisting}
ll n,p,ans[MAXN],corr[MAXN],now,fac[MAXN],inv[MAXN];
ll fastMod(ll a,ll n,ll mod){
    ll ans=1;
    a%=mod;
    while(n){
        if(n&1){
            ans=(ans*a)%mod;
        }
        a=(a*a)%mod;
        n>>=1;
    }
    return ans;
}
int main(){
    cin>>n>>p;
    fac[0]=1;
    for(int i=1;i<=n;i++){
        fac[i]=(fac[i-1]*i)%p;
    }
    inv[n]=fastMod(fac[n],p-2,p);//inv fac[i]
    for(int i=n-1;i>=1;i--){
        inv[i]=((i+1)*inv[i+1])%p;
    }
    for(int i=1;i<=n;i++){
        cout<<(fac[i-1]*inv[i])%p<<'\n';
    }
}
\end{lstlisting}

\entry{Matrix power}{Ruiming \quad $O(k^3 \log n)$}{%
\textbf{1-indexed}, fixed \code{m[MAXN][MAXN]}: set \code{row}/\code{col} and fill
\code{m[1..row][1..col]}. \code{a * b} needs \code{a.col == b.row};
\code{pow\_matrix(a, n)} needs \code{a} square.\\[0.3ex]
\textbf{Delete the header's \code{typedef vector<vi> matrix;}} --- same name. Keep
\code{MAXN} small (his file uses 106): every \code{matrix} is $8\,\texttt{MAXN}^2$ bytes
and \code{operator*} builds one on the stack. Clashes with \code{MOD} in
\emph{Binomials} if both are pasted.\\[0.3ex]
Linear recurrences: Fibonacci is \code{\{\{1,1\},\{1,0\}\}} to the $n$, entry
\code{m[1][2]}. With \code{min}/\code{+} instead of \code{+}/\code{*} it gives
``cheapest walk using exactly $n$ edges''.}

%UNDEF matrix
%FROM Ruiming/math/Matrix.cpp
\begin{lstlisting}
const ll MOD=1e9+7;
struct matrix{
    ll row,col;
    ll m[MAXN][MAXN];
};
matrix operator * (const matrix &a,const matrix &b){
    //assuming they fits a.col==b.row
    matrix c;
    memset(c.m,0,sizeof(c.m));
    c.row=a.row;
    c.col=b.col;
    for(int i=1;i<=c.row;i++){
        for(int j=1;j<=c.col;j++){
            for(int k=1;k<=a.col;k++){
                c.m[i][j]+=(a.m[i][k]*b.m[k][j])%MOD;
                c.m[i][j]%=MOD;
            }
        }
    }
    return c;
}
matrix pow_matrix(matrix a,ll n){
    //assuming a is a square matrix
    matrix res;
    memset(res.m,0,sizeof(res.m));
    res.row=res.col=a.col;
    for(int i=1;i<MAXN;i++) res.m[i][i]=1;
    while(n){
        if(n&1) res=res*a;
        a=a*a;
        n/=2;
    }
    return res;
}
\end{lstlisting}

\entry{Gaussian elimination}{Ruiming \quad $O(n^3)$}{%
A whole program, \textbf{1-indexed}, square $n\times n$: reads \code{n}, then each row
of \code{A} followed by its \code{b[i]}. Prints \code{-1} if there is no solution,
\code{0} if there are infinitely many, else \code{x1=..} to 2 decimals. Forward
elimination, then back substitution.\\[0.3ex]
Zero tests are exact (\code{==0}) and there is no largest-pivot choice: fine for integer
input, but with real input compare \code{fabs(..) < eps} instead.\\[0.3ex]
Over $\mathbb{F}_2$ use \code{bitset<N>} rows and \code{\textasciicircum=}, which is
64x faster and exact --- that is the version for XOR-basis and lights-out problems.
Mod a prime: replace \code{/} with \code{inv}.}

%FROM Ruiming/math/gaussjordanelimination.cpp
\begin{lstlisting}
ll n;
double A[MAXN][MAXN],b[MAXN],ans[MAXN];
stack<pair<ll,ll>> swaped;
int main(){
    ios_base::sync_with_stdio(false);
    cin.tie(NULL);
    cin>>n;
    for(int i=1;i<=n;i++){
        for(int j=1;j<=n;j++){
            cin>>A[i][j];
        }
        cin>>b[i];
    }
    ll col=1;
    for(int i=1;i<=n&&col<=n;i++,col++){//row
        while(A[i][col]==0&&col<=n){
            bool found=false;
            for(int j=i+1;j<=n;j++){
                if(A[j][col]!=0){
                    swap(A[j],A[i]);
                    swap(b[j],b[i]);//swap b as well
                    found=true;
                    break;
                }
            }
            if(!found) col++;
        }
        for(int j=i+1;j<=n;j++){
            if(A[i][col]==0) break;//excepet bug
            double fac=A[j][col]/A[i][col];
            for(int k=col;k<=n;k++){
                A[j][k]-=fac*A[i][k];
            }
            b[j]-=fac*b[i];
        }
    }
    bool no=false,p=false;
    for(int i=1;i<=n;i++){
        bool zero=true;
        for(int j=1;j<=n;j++){
            if(A[i][j]!=0){
                zero=false;
                break;
            }
        }
        p|=zero;
        if(zero&&b[i]!=0){
            no=true;
            break;
        }
    }
    if(p){
        if(no){
            cout<<-1<<'\n';
        }else cout<<0<<'\n';
    }else{
        for(int i=n;i>=1;i--){
            for(int j=i+1;j<=n;j++){
                b[i]-=ans[j]*A[i][j];
            }
            ans[i]=b[i]/=A[i][i];
        }
        for(int i=1;i<=n;i++){
            printf("x%d=%.2f\n",i,ans[i]);
        }
    }
}
\end{lstlisting}

\entry{0/1 fractional programming}{Ruiming \quad $O(n \log n \log \frac{n}{\varepsilon})$}{%
\textbf{1-indexed} globals: \code{n} items with values \code{a[i]} and weights
\code{b[i]}; drop \code{k} of them to \textbf{maximise} $\sum a / \sum b$. Binary
search on the answer $x$: it is achievable iff the best $n-k$ items have
$\sum (a_i - x b_i) \ge 0$, which \code{check} tests by sorting. The search range is
$[0, n]$ --- widen \code{end} if the ratio can exceed $n$.\\[0.3ex]
Returns $\lfloor 100x + 0.5 \rfloor$ (his problem printed a percentage); use
\code{start} directly for the real value. His file sets \code{eps = 1e-6} and
\code{MAXN = 100006}. Same shape for max average / min ratio cycle: change
\code{check}.}

%FROM Ruiming/math/zero_one_partition.cpp
%LABEL frac
\begin{lstlisting}
ll n,a[MAXN],b[MAXN],ans,k;
//eps=1e-6;
bool check(ld x){
    vector<ld> all;
    ld res=0;
    for(int i=1;i<=n;i++){
        all.push_back(a[i]-x*b[i]);
    }
    sort(all.begin(),all.end(),greater<ld>());
    for(int i=0;i<n-k;i++){
        res+=all[i];
    }
    return res>=-eps;
}
ll zero_one_partition(){
    ld start=0,end=n;
    while(end-start>eps){
        ld x=start+(end-start)/2;
        //if f(x)>0
        if(check(x)){
            start=x+eps;
        }else end=x-eps;
    }
    ll res = floor(100 * start + 0.5);
    return res;
}
\end{lstlisting}

\entry{Big integers}{Ruiming}{%
Decimal strings, optional leading \code{-}. \code{Bigadd}, \code{Bigsub},
\code{Bigmul} ($O(nm)$ schoolbook), \code{Bighalf} (non-negative only).
\code{Scompare}/\code{Ecompare} are $<$ / $\le$ for non-negative numbers.
\code{BigMod(x, m)} is $x \bmod m$ by binary search on the quotient (the same search
gives division); \code{BigModSmall(a, mod)} is a big number mod an \code{ll}.\\[0.3ex]
\code{Bigadd} of a positive and a negative number is \emph{not} handled (those two lines
are commented out) --- call \code{Bigsub} yourself.}

%FROM Ruiming/math/Big_Integer.cpp
\begin{lstlisting}
string Bigadd(string x,string y){
    //if(x[0]!='-'&&y[0]=='-')
    //    return Bigsub(x,y.substr(1,y.size()-1));
    //if(x[0]=='-'&&y[0]!='-')
    //    return Bigsub(y,x.substr(1,x.size()-1));
    if(x[0]=='-'&&y[0]=='-') return "-"+Bigadd(
        y.substr(1,y.size()-1),x.substr(1,x.size()-1));
    int s=max(x.size(),y.size()),a[s+6],b[s+6];
    int res[s+6];
    memset(res,0,sizeof(res));
    memset(a,0,sizeof(a));
    memset(b,0,sizeof(b));
    for(int i=0;i<x.size();i++) a[i]=x[x.size()-1-i]-'0';
    for(int i=0;i<y.size();i++) b[i]=y[y.size()-1-i]-'0';
    for(int i=0;i<s;i++){
        res[i]+=a[i]+b[i];
        res[i+1]=res[i]/10;
        res[i]%=10;
    }
    string ans="";
    bool first=false;
    for(int i=s+1;i>=0;i--){
        if(res[i]==0&&!first) continue;
        else{
            ans+=res[i]+'0';
            first=true;
        }
    }
    if(ans=="") return "0";
    return ans;
}
string Bigsub(string x,string y){
    if(x[0]!='-'&&y[0]=='-')
        return Bigadd(x,y.substr(1,y.size()-1));
    if(x[0]=='-'&&y[0]!='-')
        return '-' + Bigadd(x.substr(1,x.size()-1),y);
    if(x[0]=='-'&&y[0]=='-') return Bigsub(
        y.substr(1,y.size()-1),x.substr(1,x.size()-1));
    bool neg=false;
    if(y.size()>x.size()){
        swap(x,y);
        neg=true;
    }
    if(y.size()==x.size()&&y>x){
        swap(x,y);//ensures y>x
        neg=true;
    }
    int s=max(x.size(),y.size()),a[s+6],b[s+6];
    int res[s+6];
    memset(res,0,sizeof(res));
    memset(a,0,sizeof(a));
    memset(b,0,sizeof(b));
    for(int i=0;i<x.size();i++) a[i]=x[x.size()-1-i]-'0';
    for(int i=0;i<y.size();i++) b[i]=y[y.size()-1-i]-'0';
    for(int i=0;i<s;i++){
        res[i]+=a[i]-b[i];
        while(res[i]<0){
            res[i]+=10;
            res[i+1]--;
        }
    }
    string ans="";
    if(neg) ans+="-";
    bool first=false;
    for(int i=s+1;i>=0;i--){
        if(res[i]==0&&!first) continue;
        else{
            ans+=res[i]+'0';
            first=true;
        }
    }
    if(ans==""||ans=="-") return "0";
    return ans;
}
// multiplication not done like this
string Bigmul(string x,string y){
    if(x=="0"||y=="0") return "0";
    if(x[0]!='-'&&y[0]=='-')
        return "-"+Bigmul(x,y.substr(1,y.size()-1));
    if(x[0]=='-'&&y[0]!='-')
        return "-"+Bigmul(y,x.substr(1,x.size()-1));
    if(x[0]=='-'&&y[0]=='-'){
        x=x.substr(1,x.size()-1);
        y=y.substr(1,y.size()-1);
    }
    int s=max(x.size(),y.size()),a[2*s+6],b[2*s+6];
    int res[2*s+6];
    memset(res,0,sizeof(res));
    memset(a,0,sizeof(a));
    memset(b,0,sizeof(b));
    for(int i=0;i<x.size();i++) a[i]=x[x.size()-1-i]-'0';
    for(int i=0;i<y.size();i++) b[i]=y[y.size()-1-i]-'0';
    for(int i=0;i<x.size();i++){
        for(int j=0;j<y.size();j++){
            res[i+j]+=b[j]*a[i];
            res[i+j+1]+=res[i+j]/10;//will not increase 2.
            res[i+j]%=10;
        }
    }
    string ans="";
    bool first=false;
    for(int i=2*s+5;i>=0;i--){
        if(res[i]==0&&!first) continue;
        else{
            ans+=res[i]+'0';
            first=true;
        }
    }
    if(ans=="") return "0";
    return ans;
}
string Bighalf(string x){
    int s=x.size(),a[s+6];
    int res[s+6];
    memset(res,0,sizeof(res));
    memset(a,0,sizeof(a));
    for(int i=0;i<x.size();i++) a[i]=x[i]-'0';
    for(int i=0;i<x.size();i++){
        res[i]+=a[i]/2;
        if(a[i]%2){
            a[i+1]+=10;
        }
    }
    string ans="";
    bool first=false;
    for(int i=0;i<s;i++){
        if(res[i]==0&&!first) continue;
        else{
            ans+=res[i]+'0';
            first=true;
        }
    }
    if(ans=="") return "0";
    return ans;
}
bool Scompare(string x,string y){
    if(x.size()<y.size()) return true;
    if(x.size()>y.size()) return false;
    return x<y;
}
bool Ecompare(string x,string y){
    if(x.size()<y.size()) return true;
    if(x.size()>y.size()) return false;
    return x<=y;
}
//can get Bigdiv in the same way
string BigMod(string x,string m){
    string start="0",end="1";
    for(int i=1;i<=x.size()+6;i++) end+="0";
    while(Scompare(start,end)){
        string temp=Bigadd(end,"1");
        string mid=Bigadd(start,Bighalf(Bigsub(temp,start)));
        if(Ecompare(Bigmul(mid,m),x)){
            start=mid;
        }else{
            end=Bigsub(mid,"1");
        }
    }
    string res=Bigsub(x,Bigmul(start,m));
    return res;
}
ll BigModSmall(string a,ll mod){
    ll x=0;
    for(int i=0;i<a.size();i++){
        x=((10*x)+(a[i]-'0'))%mod;
    }
    return x;
}
\end{lstlisting}

\entry{FFT / NTT (convolution)}{$O(n \log n)$ \quad mem $n$}{%
\code{conv(a, b)} multiplies two polynomials: \code{c[i] = sum a[j]*b[i-j]}, length
\code{a.size()+b.size()-1}. This is also ``for every pair of elements, bucket their
sum'' --- count pairs summing to $s$, subset-sum counts, string matching with
wildcards.\\[0.3ex]
\textbf{NTT}, exact, mod \code{P = 998244353} (a prime with $2^{23} \mid P-1$, root 3).
So the result must be wanted mod that prime, and the transform length must be a power
of two $\le 2^{23}$. Other friendly primes: $1004535809$ (root 3),
$167772161$ (root 3).\\[0.3ex]
If the true coefficients are needed exactly and exceed \code{P}, run it under two
different primes and \code{crt} the results. For real/complex input use a
\code{complex<double>} FFT instead --- but it loses precision above $\sim10^{14}$ in the
output, so prefer NTT whenever the answer is modular.}

%LABEL ntt
\begin{lstlisting}
const ll P = 998244353, G = 3;
ll pw(ll a, ll b, ll m = P) {
    ll r = 1; a %= m;
    for (; b; b >>= 1, a = a * a % m)
        if (b & 1) r = r * a % m;
    return r;
}
void ntt(vi& a, bool inv) {
    int n = (int)a.size();
    for (int i = 1, j = 0; i < n; i++) {
        int b = n >> 1;
        for (; j & b; b >>= 1) j ^= b;
        j ^= b;
        if (i < j) swap(a[i], a[j]);
    }
    for (int len = 1; len < n; len <<= 1) {
        ll w = pw(G, (P - 1) / (2 * len));
        if (inv) w = pw(w, P - 2);
        for (int i = 0; i < n; i += 2 * len) {
            ll cur = 1;
            rep(j, 0, len) {
                ll u = a[i+j], v = a[i+j+len] * cur % P;
                a[i+j] = (u + v) % P;
                a[i+j+len] = (u - v + P) % P;
                cur = cur * w % P;
            }
        }
    }
    if (inv) {
        ll ni = pw(n, P - 2);
        for (ll& x : a) x = x * ni % P;
    }
}
vi conv(vi a, vi b) {
    if (a.empty() || b.empty()) return {};
    int need = (int)(a.size() + b.size() - 1), n = 1;
    while (n < need) n <<= 1;
    a.resize(n); b.resize(n);
    ntt(a, false); ntt(b, false);
    rep(i, 0, n) a[i] = a[i] * b[i] % P;
    ntt(a, true);
    a.resize(need);
    return a;
}
\end{lstlisting}

\section{Algebra and numerics (KACTL)}

\entry{Discrete log, modular sqrt}{KACTL \quad $O(\sqrt m)$, $O(\log^2 p)$}{%
\code{modLog(a, b, m)}: the smallest $x > 0$ with $a^x \equiv b \pmod m$, or \code{-1}.
Any \code{m}, prime or not, \code{a} need not be coprime to it. Baby-step giant-step:
store $b\cdot a^j$ for $j \le \sqrt m$, then step by $a^{\sqrt m}$ and look each up.
$x = 0$ is never returned, so if $b \equiv 1$ is a valid answer at exponent 0, check that
yourself; \code{modLog(a, 1, m)} is the \emph{order} of \code{a}. Needs
$m < 3\cdot10^9$ (\code{e*a} must fit).\\[0.3ex]
\code{modSqrt(a, p)}: an $x$ with $x^2 \equiv a \pmod p$, \textbf{p prime}; the other
root is \code{p - x}. Returns \code{-1} if \code{a} is not a square (Euler's criterion,
$a^{(p-1)/2} \ne 1$) --- KACTL asserts instead. Tonelli--Shanks. Uses \code{power}
from \emph{Modular basics}.}

%DEP modbasics
%LABEL modlog
\begin{lstlisting}
ll modLog(ll a, ll b, ll m) {
    ll n = (ll) sqrt(m) + 1, e = 1, f = 1, j = 1;
    unordered_map<ll, ll> A;
    while (j <= n && (e = f = e * a % m) != b % m)
        A[e * b % m] = j++;
    if (e == b % m) return j;
    if (__gcd(m, e) == __gcd(m, b))
        rep(i,2,n+2) if (A.count(e = e * f % m))
            return n * i - A[e];
    return -1;
}
ll modSqrt(ll a, ll p) {          // KACTL: sqrt(a, p)
    a %= p; if (a < 0) a += p;
    if (a == 0) return 0;
    if (power(a, (p-1)/2, p) != 1) return -1;  // no root
    if (p % 4 == 3) return power(a, (p+1)/4, p);
    ll s = p - 1, n = 2;
    int r = 0, m;
    while (s % 2 == 0)
        ++r, s /= 2;
    // find a non-square mod p
    while (power(n, (p - 1) / 2, p) != p - 1) ++n;
    ll x = power(a, (s + 1) / 2, p);
    ll b = power(a, s, p), g = power(n, s, p);
    for (;; r = m) {
        ll t = b;
        for (m = 0; m < r && t != 1; ++m)
            t = t * t % p;
        if (m == 0) return x;
        ll gs = power(g, 1LL << (r - m - 1), p);
        g = gs * gs % p;
        x = x * gs % p;
        b = b * g % p;
    }
}
\end{lstlisting}

\entry{FFT, convolution mod any m}{KACTL \quad $O(n \log n)$}{%
Same job as the NTT entry, \code{c[i] = sum a[j]*b[i-j]}, for the two cases NTT cannot
do.\\[0.3ex]
\code{conv(a, b)} on \code{double}s: \textbf{exact integer} results without any
modulus, e.g.\ counting pairs whose sum is $s$ when the counts reach $10^{10}$ (past
NTT's prime). Round with \code{llround}. Safe while
$(\sum a_i^2 + \sum b_i^2)\log_2 N < 9\cdot10^{14}$.\\[0.3ex]
\code{convMod<M>(a, b)}: the answer modulo \textbf{any} \code{M} up to $\sim10^9$ ---
the usual case is $10^9{+}7$, where NTT does not work because $10^9{+}6$ has no large
power of two. Inputs must already be in \code{[0, M)}. It splits every value into
$\sqrt M$-sized halves, so it is about twice as slow as one NTT.\\[0.3ex]
KACTL names the complex type \code{C}; it is \code{cd} here because \code{C(n, k)} is
the binomial.}

%LABEL fft
\begin{lstlisting}
typedef complex<double> cd;
typedef vector<double> vd;
void fft(vector<cd>& a) {
    int n = (int)a.size(), L = 31 - __builtin_clz(n);
    static vector<complex<long double>> R(2, 1);
    static vector<cd> rt(2, 1);  // (^ 10% faster if double)
    for (static int k = 2; k < n; k *= 2) {
        R.resize(n); rt.resize(n);
        auto x = polar(1.0L, acos(-1.0L) / k);
        rep(i,k,2*k) rt[i] = R[i] = i&1 ? R[i/2] * x : R[i/2];
    }
    vi rev(n);
    rep(i,0,n) rev[i] = (rev[i / 2] | (i & 1) << L) / 2;
    rep(i,0,n) if (i < rev[i]) swap(a[i], a[rev[i]]);
    for (int k = 1; k < n; k *= 2)
        for (int i = 0; i < n; i += 2 * k) rep(j,0,k) {
            cd z = rt[j+k] * a[i+j+k];
            a[i + j + k] = a[i + j] - z;
            a[i + j] += z;
        }
}
vd conv(const vd& a, const vd& b) {
    if (a.empty() || b.empty()) return {};
    vd res(a.size() + b.size() - 1);
    int L = 32 - __builtin_clz((int)res.size()), n = 1 << L;
    vector<cd> in(n), out(n);
    copy(all(a), begin(in));
    rep(i,0,(ll)b.size()) in[i].imag(b[i]);
    fft(in);
    for (cd& x : in) x *= x;
    rep(i,0,n) out[i] = in[-i & (n - 1)] - conj(in[i]);
    fft(out);
    rep(i,0,(ll)res.size()) res[i] = imag(out[i]) / (4 * n);
    return res;
}
template<int M> vi convMod(const vi &a, const vi &b) {
    if (a.empty() || b.empty()) return {};
    vi res(a.size() + b.size() - 1);
    int B = 32 - __builtin_clz((int)res.size()), n = 1 << B;
    int cut = int(sqrt(M));
    vector<cd> L(n), R(n), outs(n), outl(n);
    rep(i,0,(ll)a.size())
        L[i] = cd((int)a[i] / cut, (int)a[i] % cut);
    rep(i,0,(ll)b.size())
        R[i] = cd((int)b[i] / cut, (int)b[i] % cut);
    fft(L), fft(R);
    rep(i,0,n) {
        int j = -i & (n - 1);
        outl[j] = (L[i] + conj(L[j])) * R[i] / (2.0 * n);
        outs[j] = (L[i] - conj(L[j])) * R[i] / (2.0 * n) / 1i;
    }
    fft(outl), fft(outs);
    rep(i,0,(ll)res.size()) {
        ll av = ll(real(outl[i])+.5);
        ll cv = ll(imag(outs[i])+.5);
        ll bv = ll(imag(outl[i])+.5) + ll(real(outs[i])+.5);
        res[i] = ((av % M * cut + bv) % M * cut + cv) % M;
    }
    return res;
}
\end{lstlisting}

\entry{AND / OR / XOR convolution}{KACTL \quad $O(n \log n)$}{%
\code{c[z] = sum over x op y == z of a[x]*b[y]}, with op one of AND, OR, XOR --- the
bitwise cousin of FFT. Typical: ``how many pairs have XOR exactly $k$'' (put the counts
in \code{a} and \code{b}, read \code{c[k]}), or combining two independent choices
described by bitmasks. Pad both to the same power of two.\\[0.3ex]
\textbf{Pick the op by leaving exactly one of the three \code{tie} lines active}; for
XOR also uncomment the final division. Mod a prime: reduce after every \code{+}/\code{-},
and for XOR multiply by $n^{-1}$ instead of dividing. Called \code{conv} in KACTL,
renamed so it does not clash with NTT's.}

%LABEL fst
\begin{lstlisting}
void FST(vi& a, bool inv) {
    for (int n = (int)a.size(), step = 1; step < n; step *= 2)
        for (int i = 0; i < n; i += 2*step) rep(j,i,i+step) {
            ll &u = a[j], &v = a[j + step]; tie(u, v) =
                inv ? pii(v - u, u) : pii(v, u + v); // AND
             // inv ? pii(v, u - v) : pii(u + v, u); // OR
             // pii(u + v, u - v);                   // XOR
        }
    // if (inv) for (ll& x : a) x /= (ll)a.size(); // XOR only
}
vi fstConv(vi a, vi b) {
    FST(a, 0); FST(b, 0);
    rep(i,0,(ll)a.size()) a[i] *= b[i];
    FST(a, 1); return a;
}
\end{lstlisting}

\entry{Berlekamp--Massey + k-th term}{KACTL \quad $O(n^2)$, $O(n^2 \log k)$}{%
The ``I can compute the first 50 answers but need the $10^{18}$-th'' tool. If a sequence
mod a prime satisfies \emph{some} linear recurrence
$S_i = \sum_j tr_j\, S_{i-j-1}$ of order $\le n$, then \code{berlekampMassey} recovers
\code{tr} from the first $2n$ terms. Any DP whose step is a fixed linear map on a
fixed-size state has one (tilings of a $k\times n$ board, walks in a fixed graph,
counting strings avoiding a pattern), of order at most the state size, so you never derive
the matrix: brute-force the first terms, feed them in.\\[0.3ex]
Then \code{linearRec(S, tr, k)} gives term $k$ (0-indexed) from the first
\code{tr.size()} terms --- faster than matrix power. Example:
\code{berlekampMassey(\{0,1,1,2,3,5\})} $=$ \code{\{1,1\}}, and
\code{linearRec(\{0,1\}, \{1,1\}, k)} is Fibonacci.\\[0.3ex]
Give it \textbf{plenty} of terms: if the result's length is close to half of what you
supplied, the recurrence may be longer --- add terms until it stops changing. An all-zero
input gives an empty \code{tr}; handle that before calling \code{linearRec}. Uses
\code{power} from \emph{Modular basics}; \code{mod} is shared with the next entry.}

%DEP modbasics
%LABEL bm
\begin{lstlisting}
const ll mod = 1000000007;          // must be prime for BM
vi berlekampMassey(vi s) {
    int n = (int)s.size(), L = 0, m = 0;
    vi C(n), B(n), T;
    C[0] = B[0] = 1;
    ll b = 1;
    rep(i,0,n) { ++m;
        ll d = s[i] % mod;
        rep(j,1,L+1) d = (d + C[j] * s[i - j]) % mod;
        if (!d) continue;
        T = C; ll coef = d * power(b, mod-2, mod) % mod;
        rep(j,m,n) C[j] = (C[j] - coef * B[j - m]) % mod;
        if (2 * L > i) continue;
        L = i + 1 - L; B = T; b = d; m = 0;
    }
    C.resize(L + 1); C.erase(C.begin());
    for (ll& x : C) x = (mod - x) % mod;
    return C;
}
ll linearRec(vi S, vi tr, ll k) {
    int n = (int)tr.size();
    auto combine = [&](vi a, vi b) {
        vi res(n * 2 + 1);
        rep(i,0,n+1) rep(j,0,n+1)
            res[i + j] = (res[i + j] + a[i] * b[j]) % mod;
        for (int i = 2 * n; i > n; --i) rep(j,0,n)
            res[i-1-j] = (res[i-1-j] + res[i] * tr[j]) % mod;
        res.resize(n + 1);
        return res;
    };
    vi pol(n + 1), e(pol);
    pol[0] = e[1] = 1;
    for (++k; k; k /= 2) {
        if (k % 2) pol = combine(pol, e);
        e = combine(e, e);
    }
    ll res = 0;
    rep(i,0,n) res = (res + pol[i + 1] * S[i]) % mod;
    return res;
}
\end{lstlisting}

\entry{Determinant}{KACTL \quad $O(n^3)$}{%
Both versions \textbf{destroy} the matrix --- pass a copy if you need it again.
\code{double}: Gaussian elimination with partial pivoting. \code{ll}: the exact answer
mod \code{mod}, and \code{mod} need \emph{not} be prime, because it eliminates with
Euclid-style row subtractions instead of inverses (remove every \code{\% mod} for plain
integers if nothing overflows).\\[0.3ex]
Main use, \textbf{Kirchhoff's matrix-tree theorem}: the number of spanning trees is the
determinant of the Laplacian (\code{L[v][v]} = degree, \code{L[u][v]} = $-$(number of
edges $uv$)) with any one row and the same column deleted. Also the
Lindstr\"om--Gessel--Viennot lemma: the number of vertex-disjoint path systems in a DAG
is the determinant of the path-count matrix.}

%DEP modbasics bm
%LABEL det
\begin{lstlisting}
double det(vector<vector<double>>& a) {
    int n = (int)a.size(); double res = 1;
    rep(i,0,n) {
        int b = i;
        rep(j,i+1,n) if (fabs(a[j][i]) > fabs(a[b][i])) b = j;
        if (i != b) swap(a[i], a[b]), res *= -1;
        res *= a[i][i];
        if (res == 0) return 0;
        rep(j,i+1,n) {
            double v = a[j][i] / a[i][i];
            if (v != 0) rep(k,i+1,n) a[j][k] -= v * a[i][k];
        }
    }
    return res;
}
ll det(vector<vector<ll>>& a) {          // mod `mod`
    int n = (int)a.size(); ll ans = 1;
    rep(i,0,n) {
        rep(j,i+1,n) {
            while (a[j][i] != 0) { // gcd step
                ll t = a[i][i] / a[j][i];
                if (t) rep(k,i,n)
                    a[i][k] = (a[i][k] - a[j][k] * t) % mod;
                swap(a[i], a[j]);
                ans *= -1;
            }
        }
        ans = ans * a[i][i] % mod;
        if (!ans) return 0;
    }
    return (ans + mod) % mod;
}
\end{lstlisting}

\entry{Simplex (linear programming)}{KACTL \quad fast in practice, $2^n$ worst}{%
Maximise $c^T x$ subject to $Ax \le b$, $x \ge 0$ (real $x$). Returns the optimum and
sets \code{x}; \code{-inf} if infeasible, \code{inf} if unbounded. Recognise it when the
problem is ``choose real amounts, linear limits, linear objective'' (diets, mixing,
fractional scheduling) and $n, m \lesssim$ a few hundred.\\[0.3ex]
Rewriting into that form: \textbf{minimise} $\to$ negate \code{c} and the result;
$\ge$ row $\to$ negate the row and its \code{b}; \textbf{equality} $\to$ both $\le$ and
$\ge$; a \textbf{free} variable $\to$ $x = x^+ - x^-$, two columns. The answer is a
\code{double}; it is not an integer program, rounding \code{x} is not valid in
general.\\[0.3ex]
Usage: \code{vvd A = \{\{1,-1\},\{-1,1\},\{-1,-2\}\}; vd b = \{1,1,-4\}, c = \{-1,-1\}, x;}
\code{T val = LPSolver(A, b, c).solve(x);}\\[0.3ex]
\code{eps} renamed to \code{EPS} because the header defines \code{eps}.}

%LABEL simplex
\begin{lstlisting}
typedef double T; // long double, Rational, double + mod<P>
typedef vector<T> vd;
typedef vector<vd> vvd;
const T EPS = 1e-8, inf = 1/.0;
#define MP make_pair
#define ltj(X) if(s==-1 || MP(X[j],N[j]) < MP(X[s],N[s])) s=j
struct LPSolver {
    int m, n;
    vi N, B;
    vvd D;
    LPSolver(const vvd& A, const vd& b, const vd& c) :
        m((int)b.size()), n((int)c.size()),
        N(n+1), B(m), D(m+2, vd(n+2)) {
        rep(i,0,m) rep(j,0,n) D[i][j] = A[i][j];
        rep(i,0,m) { B[i] = n+i; D[i][n] = -1;
                     D[i][n+1] = b[i]; }
        rep(j,0,n) { N[j] = j; D[m][j] = -c[j]; }
        N[n] = -1; D[m+1][n] = 1;
    }
    void pivot(int r, int s) {
        T *a = D[r].data(), inv = 1 / a[s];
        rep(i,0,m+2) if (i != r && abs(D[i][s]) > EPS) {
            T *b = D[i].data(), inv2 = b[s] * inv;
            rep(j,0,n+2) b[j] -= a[j] * inv2;
            b[s] = a[s] * inv2;
        }
        rep(j,0,n+2) if (j != s) D[r][j] *= inv;
        rep(i,0,m+2) if (i != r) D[i][s] *= -inv;
        D[r][s] = inv;
        swap(B[r], N[s]);
    }
    bool simplex(int phase) {
        int x = m + phase - 1;
        for (;;) {
            int s = -1;
            rep(j,0,n+1) if (N[j] != -phase) ltj(D[x]);
            if (D[x][s] >= -EPS) return true;
            int r = -1;
            rep(i,0,m) {
                if (D[i][s] <= EPS) continue;
                if (r == -1 || MP(D[i][n+1] / D[i][s], B[i])
                            < MP(D[r][n+1] / D[r][s], B[r]))
                    r = i;
            }
            if (r == -1) return false;
            pivot(r, s);
        }
    }
    T solve(vd &x) {
        int r = 0;
        rep(i,1,m) if (D[i][n+1] < D[r][n+1]) r = i;
        if (D[r][n+1] < -EPS) {
            pivot(r, n);
            if (!simplex(2) || D[m+1][n+1] < -EPS)
                return -inf;
            rep(i,0,m) if (B[i] == -1) {
                int s = 0;
                rep(j,1,n+1) ltj(D[i]);
                pivot(i, s);
            }
        }
        bool ok = simplex(1); x = vd(n);
        rep(i,0,m) if (B[i] < n) x[B[i]] = D[i][n+1];
        return ok ? D[m][n+1] : inf;
    }
};
\end{lstlisting}

\entry{Numerical integration}{KACTL \quad $O(n)$ evaluations}{%
Simpson's rule on $2n$ slices, error $\propto h^4$. For areas with no closed form:
union of circles (integrate the covered length of each vertical line), volumes of
revolution. Re-run with a bigger \code{n} and check the digits you need stay put; if
\code{f} has kinks (e.g.\ where circles start and end), integrate piecewise between
them instead of one big call.}

%LABEL quad
\begin{lstlisting}
template<class F>
double quad(double a, double b, F f, const int n = 1000) {
    double h = (b - a) / 2 / n, v = f(a) + f(b);
    rep(i,1,n*2)
        v += f(a + i*h) * (i&1 ? 4 : 2);
    return v * h / 3;
}
\end{lstlisting}

\section{Strings}

\entry{KMP}{$O(n + m)$ \quad mem $n$}{%
\code{failure(p)[i]} is the length of the longest proper prefix of \code{p[0..i]} that is
also a suffix --- so \code{m - failure.back()} is the shortest period of \code{p}, and
\code{p} is a repetition of that period iff the period divides \code{m}.
\code{kmp(txt, pat)} returns every 0-indexed start where \code{pat} occurs, overlaps
included.\\[0.3ex]
Borders of \code{p} are \code{m}, \code{f[m-1]}, \code{f[f[m-1]-1]}, \dots\ --- walk the
chain. Counting occurrences of every prefix: add 1 at each \code{f} position and push
the counts down the chain in reverse order.}

%LABEL kmp
\begin{lstlisting}
vector<int> failure(const string& p) {
    int m = (int)p.size(); vector<int> f(m, 0);
    for (int i = 1, k = 0; i < m; i++) {
        while (k && p[i] != p[k]) k = f[k-1];
        if (p[i] == p[k]) k++;
        f[i] = k;
    }
    return f;
}
vector<int> kmp(const string& txt, const string& pat) {
    vector<int> r;
    if (pat.empty()) return r;
    vector<int> f = failure(pat);
    int n = (int)txt.size(), m = (int)pat.size();
    for (int i = 0, k = 0; i < n; i++) {
        while (k && txt[i] != pat[k]) k = f[k-1];
        if (txt[i] == pat[k]) k++;
        if (k == m) { r.push_back(i - m + 1); k = f[k-1]; }
    }
    return r;
}
\end{lstlisting}

\entry{Z function}{$O(n)$ \quad mem $n$}{%
\code{z[i]} is the length of the longest common prefix of \code{s} and \code{s[i..]}
(\code{z[0]} is left 0). Pattern matching: build \code{pat + '\textbackslash1' + txt} and
look for \code{z[i] == pat.size()}; the separator must not occur in either
string.\\[0.3ex]
Often replaces KMP, and is the easier one to get right. Distinct substrings, smallest
rotation, and ``is this a repetition'' all fall out of it.}

%LABEL zfunc
\begin{lstlisting}
vector<int> zfunc(const string& s) {
    int n = (int)s.size();
    vector<int> z(n, 0);
    for (int i = 1, l = 0, r = 0; i < n; i++) {
        if (i < r) z[i] = min(r - i, z[i - l]);
        while (i + z[i] < n && s[z[i]] == s[i + z[i]]) z[i]++;
        if (i + z[i] > r) { l = i; r = i + z[i]; }
    }
    return z;
}
\end{lstlisting}

\entry{Manacher}{$O(n)$ \quad mem $n$}{%
Every palindromic substring at once. \code{d1[i]} is the radius of the longest
\emph{odd} palindrome centred at \code{i} (always $\ge 1$), \code{d2[i]} the radius of
the longest \emph{even} one centred between \code{i-1} and \code{i} (may be 0). The
number of palindromic substrings is $\sum d1 + \sum d2$.\\[0.3ex]
\code{isPal(l, r)} answers ``is \code{s[l..r)} a palindrome'' in $O(1)$, which is what
you usually actually want --- e.g. for ``split the string into the fewest
palindromes'' DP. Longest palindromic substring:
$\max(2 d1[i] - 1, 2 d2[i])$.}

%LABEL manacher
\begin{lstlisting}
struct Manacher {
    int n; vector<int> d1, d2;
    Manacher(const string& s) : n((int)s.size()),
                                d1(n, 0), d2(n, 0) {
        for (int i = 0, l = 0, r = -1; i < n; i++) {
            int k = i > r ? 1 : min(d1[l+r-i], r-i+1);
            while (i-k >= 0 && i+k < n && s[i-k] == s[i+k])
                k++;
            d1[i] = k--;
            if (i + k > r) { l = i - k; r = i + k; }
        }
        for (int i = 0, l = 0, r = -1; i < n; i++) {
            int k = i > r ? 0 : min(d2[l+r-i+1], r-i+1);
            while (i - k - 1 >= 0 && i + k < n
                   && s[i-k-1] == s[i+k]) k++;
            d2[i] = k--;
            if (i + k > r) { l = i - k - 1; r = i + k; }
        }
    }
    bool isPal(int l, int r) {                // s[l..r)
        int len = r - l, c = l + len / 2;
        return len % 2 ? 2*d1[c]-1 >= len : 2*d2[c] >= len;
    }
};
\end{lstlisting}

\entry{String hashing}{build $O(n)$, compare $O(1)$ \quad mem $n$}{%
\code{get(l, r)} is the hash of \code{s[l..r)}; two substrings are equal iff their
hashes match. Two independent moduli, so a collision needs both to fail at once ---
with a single 64-bit overflow hash an adversarial test \emph{will} break you.\\[0.3ex]
This is the workhorse: compare any two substrings, binary-search the longest common
prefix of two suffixes, dedupe substrings by putting \code{get} results in a
\code{set}, Rabin-Karp matching. Often replaces a suffix array
outright.\\[0.3ex]
For randomised bases (safer against hash-killers) pick \code{B1}, \code{B2} at run time
from \code{mt19937} in $[256, M)$.}

%LABEL hash
\begin{lstlisting}
struct Hash {
    static constexpr ll M1 = 1000000007, M2 = 998244353,
                        B1 = 131, B2 = 137;
    int n; vi h1, h2, p1, p2;
    Hash(const string& s) : n((int)s.size()), h1(n+1, 0),
            h2(n+1, 0), p1(n+1, 1), p2(n+1, 1) {
        rep(i, 0, n) {
            h1[i+1] = (h1[i] * B1 + s[i]) % M1;
            h2[i+1] = (h2[i] * B2 + s[i]) % M2;
            p1[i+1] = p1[i] * B1 % M1;
            p2[i+1] = p2[i] * B2 % M2;
        }
    }
    pii get(int l, int r) {                   // s[l..r)
        ll a = (h1[r] - h1[l] * p1[r-l]) % M1;
        ll b = (h2[r] - h2[l] * p2[r-l]) % M2;
        return {(a + M1) % M1, (b + M2) % M2};
    }
};
\end{lstlisting}

\entry{Suffix array + LCP}{$O(n \log n)$ \quad mem $n$}{%
\code{sa[i]} is the start of the \code{i}-th smallest suffix; \code{rnk} is its inverse;
\code{lcp[i]} is the longest common prefix of \code{sa[i]} and \code{sa[i+1]}.
The input must not contain the byte \code{1}, which is used as a
sentinel.\\[0.3ex]
Number of distinct substrings is $\frac{n(n+1)}{2} - \sum \mathrm{lcp}$. Longest
repeated substring is $\max \mathrm{lcp}$. Longest common substring of two strings: join
with a separator and take the best \code{lcp} between suffixes from different sides.
LCP of two arbitrary suffixes = range-min of \code{lcp} between their ranks --- feed
\code{lcp} to \code{SpTab}.\\[0.3ex]
Consider \code{Hash} plus binary search first: shorter to write, usually fast enough.}

%LABEL sufarr
\begin{lstlisting}
struct SuffixArray {
    int n; vector<int> sa, rnk, lcp;
    SuffixArray(const string& str) {
        string s = str; s += char(1);       // sentinel
        n = (int)s.size();
        vector<int> p(n), c(n), pn(n), cn(n);
        vector<int> cnt(max(256, n), 0);
        for (int i = 0; i < n; i++)
            cnt[(unsigned char)s[i]]++;
        for (int i = 1; i < 256; i++) cnt[i] += cnt[i-1];
        for (int i = 0; i < n; i++)
            p[--cnt[(unsigned char)s[i]]] = i;
        int cls = 1; c[p[0]] = 0;
        for (int i = 1; i < n; i++) {
            if (s[p[i]] != s[p[i-1]]) cls++;
            c[p[i]] = cls - 1;
        }
        for (int h = 0; (1 << h) < n; h++) {
            for (int i = 0; i < n; i++) {
                pn[i] = p[i] - (1 << h);
                if (pn[i] < 0) pn[i] += n;
            }
            fill(cnt.begin(), cnt.begin() + cls, 0);
            for (int i = 0; i < n; i++) cnt[c[pn[i]]]++;
            for (int i = 1; i < cls; i++) cnt[i] += cnt[i-1];
            for (int i = n - 1; i >= 0; i--)
                p[--cnt[c[pn[i]]]] = pn[i];
            cn[p[0]] = 0; cls = 1;
            for (int i = 1; i < n; i++) {
                int a = p[i], b = p[i-1], d = 1 << h;
                if (c[a] != c[b]
                    || c[(a+d) % n] != c[(b+d) % n]) cls++;
                cn[a] = cls - 1;
            }
            c = cn;
        }
        sa.assign(p.begin() + 1, p.end());   // drop sentinel
        n--;
        rnk.assign(n, 0);
        for (int i = 0; i < n; i++) rnk[sa[i]] = i;
        lcp.assign(max(n - 1, 0), 0);
        for (int i = 0, k = 0; i < n; i++) {
            if (rnk[i] == n - 1) { k = 0; continue; }
            int j = sa[rnk[i] + 1];
            while (i + k < n && j + k < n
                   && str[i+k] == str[j+k]) k++;
            lcp[rnk[i]] = k;
            if (k) k--;
        }
    }
};
\end{lstlisting}

\entry{Aho-Corasick}{build $O(\Sigma \cdot \mathrm{tot})$, scan $O(n + \mathrm{hits})$ \quad mem $\Sigma \cdot \mathrm{tot}$}{%
Many patterns at once. \code{add(s, id)} before \code{build()}, then \code{match(txt, f)}
calls \code{f(i, id)} for every pattern \code{id} ending at text position \code{i}.
Lowercase \code{a-z}; change the 26 and the \code{- 'a'} for another
alphabet.\\[0.3ex]
\code{out} is the \emph{output link} --- the nearest fail-ancestor that is itself a
pattern end --- so reporting costs $O(1)$ per hit instead of walking the whole fail
chain, which is the difference between $O(n + \mathrm{hits})$ and $O(n \cdot
\mathrm{depth})$. Empty patterns are not allowed (the root is never a match).\\[0.3ex]
To \emph{count} hits per pattern without listing them, mark \code{cnt[u]++} per text
position and push counts along \code{fail} in reverse BFS order --- $O(n +
\mathrm{tot})$ regardless of how many matches there are.}

%LABEL aho
\begin{lstlisting}
struct Aho {
    struct Nd { int go[26], fail = 0, out = 0, head = -1; };
    vector<Nd> t;
    vector<int> idv, nxt;        // id lists on terminal nodes
    Aho() { newNode(); }
    int newNode() {
        t.push_back(Nd{});
        fill(t.back().go, t.back().go + 26, -1);
        return (int)t.size() - 1;
    }
    void add(const string& s, int id) {  // nonempty s
        int u = 0;
        for (char ch : s) {
            int c = ch - 'a';
            if (t[u].go[c] < 0) {
                int v = newNode(); t[u].go[c] = v;
            }
            u = t[u].go[c];
        }
        idv.push_back(id); nxt.push_back(t[u].head);
        t[u].head = (int)idv.size() - 1;
    }
    void build() {
        vector<int> q;
        for (int c = 0; c < 26; c++) {
            int v = t[0].go[c];
            if (v < 0) t[0].go[c] = 0;
            else { t[v].fail = 0; q.push_back(v); }
        }
        for (int i = 0; i < (int)q.size(); i++) {
            int u = q[i], f = t[u].fail;
            t[u].out = t[f].head >= 0 ? f : t[f].out;
            for (int c = 0; c < 26; c++) {
                int v = t[u].go[c], w = t[f].go[c];
                if (v < 0) t[u].go[c] = w;
                else { t[v].fail = w; q.push_back(v); }
            }
        }
    }
    template<class F>       // f(end index, pattern id)
    void match(const string& s, F f) {
        int u = 0;
        for (int i = 0; i < (int)s.size(); i++) {
            u = t[u].go[s[i] - 'a'];
            for (int v = t[u].head >= 0 ? u : t[u].out; v;
                 v = t[v].out)
                for (int e = t[v].head; e >= 0; e = nxt[e])
                    f(i, idv[e]);
        }
    }
};
\end{lstlisting}

\entry{XOR trie}{insert / query $O(\mathrm{BITS})$ \quad mem $n \cdot \mathrm{BITS}$}{%
\code{maxXor(x)} is the largest \code{x \textasciicircum{} y} over everything inserted:
walk from the top bit and always take the opposite branch if it exists. Raise
\code{BITS} to 62 for \code{ll} values.\\[0.3ex]
Keep a counter per node and you can also answer ``how many \code{y} with
\code{x \textasciicircum{} y < k}'', and delete by decrementing. Maximum XOR
\emph{subarray}: insert prefix xors as you sweep.\\[0.3ex]
A plain prefix trie over letters is the same shape with 26 children --- see
\code{Aho} without the fail links.}

\begin{lstlisting}
struct XorTrie {
    static const int BITS = 30;
    vector<array<int,2>> ch{{-1, -1}};
    void insert(ll x) {
        int u = 0;
        down(b, BITS + 1, 0) {
            int c = x >> b & 1;
            if (ch[u][c] < 0) {
                ch.push_back({-1, -1});
                ch[u][c] = (int)ch.size() - 1;
            }
            u = ch[u][c];
        }
    }
    ll maxXor(ll x) {         // needs >= 1 inserted value
        int u = 0; ll r = 0;
        down(b, BITS + 1, 0) {
            int c = (x >> b & 1) ^ 1;
            if (ch[u][c] >= 0) { r |= 1LL<<b; u = ch[u][c]; }
            else u = ch[u][c ^ 1];
        }
        return r;
    }
};
\end{lstlisting}

\entry{Minimum rotation}{KACTL \quad $O(n)$}{%
Index where the lexicographically smallest rotation starts:
\code{rotate(s.begin(), s.begin() + minRotation(s), s.end())}. Canonical form for
cyclic strings and necklaces --- two strings are rotations of each other iff their
minimal rotations are equal.}

%LABEL minrot
\begin{lstlisting}
int minRotation(string s) {
    int a=0, N=(int)s.size(); s += s;
    rep(b,0,N) rep(k,0,N) {
        if (a+k == b || s[a+k] < s[b+k]) {
            b += max(0LL, k-1); break; }
        if (s[a+k] > s[b+k]) { a = b; break; }
    }
    return a;
}
\end{lstlisting}

\section{Geometry}

\entry{Point / vector}{Ruiming}{%
Ruiming's geometry base: every other entry here builds on it, so it is printed once.
Floating point, compared through \code{ldcmp(x, y)}, which returns \code{-1}, \code{0},
\code{1} with tolerance \code{eps}. His geometry files use \code{eps = 1e-6} (the header
has \code{1e-9}) --- set it in \code{solve()}; for large coordinates make it relative:
\code{eps = max(eps, 1e-6 * maxabs)}.\\[0.3ex]
\begin{ctab}{0.27}{0.68}
\code{Cross(A,B)} & $A_x B_y - A_y B_x$. Sign = which side; $|\cdot|$ = parallelogram
area.\\
\code{Area(A,B,C)} & signed triangle area, $>0$ if counter-clockwise.\\
\code{Angle(A,B)} & unsigned, in $[0,\pi]$.\\
\code{DirectedAngle} & signed, in $(-\pi,\pi]$, from \code{A} to \code{B}.\\
\code{Rotate(A,rad)} & counter-clockwise.\\
\end{ctab}\\[0.3ex]
\code{POF} is the ``no point'' sentinel \code{Cross\_Segment} compares against.}

%LABEL point
%FROM Ruiming/geomtery/Line_and_Vectors.cpp
\begin{lstlisting}
const ld pi=acos(-1.0);
//eps=1e-6;//MIGHT GET WA PGA EPS TOO SMALL!!!!!!!!!!!!!!!!!
ll ldcmp(ld x,ld y){
    if(fabs(x-y)<eps) return 0;
    else return x<y?-1:1;
    //is x<y return -1, x>y return 1;
}
struct Point{
    ld x,y;
    Point(): x(0),y(0){}
    Point(ld x,ld y): x(x),y(y) {}
    Point operator + (Point B){return Point(x+B.x,y+B.y);}
    Point operator - (Point B){return Point(x-B.x,y-B.y);}
    Point operator * (ld k){return Point(x*k,y*k);}
    Point operator / (ld k){return Point(x/k,y/k);}
    bool operator == (Point B){
        return ((ldcmp(x-B.x,0)==0)&&(ldcmp(y-B.y,0)==0));}
    bool operator < (Point B){
        if (x != B.x) return x < B.x;
        return y < B.y;
    }
};
Point POF=Point(INF,INF);
ld Dist(Point A,Point B){
    return sqrt((A.x-B.x)*(A.x-B.x)+(A.y-B.y)*(A.y-B.y));
}
typedef Point Vector;
ld Dot(Vector A,Vector B){return A.x*B.x+A.y*B.y;}
ld Len(Vector A){return sqrt(Dot(A,A));}
ld Cross(Vector A,Vector B){return A.x*B.y-A.y*B.x;}
bool Parallel(Vector A,Vector B){
    return ldcmp(Cross(A,B),0)==0;}
ld Angle(Vector A,Vector B){
    if(Parallel(A,B)) return 0;
    return acos(Dot(A,B)/Len(A)/Len(B));
}
ld DirectedAngle(Vector A, Vector B) {
    return atan2l(Cross(A,B), Dot(A,B));
}
ld Area(Point A,Point B,Point C){return Cross(B-A,C-A)/2;}
Vector Rotate(Vector A,ld rad){
    return Vector(A.x*cos(rad)-A.y*sin(rad),
                  A.x*sin(rad)+A.y*cos(rad));
}
\end{lstlisting}

\entry{Lines and segments}{Ruiming}{%
A \code{Line} is two points \code{p1}, \code{p2}; the same struct is used as a
\emph{segment}. Build it from two points, from a point and an angle, or from
$ax+by+c=0$.\\[0.3ex]
\code{Point\_line\_relation}: \code{1} left of the directed line \code{p1->p2},
\code{2} right, \code{0} on it. \code{Line\_relation}: \code{0} parallel, \code{1} the
same line, \code{2} intersecting.\\[0.3ex]
\code{Cross\_point(a,b,c,d)} intersects the \emph{lines} \code{ab} and \code{cd}; it
divides by zero when they are parallel, so check \code{Line\_relation} first.
\code{Cross\_Segment} handles the collinear cases through its endpoint checks.}

%DEP point
%LABEL seg
%FROM Ruiming/geomtery/Line_and_Vectors.cpp
\begin{lstlisting}
struct Line{
    Point p1,p2;
    Line(){}
    Line(Point p1,Point p2): p1(p1),p2(p2){}
    Line(Point P,ld angle){
        p1=P;
        if(ldcmp(angle-pi/2,0)==0){
            p2=p1+Point(0,1);
        }else{
            p2=p1+Point(1,tan(angle));
        }
    }
    Line(ld a,ld b,ld c){
        if(ldcmp(a,0)==0){
            p1=Point(0,-c/b);
            p2=Point(1,-c/b);
        }else if(ldcmp(b,0)==0){
            p1=Point(-c/a,0);
            p2=Point(-c/a,1);
        }else{
            p1=Point(0,-c/b);
            p2=Point(1,(-c-a)/b);
        }
    }
};
//for directed lines from p1 to p2
ll Point_line_relation(Point P, Line v){
    ll c=ldcmp(Cross(P-v.p1,v.p2-v.p1),0);
    if(c<0) return 1;//P on the leftside of v
    if(c>0) return 2;//P on the rightside of v
    return 0; //P on v
}
bool Point_on_seg(Point P, Line v){
    return ldcmp(Cross(P-v.p1,v.p2-v.p1),0)==0
         &&ldcmp(Dot(P-v.p1,P-v.p2),0)<=0;
}
ld Dis_point_line(Point P, Line v){
    return fabs(Cross(P-v.p1,v.p2-v.p1)/Dist(v.p1,v.p2));
}
Point Point_line_proj(Point P, Line v){
    ld k=Dot(v.p2-v.p1,P-v.p1)/Dot(v.p2-v.p1,v.p2-v.p1);
    return v.p1+(v.p2-v.p1)*k;
}
Point Point_line_symmetry(Point P, Line v){
    Point q=Point_line_proj(P,v);
    return Point(2*q.x-P.x,2*q.y-P.y);
}
ld Dis_point_seg(Point P, Line v){
    if(ldcmp(Dot(P-v.p1,v.p2-v.p1),0)<0
     ||ldcmp(Dot(P-v.p2,v.p1-v.p2),0)<0){
        return min(Dist(P,v.p1),Dist(P,v.p2));
    }
    return Dis_point_line(P,v);
}
//Directed lines
ll Line_relation(Line v1, Line v2){
    if(ldcmp(Cross(v1.p2-v1.p1,v2.p2-v2.p1),0)==0){
        //overlapp
        if(Point_line_relation(v1.p1,v2)==0) return 1;
        //parallell
        else return 0;
    }
    //intersect
    return 2;
}
//Must have ab,cd not the same
Point Cross_point(Point a,Point b,Point c,Point d){
    ld s1=Cross(b-a,c-a);
    ld s2=Cross(b-a,d-a);
    return Point(c.x*s2-d.x*s1,c.y*s2-d.y*s1)/(s2-s1);
}
bool Cross_Segment(Point a,Point b,Point c,Point d){
    Point P=Cross_point(a,b,c,d);
    if(P==POF) return 0;
    //special case if all share a line, in line case
    //whatever works, not in segment case
    if(Point_on_seg(a,Line(c,d))||Point_on_seg(b,Line(c,d))||
       Point_on_seg(c,Line(a,b))||Point_on_seg(d,Line(a,b)))
        return 1;
    if(Point_on_seg(P,Line(a,b))&&Point_on_seg(P,Line(c,d)))
        return 1;
    else return 0;
}
\end{lstlisting}

\entry{Polygons}{Ruiming \quad $O(n)$}{%
\textbf{1-indexed}: the polygon is \code{p[1..n]}, in order (either direction).
\code{Point\_in\_polygon} returns \code{3} on a vertex, \code{2} on an edge,
\code{1} inside, \code{0} outside; it works for non-convex polygons (winding number).
\code{Polygon\_area} returns the \emph{absolute} area.\\[0.3ex]
\textbf{Pick's theorem} (lattice polygon): $A = I + B/2 - 1$, with the boundary count
$B = \sum \gcd(|dx|, |dy|)$ over the edges --- so the interior count is
$I = (\texttt{Polygon\_area2} - B + 2)/2$ using the exact version in
\emph{Lattice points}.}

%DEP point seg
%LABEL poly
%FROM Ruiming/geomtery/Polygon.cpp
\begin{lstlisting}
//list of points p representing polygon
ll Point_in_polygon(Point pt, Point *p,ll n){
    for(int i=1;i<=n;i++){
        //Point on vetices
        if(ldcmp(Dist(p[i],pt),0)==0) return 3;
    }
    for(int i=1;i<=n;i++){
        Line v=Line(p[i],p[i%n+1]);
        //Point on Edge
        if(Point_on_seg(pt,v)) return 2;
    }
    ll num=0;
    for(int i=1;i<=n;i++){
        ll j=i%n+1;
        ll c=ldcmp(Cross(pt-p[j],p[i]-p[j]),0);
        ll u=ldcmp(p[i].y-pt.y,0);
        ll v=ldcmp(p[j].y-pt.y,0);
        if(c>0&&u<0&&v>=0) num++;
        if(c<0&&u>=0&&v<0) num--;
    }
    //return 1 is inside, 0 is outside
    return num!=0;
}
ld Polygon_area(Point *p,ll n){
    ld area=0;
    for(int i=1;i<=n;i++){
        area+=Cross(p[i],p[i%n+1]);
    }
    return abs(area/2.0);
}
\end{lstlisting}

\entry{Convex hull}{Ruiming \quad $O(n \log n)$}{%
Andrew's monotone chain, \textbf{1-indexed}: input \code{p[1..n]} (sorted and
deduplicated in place), output \code{ch[1..v]}, returns \code{v}. Collinear points on
the boundary are \textbf{kept}; change both \code{<0} to \code{<=0} to drop them.\\[0.3ex]
The final \code{sort}/\code{unique} leaves \code{ch} sorted by $(x,y)$, \emph{not} in
hull order --- fine for counting or membership (``are all points on the hull'' is
\code{Convex\_Hull(p,n,ch)==n}). For the actual cycle, delete those two lines; then
\code{ch[v]} repeats \code{ch[1]}, so use \code{v-1} points.\\[0.3ex]
Diameter of a point set = farthest pair on the hull (rotating calipers, or $O(h^2)$ when
$h$ is small).}

%DEP point
%LABEL hull
%FROM Ruiming/geomtery/Convex_Hull.cpp
\begin{lstlisting}
ll Convex_Hull(Point *p, ll n,Point *ch){
    sort(p+1,p+n+1);
    n=unique(p+1,p+n+1)-(p+1);
    if (n == 1) { ch[1] = p[1]; return 1; }
    if (n == 2) { ch[1] = p[1]; ch[2] = p[2]; return 2; }
    ll v=0;
    for(int i=1;i<=n;i++){
        while(v>1&&
              ldcmp(Cross(ch[v]-ch[v-1],p[i]-ch[v]),0)<0){
            v--;
        }
        ch[++v]=p[i];
    }
    ll j=v;
    for(int i=n-1;i>=1;i--){
        //<=0 do not count points on the side of convex hull
        while(v>j&&
              ldcmp(Cross(ch[v]-ch[v-1],p[i]-ch[v]),0)<0){
            v--;
        }
        ch[++v]=p[i];
    }
    sort(ch+1,ch+v+1);
    v=unique(ch+1,ch+v+1)-(ch+1);
    //might include double points on the line case
    return v;
}
\end{lstlisting}

\entry{Half-plane intersection}{Ruiming \quad $O(n \log n)$}{%
Each \code{Line} keeps the half-plane on its \textbf{left} (\code{p1->p2}).
\code{HPI(L)} returns the vertices of the intersection, counter-clockwise, or
\code{\{\}} if it is empty or has fewer than 3 vertices. It assumes the result is
\textbf{bounded} --- add four big box edges if it might not be.\\[0.3ex]
This file has its \textbf{own \code{Line}} (adds direction \code{v}, angle \code{ang} and
\code{operator<}): use it \emph{instead of} the one in \emph{Lines and segments}. The
functions there are textually identical in his HPI file and work with this
\code{Line} too. (Its \code{Point::operator<} also compares \code{x} through
\code{ldcmp}; \code{HPI} itself does not need it.)\\[0.3ex]
Area of the intersection of convex polygons: push every edge of every polygon (each
counter-clockwise), then the shoelace sum over \code{HPI}'s output.}

%DEP point
%FROM Ruiming/geomtery/Half_Plane_Intersect.cpp
%LABEL hpi
\begin{lstlisting}
struct Line {
    Point p1, p2;
    Vector v;
    ld ang;
    Line() : p1(0,0), p2(0,0), v(0,0), ang(0) {}
    Line(Point p1, Point p2) : p1(p1), p2(p2) {
        v = p2 - p1;
        ang = atan2l(v.y, v.x);
    }
    Line(Point P, ld angle) : p1(P) {
        p2 = p1 + Point(cosl(angle), sinl(angle));
        v = p2 - p1;
        ang = atan2l(v.y, v.x);
    }
    Line(ld a, ld b, ld c) {
        if (ldcmp(a, 0) == 0) {
            p1 = Point(0, -c/b);
            p2 = Point(1, -c/b);
        } else if (ldcmp(b, 0) == 0) {
            p1 = Point(-c/a, 0);
            p2 = Point(-c/a, 1);
        } else {
            p1 = Point(0, -c/b);
            p2 = Point(1, (-c-a)/b);
        }
        v = p2 - p1;
        ang = atan2l(v.y, v.x);
    }
    bool operator<(Line L) {
        if (ldcmp(ang, L.ang) != 0) return ang < L.ang;
        return ldcmp(Cross(L.p1 - p1, L.p2 - p1), 0) > 0;
    }
};
bool OnLeft(Line L, Point p) {
    return ldcmp(Cross(L.v, p - L.p1), 0) >= 0;
}
Point Cross_point(Line a, Line b) {
    Vector u = a.p1 - b.p1;
    ld t = Cross(b.v, u) / Cross(a.v, b.v);
    return a.p1 + a.v * t;
}
vector<Point> HPI(vector<Line> L) {
    int n = L.size();
    if (n == 0) return {};
    sort(L.begin(), L.end());
    // dedupe by polar angle, keep only the leftmost one
    vector<Line> lines;
    for (int i = 0; i < n; i++) {
        if (i > 0 && ldcmp(L[i].ang, L[i-1].ang) == 0)
            continue;
        lines.push_back(L[i]);
    }
    n = lines.size();
    if (n <= 2) return {};
    vector<Line> q(n);
    vector<Point> p(n);
    int first = 0, last = 0;
    q[0] = lines[0];
    for (int i = 1; i < n; i++) {
        while (first < last && !OnLeft(lines[i], p[last-1]))
            last--;
        while (first < last && !OnLeft(lines[i], p[first]))
            first++;
        q[++last] = lines[i];
        if (ldcmp(Cross(q[last].v, q[last-1].v), 0) == 0) {
            last--;
            if (OnLeft(q[last], lines[i].p1))
                q[last] = lines[i];
        }
        if (first < last) {
            p[last-1] = Cross_point(q[last-1], q[last]);
        }
    }
    while (first < last && !OnLeft(q[first], p[last-1]))
        last--;
    if (last - first <= 1) return {};
    p[last] = Cross_point(q[last], q[first]);
    vector<Point> ans;
    for (int i = first; i <= last; i++) ans.push_back(p[i]);
    return ans;
}
\end{lstlisting}

\entry{Lattice points (exact)}{Ruiming}{%
The integer versions: \code{PointLL} / \code{LineLL}, no \code{eps}, so every sign
test is \textbf{exact}. Same names and return codes as the floating-point functions
(they overload), plus \code{Dist2} and \code{Area2} / \code{Polygon\_area2}, which are
\emph{twice} the signed area. Projection, reflection, intersection and rotation return a
floating \code{Point}.\\[0.3ex]
\code{Point\_in\_polygon} is 1-indexed and returns \code{3} vertex, \code{2} edge,
\code{1} inside, \code{0} outside. Coordinates up to $10^9$ keep \code{Cross} inside
\code{ll}.}

%DEP point
%LABEL lattice
%FROM Ruiming/geomtery/Lattic_points.cpp
\begin{lstlisting}
// ==================== integer point ====================
struct PointLL {
    ll x, y;
    PointLL() : x(0), y(0) {}
    PointLL(ll x, ll y) : x(x), y(y) {}
    PointLL operator+(const PointLL& B) const {
        return PointLL(x + B.x, y + B.y); }
    PointLL operator-(const PointLL& B) const {
        return PointLL(x - B.x, y - B.y); }
    PointLL operator*(ll k) const {
        return PointLL(x * k, y * k); }
    PointLL operator/(ll k) const {
        return PointLL(x / k, y / k); }
    bool operator==(const PointLL& B) const {
        return x == B.x && y == B.y; }
    bool operator!=(const PointLL& B) const {
        return !(*this == B); }
};
typedef PointLL VectorLL;
// dot product
ll Dot(const VectorLL& A, const VectorLL& B) {
    return A.x * B.x + A.y * B.y;
}
// cross product
ll Cross(const VectorLL& A, const VectorLL& B) {
    return A.x * B.y - A.y * B.x;
}
// squared distance
ll Dist2(const PointLL& A, const PointLL& B) {
    ll dx = A.x - B.x, dy = A.y - B.y;
    return dx * dx + dy * dy;
}
// twice the signed area
ll Area2(const PointLL& A, const PointLL& B,
         const PointLL& C) {
    return Cross(B - A, C - A);
}
// parallel test
bool Parallel(const VectorLL& A, const VectorLL& B) {
    return Cross(A, B) == 0;
}
// line (integer points)
struct LineLL {
    PointLL p1, p2;
    LineLL() {}
    LineLL(const PointLL& p1, const PointLL& p2)
        : p1(p1), p2(p2) {}
};
// point vs line: 1 left, 2 right, 0 on the line
ll Point_line_relation(const PointLL& P, const LineLL& v) {
    ll c = Cross(P - v.p1, v.p2 - v.p1);
    if (c < 0) return 1;
    if (c > 0) return 2;
    return 0;
}
// is the point on the segment
bool Point_on_seg(const PointLL& P, const LineLL& v) {
    return Cross(P - v.p1, v.p2 - v.p1) == 0 &&
           min(v.p1.x, v.p2.x) <= P.x &&
           P.x <= max(v.p1.x, v.p2.x) &&
           min(v.p1.y, v.p2.y) <= P.y &&
           P.y <= max(v.p1.y, v.p2.y);
}
// two lines: 0 parallel, 1 same line, 2 intersecting
ll Line_relation(const LineLL& v1, const LineLL& v2) {
    if (Cross(v1.p2 - v1.p1, v2.p2 - v2.p1) == 0) {
        if (Cross(v1.p1 - v2.p1, v2.p2 - v2.p1) == 0)
            return 1;
        return 0;
    }
    return 2;
}
// point vs polygon (exact integers)
// returns: 3 vertex, 2 edge, 1 inside, 0 outside
ll Point_in_polygon(const PointLL& pt, const PointLL* p,
                    ll n) {
    for (int i = 1; i <= n; i++) {
        if (p[i] == pt) return 3;
    }
    for (int i = 1; i <= n; i++) {
        int nxt = (i % n) + 1;
        LineLL v(p[i], p[nxt]);
        if (Point_on_seg(pt, v)) return 2;
    }
    bool inside = false;
    for (int i = 1, j = n; i <= n; j = i++) {
        if ((p[i].y > pt.y) != (p[j].y > pt.y)) {
            ll dy = p[j].y - p[i].y;
            ll dx = p[j].x - p[i].x;
            ll lhs = (pt.x - p[i].x) * dy;
            ll rhs = (pt.y - p[i].y) * dx;
            if (dy > 0) {
                if (lhs < rhs) inside = !inside;
            } else {
                if (lhs > rhs) inside = !inside;
            }
        }
    }
    return inside ? 1 : 0;
}
// twice the polygon area (signed), as ll
ll Polygon_area2(const PointLL* p, ll n) {
    ll area = 0;
    for (int i = 1; i <= n; i++) {
        int nxt = (i % n) + 1;
        area += p[i].x * p[nxt].y - p[i].y * p[nxt].x;
    }
    return area;
}
// distance
ld Dist(const PointLL& A, const PointLL& B) {
    return sqrt((ld)Dist2(A, B));
}
// point to line distance
ld Dis_point_line(const PointLL& P, const LineLL& v) {
    return fabs((ld)Cross(P - v.p1, v.p2 - v.p1))
           / Dist(v.p1, v.p2);
}
// point to segment distance
ld Dis_point_seg(const PointLL& P, const LineLL& v) {
    if (Dot(P - v.p1, v.p2 - v.p1) < 0
        || Dot(P - v.p2, v.p1 - v.p2) < 0) {
        return min(Dist(P, v.p1), Dist(P, v.p2));
    }
    return Dis_point_line(P, v);
}
// projection onto a line (floating point)
Point Point_line_proj(const PointLL& P, const LineLL& v) {
    ld k = (ld)Dot(v.p2 - v.p1, P - v.p1)
           / Dot(v.p2 - v.p1, v.p2 - v.p1);
    return Point(v.p1.x + (v.p2.x - v.p1.x) * k,
                 v.p1.y + (v.p2.y - v.p1.y) * k);
}
// reflection in a line (floating point)
Point Point_line_symmetry(const PointLL& P, const LineLL& v) {
    Point q = Point_line_proj(P, v);
    return Point(2 * q.x - P.x, 2 * q.y - P.y);
}
// line intersection (floating point). Lines must not be
// parallel.
Point Cross_point(const PointLL& a, const PointLL& b,
                  const PointLL& c, const PointLL& d) {
    ld s1 = (ld)Cross(b - a, c - a);
    ld s2 = (ld)Cross(b - a, d - a);
    ld denom = s2 - s1;
    return Point(
        (c.x * s2 - d.x * s1) / denom,
        (c.y * s2 - d.y * s1) / denom
    );
}
// rotate a vector (floating point)
Point Rotate(const PointLL& A, ld rad) {
    ld c = cos(rad), s = sin(rad);
    return Point(A.x * c - A.y * s, A.x * s + A.y * c);
}
\end{lstlisting}

\entry{Closest pair}{$O(n \log n)$ \quad mem $n$}{%
Not Ruiming's --- written on his \code{PointLL}. Returns the smallest \textbf{squared}
distance and the two points, exact. Needs $n \ge 2$; duplicate points give 0.\\[0.3ex]
Sweep in \code{x}, keep a \code{set} ordered by \code{y} holding everything within the
current best distance, and only compare against the \code{y}-band.}

%DEP point lattice
%LABEL closest
\begin{lstlisting}
pair<ll, pair<PointLL,PointLL>> closest(vector<PointLL> p) {
    sort(all(p), [](const PointLL& a, const PointLL& b) {
        return tie(a.x, a.y) < tie(b.x, b.y); });
    ll best = LLONG_MAX;
    pair<PointLL,PointLL> bp{p[0], p[1]};
    set<pii> s;                          // (y, x)
    int j = 0;
    rep(i, 0, (ll)p.size()) {
        while (j < i) {
            ll dx = p[i].x - p[j].x;
            if (dx * dx <= best) break;
            s.erase({p[j].y, p[j].x}); j++;
        }
        ll d = (ll)sqrtl((ld)best) + 1;
        auto lo = s.lower_bound({p[i].y - d, LLONG_MIN});
        auto hi = s.upper_bound({p[i].y + d, LLONG_MAX});
        for (auto it = lo; it != hi; ++it) {
            PointLL q(it->second, it->first);
            ll dd = Dist2(p[i], q);
            if (dd < best) { best = dd; bp = {q, p[i]}; }
        }
        s.insert({p[i].y, p[i].x});
    }
    return {best, bp};
}
\end{lstlisting}

\entry{Circles}{}{%
Not Ruiming's --- written on his \code{Point}. \code{circleLine} intersects a circle
with the \emph{infinite} line through \code{a,b}; \code{circleCircle} with another
circle. Both return 0, 1 (tangent) or 2 points. For a segment, filter the results with
\code{Point\_on\_seg}. Identical circles return \code{\{\}}.\\[0.3ex]
Tangent length from outside point \code{p} to circle \code{(c,r)} is
$\sqrt{|pc|^2 - r^2}$. Circumcentre of \code{a,b,c}: \code{Cross\_point} of the two
perpendicular bisectors.}

%DEP point
%LABEL circle
\begin{lstlisting}
vector<Point> circleLine(Point c, ld r, Point a, Point b) {
    Vector ab = b - a;
    Point m = a + ab * (Dot(c - a, ab) / Dot(ab, ab));
    ld h2 = r * r - Dot(m - c, m - c);
    if (h2 < -1e-12) return {};
    if (h2 <= 1e-12) return {m};
    Vector h = ab / Len(ab) * sqrtl(h2);
    return {m - h, m + h};
}
vector<Point> circleCircle(Point c1, ld r1, Point c2, ld r2) {
    Vector d = c2 - c1;
    ld dd = Dot(d, d), dist = sqrtl(dd);
    if (dist < 1e-12) return {};          // concentric
    if (dist > r1 + r2 + 1e-12) return {};
    if (dist < fabsl(r1 - r2) - 1e-12) return {};
    ld a = (dd + r1 * r1 - r2 * r2) / (2 * dist);
    ld h2 = r1 * r1 - a * a;
    Vector u = d / dist;
    Point m = c1 + u * a;
    if (h2 <= 1e-12) return {m};
    Vector per = Vector(-u.y, u.x) * sqrtl(h2);
    return {m - per, m + per};
}
\end{lstlisting}

\entry{Hull diameter, point in convex hull}{KACTL \quad $O(n \log n)$, $O(n)$, $O(\log n)$}{%
KACTL's, ported to Ruiming's exact \code{PointLL} (\code{a.cross(b,c)} is
\code{Area2(a,b,c)}). All three are 0-indexed \code{vector}s.\\[0.3ex]
\code{hullCCW}: the hull as an actual \textbf{counter-clockwise cycle}, collinear points
dropped, first point not repeated --- what the two below need, and what Ruiming's
\code{Convex\_Hull} (sorted output) is not.\\[0.3ex]
\code{hullDiameter}: the farthest pair of points (rotating calipers); its squared
distance is \code{Dist2(r[0], r[1])}. Run it on the hull, not the raw points.\\[0.3ex]
\code{inHull(h, p, strict)}: is \code{p} inside the CCW hull \code{h}, by binary search
around \code{h[0]} --- for many queries against one polygon, where Ruiming's $O(n)$
\code{Point\_in\_polygon} is too slow. \code{strict = true} counts the boundary as
outside.}

%DEP point lattice
%LABEL hullkactl
\begin{lstlisting}
typedef PointLL PL;
vector<PL> hullCCW(vector<PL> pts) {
    if (pts.size() <= 1) return pts;
    sort(all(pts), [](const PL& a, const PL& b) {
        return tie(a.x, a.y) < tie(b.x, b.y); });
    vector<PL> h(pts.size() + 1);
    int s = 0, t = 0;
    for (int it = 2; it--; s = --t, reverse(all(pts)))
        for (PL p : pts) {
            while (t >= s+2 && Area2(h[t-2], h[t-1], p) <= 0)
                t--;
            h[t++] = p;
        }
    return {h.begin(), h.begin() + t - (t==2 && h[0]==h[1])};
}
array<PL, 2> hullDiameter(vector<PL> S) {
    int n = (int)S.size(), j = n < 2 ? 0 : 1;
    pair<ll, array<PL, 2>> res({0, {S[0], S[0]}});
    rep(i,0,j)
        for (;; j = (j + 1) % n) {
            ll d = Dist2(S[i], S[j]);
            if (d > res.first) res = {d, {S[i], S[j]}};
            if (Cross(S[(j+1)%n] - S[j], S[i+1] - S[i]) >= 0)
                break;
        }
    return res.second;
}
int side(PL s, PL e, PL p) {             // KACTL sideOf
    ll c = Area2(s, e, p); return (c > 0) - (c < 0); }
bool inHull(const vector<PL>& l, PL p, bool strict = true) {
    int a = 1, b = (int)l.size() - 1, r = !strict;
    if (l.size() < 3)
        return r && Point_on_seg(p, LineLL(l[0], l.back()));
    if (side(l[0], l[a], l[b]) > 0) swap(a, b);
    if (side(l[0], l[a], p) >= r || side(l[0], l[b], p) <= -r)
        return false;
    while (abs(a - b) > 1) {
        int c = (a + b) / 2;
        (side(l[0], l[c], p) > 0 ? b : a) = c;
    }
    return side(l[a], l[b], p) < r;
}
\end{lstlisting}

\entry{Enclosing circle, polygon cut, tangents}{KACTL}{%
KACTL's, ported to Ruiming's floating \code{Point}. \code{perp} is the vector turned
$90^\circ$ counter-clockwise.\\[0.3ex]
\code{mec(ps)}: the \textbf{smallest circle containing all points}, \code{\{centre,
radius\}}. Randomised incremental, expected $O(n)$ --- the shuffle is what makes it
fast, keep it. \code{ccCenter} is the circumcentre of a triangle (divides by zero if the
three points are collinear).\\[0.3ex]
\code{polygonCut(poly, s, e)}: the part of a polygon (any orientation, need not be
convex) lying to the \textbf{right} of the directed line \code{s -> e}; the part to the
left is cut away. Swap \code{s, e} to keep the other side. Several cuts in a row on a
convex polygon = half-plane intersection when the box is known.\\[0.3ex]
\code{tangents(c1, r1, c2, r2)}: the \textbf{outer} common tangents of two circles, as
pairs of touching points (0, 1 or 2 of them). \textbf{Inner} tangents: pass
\code{-r2}. Tangents from a point: \code{r2 = 0}. Empty if one circle is inside the
other or they coincide.}

%DEP point
%LABEL mec
\begin{lstlisting}
Vector perp(Vector a) { return Vector(-a.y, a.x); }
Point ccCenter(Point A, Point B, Point C) {
    Vector b = C - A, c = B - A;
    return A + perp(b * Dot(c, c) - c * Dot(b, b))
               / Cross(b, c) / 2;
}
pair<Point, ld> mec(vector<Point> ps) {
    shuffle(all(ps), mt19937(time(0)));
    Point o = ps[0];
    ld r = 0, EPS = 1 + 1e-8;
    rep(i,0,(ll)ps.size()) if (Dist(o, ps[i]) > r * EPS) {
        o = ps[i], r = 0;
        rep(j,0,i) if (Dist(o, ps[j]) > r * EPS) {
            o = (ps[i] + ps[j]) / 2;
            r = Dist(o, ps[i]);
            rep(k,0,j) if (Dist(o, ps[k]) > r * EPS) {
                o = ccCenter(ps[i], ps[j], ps[k]);
                r = Dist(o, ps[i]);
            }
        }
    }
    return {o, r};
}
vector<Point> polygonCut(const vector<Point>& poly,
                         Point s, Point e) {
    vector<Point> res;
    rep(i,0,(ll)poly.size()) {
        Point cur = poly[i];
        Point prev = i ? poly[i-1] : poly.back();
        ld a = Cross(e-s, cur-s), b = Cross(e-s, prev-s);
        if ((a < 0) != (b < 0))
            res.push_back(cur + (prev - cur) * (a / (a - b)));
        if (a < 0)
            res.push_back(cur);
    }
    return res;
}
vector<pair<Point, Point>> tangents(Point c1, ld r1,
                                    Point c2, ld r2) {
    Vector d = c2 - c1;
    ld dr = r1 - r2, d2 = Dot(d, d), h2 = d2 - dr * dr;
    if (d2 == 0 || h2 < 0) return {};
    vector<pair<Point, Point>> out;
    for (ld sign : {-1, 1}) {
        Vector v = (d * dr + perp(d) * sqrtl(h2) * sign) / d2;
        out.push_back({c1 + v * r1, c2 + v * r2});
    }
    if (h2 == 0) out.pop_back();
    return out;
}
\end{lstlisting}


\section{Searching}

\entry{Binary search on a predicate}{$O(\log)$}{%
\code{pred} must be monotone false$\dots$false true$\dots$true. Returns the first
\code{x} in \code{[lo, hi]} with \code{pred(x)}, or \code{hi+1} if none --- so test the
return value before using it.\\[0.3ex]
``Largest \code{x} that still works'' is the same call on the negated predicate minus
one, or just search for the first failing \code{x} and subtract 1. Almost every
``minimise the maximum'' / ``can we do it within $k$'' problem is this function plus a
greedy check.\\[0.3ex]
Reals: fix the iteration count (100 halvings covers any range to $10^{-30}$) rather than
looping on a tolerance, which can spin forever.}

\begin{lstlisting}
template<class F>
ll bsearch(ll lo, ll hi, F pred) {   // first true in [lo,hi]
    hi++;                            // hi = "none found"
    while (lo < hi) {
        ll m = lo + (hi - lo) / 2;   // no overflow
        if (pred(m)) hi = m; else lo = m + 1;
    }
    return lo;
}
template<class F>
ld bsearchReal(ld lo, ld hi, F pred) {
    rep(it, 0, 100) {
        ld m = (lo + hi) / 2;
        if (pred(m)) hi = m; else lo = m;
    }
    return lo;
}
\end{lstlisting}

\entry{Ternary search}{$O(\log)$}{%
For a \textbf{strictly unimodal} function --- one valley, no flat stretches. A flat
region breaks it silently, so if ties are possible make the comparison
\code{<=} vs \code{<} deliberately, or binary-search on
\code{f(m+1) - f(m)} instead, which is the more robust trick.\\[0.3ex]
The integer version finishes by scanning the last three candidates, which is where
off-by-ones normally live. For a maximum, negate \code{f}.}

\begin{lstlisting}
template<class F>              // argmin of a unimodal f
ll tsearch(ll lo, ll hi, F f) {
    while (hi - lo > 2) {
        ll m1 = lo + (hi - lo) / 3, m2 = hi - (hi - lo) / 3;
        if (f(m1) < f(m2)) hi = m2; else lo = m1;
    }
    ll best = lo;
    for (ll x = lo + 1; x <= hi; x++)
        if (f(x) < f(best)) best = x;
    return best;
}
template<class F>
ld tsearchReal(ld lo, ld hi, F f) {
    rep(it, 0, 200) {
        ld m1 = lo + (hi - lo) / 3, m2 = hi - (hi - lo) / 3;
        if (f(m1) < f(m2)) hi = m2; else lo = m1;
    }
    return (lo + hi) / 2;
}
\end{lstlisting}

\entry{Coordinate compression}{$O(n \log n)$}{%
Maps arbitrary values onto \code{0 .. k-1} keeping order, so a \code{BIT} or \code{Seg}
can index them. This is what lets you avoid \code{SparseSeg} whenever the whole input is
readable up front --- which is almost always.\\[0.3ex]
Remember to compress \emph{every} value you will ever look up, queries included, not
just the ones in the array. For half-open interval problems add both \code{l} and
\code{r} (and sometimes \code{r+1}).}

\begin{lstlisting}
// returns the sorted distinct values; id(v) is v's index
vi compress(vi vals) {
    sort(all(vals));
    vals.erase(unique(all(vals)), vals.end());
    return vals;
}
int id(const vi& c, ll v) {
    return (int)(lower_bound(all(c), v) - c.begin());
}
\end{lstlisting}

\section{Reference}

\entry{What fits in the time limit}{$\sim10^8$ simple ops/s at \code{-O2}}{%
{\ttfamily\footnotesize
\begin{ctab}{0.20}{0.75}
n <= 10 & O(n!) \hfill brute force permutations\\
n <= 20 & O(2\^{}n n) \hfill bitmask DP, PathDP\\
n <= 40 & O(2\^{}(n/2)) \hfill meet in the middle\\
n <= 100 & O(n\^{}4) \hfill 4 nested loops, Hungarian\\
n <= 500 & O(n\^{}3) \hfill Floyd, matrix, interval DP\\
n <= 5000 & O(n\^{}2) \hfill pairwise DP\\
n <= 1e5 & O(n log\^{}2 n) \hfill HLD + Seg\\
n <= 1e6 & O(n log n) \hfill sort, Seg, Dijkstra\\
n <= 1e7 & O(n) \hfill sieve, two pointers\\
n <= 1e18 & O(log n) \hfill binary search, matrix power\\
\end{ctab}}\\[0.3ex]
A \code{map} / \code{set} operation costs $\sim$10x a \code{vector} one; sorting
$10^6$ \code{ll} is $\sim$0.1 s; $10^8$ additions is $\sim$0.1 s but $10^8$ random
memory accesses is $\sim$1 s. If $n \le 18$ and nothing else fits, it is a bitmask.}

\entry{Overflow and precision}{read this when a big test fails}{%
\code{ll} holds $9.2\cdot10^{18}$, \code{int} only $2.1\cdot10^9$ --- and
\code{n*(n-1)/2} already overflows \code{int} at \code{n = 1e5}.\\[0.3ex]
Use \code{INF = 1e18}, never \code{LLONG\_MAX}: \code{INF + INF} still fits, so a
forgotten guard gives a wrong answer instead of undefined behaviour. Guard anyway.\\[0.3ex]
\code{a * b \% m} overflows once $m > 3\cdot10^9$ --- cast to \code{i128}.
\code{cross} on coordinates up to $10^9$ reaches $2\cdot10^{18}$: fine, barely.\\[0.3ex]
\code{-7 \% 3 == -1} in C++. Normalise with \code{((x \% m) + m) \% m}.\\[0.3ex]
\code{pow(10, 2)} can return \code{99.999...}; never use \code{pow} for integers.
\code{sqrt(x)} likewise --- take \code{(ll)sqrtl(x)} then fix up by $\pm1$ in a loop.\\[0.3ex]
Compare squared distances, never \code{sqrt}ed ones. \code{double} has 15--16 significant
digits, \code{long double} 18--19 on x86.}

\entry{Pitfalls}{the usual wrong answers}{%
\textbf{Multiple test cases:} clear every global, resize every vector, reset
\code{cnt}/\code{timer}. The most common repeat-offender bug there is.\\[0.3ex]
\textbf{Recursion depth:} a DFS on a path of $2\cdot10^5$ vertices can blow the stack.
The iterative \code{euler}, \code{SCC}, \code{Bridges}, \code{centroids} and
\code{eulerPath} here are safe; \code{HLD}, \code{Dinic::dfs} and \code{Matching} are
not.\\[0.3ex]
\textbf{Output:} \code{'\textbackslash n'} not \code{endl} (\code{endl} flushes).
Keep \code{sync\_with\_stdio(false)}, and then never mix \code{printf} with
\code{cout}.\\[0.3ex]
\textbf{\code{unordered\_map}} is $O(n)$ per operation against a crafted test; use
\code{map}, or add a random offset to every key.\\[0.3ex]
\textbf{Comparators} must be a strict weak ordering --- \code{<=} in a \code{sort}
comparator is a crash, not a wrong answer.\\[0.3ex]
\textbf{\code{vector<bool>}} is a proxy type; \code{auto\& x = v[i]} does not do what you
expect. \textbf{Iterators} invalidate on \code{push\_back}.\\[0.3ex]
\textbf{Reading:} after \code{cin >> n}, a following \code{getline} reads the rest of that
line. Interactive problems need an explicit flush.\\[0.3ex]
\textbf{Edge cases:} $n = 1$, $n = 0$, an empty string, all elements equal, a
disconnected graph, self-loops, multi-edges, duplicate points, collinear points.}

\entry{Combinatorics}{}{%
$\binom{n}{k} = \binom{n-1}{k-1} + \binom{n-1}{k}$; hockey stick
$\sum_{i=r}^{n}\binom{i}{r} = \binom{n+1}{r+1}$; Vandermonde
$\sum_k \binom{m}{k}\binom{n}{p-k} = \binom{m+n}{p}$.\\[0.3ex]
\textbf{Stars and bars.} $n$ identical items into $k$ ordered boxes:
$\binom{n+k-1}{k-1}$, and $\binom{n-1}{k-1}$ if every box must be non-empty.\\[0.3ex]
\textbf{Catalan} $C_n = \frac{1}{n+1}\binom{2n}{n}$: balanced bracket strings,
binary trees on $n$ nodes, triangulations of an $(n{+}2)$-gon, paths under the
diagonal. $1, 1, 2, 5, 14, 42, 132, 429$.\\[0.3ex]
\textbf{Derangements} $D_n = (n-1)(D_{n-1} + D_{n-2})$, $D_0 = 1$, $D_1 = 0$.\\[0.3ex]
\textbf{Inclusion-exclusion}
$|\bigcup A_i| = \sum|A_i| - \sum|A_i \cap A_j| + \dots$ --- over $\le 20$ sets this is
a mask loop with $(-1)^{\mathrm{popcount}}$.\\[0.3ex]
\textbf{Burnside.} Orbits = average number of fixed points over the group. Necklaces of
$n$ beads in $k$ colours under rotation: $\frac{1}{n}\sum_{i=0}^{n-1} k^{\gcd(i,n)}$.\\[0.3ex]
\textbf{Lucas.} $\binom{n}{k} \bmod p$ for prime $p$ = product of
$\binom{n_i}{k_i}$ over the base-$p$ digits (0 if any $k_i > n_i$).\\[0.3ex]
$\sum_{i\le n} i = \frac{n(n+1)}{2}$, $\sum i^2 = \frac{n(n+1)(2n+1)}{6}$,
$\sum i^3 = \left(\frac{n(n+1)}{2}\right)^2$.}

\entry{Number theory}{}{%
$\gcd(a,b)\cdot\mathrm{lcm}(a,b) = ab$ --- compute \code{a / gcd * b} to avoid
overflow.\\[0.3ex]
$\varphi(n) = n\prod_{p\mid n}(1-\frac1p)$; $\sum_{d\mid n}\varphi(d) = n$.
Fermat: $a^{p-1} \equiv 1$ for prime $p \nmid a$. Euler:
$a^{\varphi(m)} \equiv 1$ when $\gcd(a,m)=1$.\\[0.3ex]
$d(n) = \prod(e_i+1)$ divisors, $\sigma(n) = \prod\frac{p_i^{e_i+1}-1}{p_i-1}$.
The number of divisors stays small: at most 240 below $10^9$, 1344 below
$10^{18}$.\\[0.3ex]
\textbf{Legendre.} The exponent of $p$ in $n!$ is
$\sum_{i\ge1}\lfloor n/p^i \rfloor$ --- that is how you count trailing zeros of a
factorial.\\[0.3ex]
\textbf{M\"obius} $\mu(n)$ is $0$ with a squared factor, else $(-1)^{\#\text{primes}}$;
$\sum_{d\mid n}\mu(d) = [n=1]$, which turns ``count coprime pairs'' into a divisor
sum.\\[0.3ex]
Primes below $10^k$: $4$, $25$, $168$, $1229$, $9592$, $78498$, $664579$, $5761455$.
The gap between consecutive primes near $10^{18}$ stays under 1500.}

\entry{Graph facts}{}{%
\textbf{Max flow = min cut.} \textbf{K\"onig} (bipartite): max matching = min vertex
cover, and max independent set $= n - $ max matching. \textbf{Hall:} a perfect matching
on the left exists iff every subset $S$ of the left has
$|N(S)| \ge |S|$.\\[0.3ex]
\textbf{Dilworth:} the minimum number of chains covering a poset equals the largest
antichain --- in practice a bipartite matching or an LIS.\\[0.3ex]
Bipartite iff no odd cycle. A tree has $n-1$ edges and a unique path between any two
vertices. A connected graph has an Euler circuit iff every degree is even (directed: in
= out everywhere).\\[0.3ex]
\textbf{Planar:} $V - E + F = 2$, so $E \le 3V - 6$ and some vertex has degree
$\le 5$. Four colours always suffice, five are easy to find.\\[0.3ex]
\textbf{Spanning trees:} Cayley gives $n^{n-2}$ labelled trees on $n$ vertices; for a
general graph it is any cofactor of the Laplacian (Matrix-Tree), i.e. one
determinant by Gaussian elimination.\\[0.3ex]
\textbf{Pick:} a lattice polygon has $A = I + B/2 - 1$, with
$B = \sum \gcd(|dx|,|dy|)$ over the edges.}

\entry{When stuck}{}{%
Re-read the statement and the constraints. The bound on $n$ usually names the intended
complexity outright --- check the table above before designing anything.\\[0.3ex]
Write the brute force first, then stress test against it. A wrong answer on $n=3$ is
worth more than an hour of staring.\\[0.3ex]
Reframe: is it a graph? a flow? a matching? Can you binary-search the answer? Sort the
input? Process queries offline, sorted? Think backwards from the end? Count the
complement? Fix one variable and see what is left?\\[0.3ex]
Check the samples \emph{by hand}. If the sample passes but the judge does not, the bug is
almost always an edge case from the Pitfalls list --- most often multiple test cases, or
an overflow.\\[0.3ex]
\footnotesize
\emph{Sources: templates by Ruben, Ruiming and Martin, consolidated;
\code{LineContainer} and every entry marked KACTL are adapted from KACTL (KTH; licences per file in the
README) --- \code{sz(x)} written out,
\code{int} widened to \code{ll} where the header's \code{vi} needed it. Provenance and the per-entry diff
against the original files are in \code{notebook/README.md}.}}


\end{multicols*}
\end{document}

